\documentclass[11pt]{article}
% Shared preamble for all solution files.
\usepackage[margin=1.1in]{geometry}
\usepackage{enumitem}
\usepackage{tikz}
\usepackage[hidelinks]{hyperref}
\usepackage{amsmath,amssymb,amsthm}
\newcommand{\fa}{\mathfrak a}
\newcommand{\fb}{\mathfrak b}
\newcommand{\fc}{\mathfrak c}
\newcommand{\fp}{\mathfrak p}
\newcommand{\fq}{\mathfrak q}
\newcommand{\fm}{\mathfrak m}
\newcommand{\fn}{\mathfrak n}
\newcommand{\fr}{\mathfrak r}
\newcommand{\Nil}{\mathfrak N}
\newcommand{\Jac}{\mathfrak R}
\newcommand{\rad}{r}
\newcommand{\Spec}{\operatorname{Spec}}
\newcommand{\Max}{\operatorname{Max}}
\newcommand{\Ann}{\operatorname{Ann}}
\newcommand{\Ker}{\operatorname{Ker}}
\newcommand{\Z}{\mathbf Z}
\newcommand{\Q}{\mathbf Q}
\newcommand{\R}{\mathbf R}
\newcommand{\C}{\mathbf C}
\newcommand{\ci}{\mathrm{c}}
\newcommand{\ex}{\mathrm{e}}

\setlength{\parindent}{0pt}
\setlength{\parskip}{0.55em}
% \exercise{label}{statement}  -- followed by the solution text
\newcommand{\exercise}[2]{\par\bigskip\noindent{\large\textbf{Exercise #1}}\par\smallskip
  \begin{quote}\small\textbf{Problem.}\ #2\end{quote}\par}
\newcommand{\pf}{\noindent\textbf{Solution.}\ }

\title{Hartshorne, Chapter I: Varieties\\[2pt]\large Solutions to the exercises}
\author{}
\date{}
\begin{document}
\maketitle
\noindent\textbf{Conventions.} Throughout, $k$ is an algebraically closed field, varieties live in $\mathbf A^n$ or $\mathbf P^n$ with the Zariski topology, and results of the text are cited by their numbers in Hartshorne (``Thm.~1.8A'', ``Prop.~3.5'', \dots). The numbering and statements of the exercises follow the edition scanned in this repository. Exercises are grouped by section, and each is restated before its solution; a starred exercise in the book is marked ``Starred in the book'' at the end of its statement.

Honest notes. A few exercises are not solvable as printed or go beyond what Chapter I can prove: Ex.~3.19(b) is the Jacobian Conjecture (open), so only what can be proved is given; Ex.~2.17(d) is an open-problem remark. In several places the printed statement needs a hypothesis (for example a characteristic condition in Ex.~5.1, 5.5, 5.11, 5.15, or ``projective \emph{variety}'' in Ex.~2.8) and the solution says what is assumed. A few solutions use standard commutative algebra beyond Chapter I (Krull's theorem on normal Noetherian domains, Noether normalization, Weierstrass preparation); these are named where used.

\section*{1. Affine varieties}

Throughout, $k$ is an algebraically closed field (hence infinite), $A=k[x_1,\dots,x_n]$, and I use the numbering of Hartshorne's Chapter~I (Thm.~1.3A is the Nullstellensatz, Cor.~1.4 the correspondence between algebraic sets and radical ideals, Prop.~1.7 and Thm.~1.8A the dimension theory, Thm.~1.11A Krull's Hauptidealsatz).

\exercise{1.1}{
\begin{enumerate}[label=(\alph*),leftmargin=*]
\item Let $Y$ be the plane curve $y=x^2$ (i.e.\ $Y$ is the zero set of $f=y-x^2$). Show that $A(Y)$ is isomorphic to a polynomial ring in one variable over $k$.
\item Let $Z$ be the plane curve $xy=1$. Show that $A(Z)$ is not isomorphic to a polynomial ring in one variable over $k$.
\item Let $f$ be any irreducible quadratic polynomial in $k[x,y]$, and let $W$ be the conic defined by $f$. Show that $A(W)$ is isomorphic to $A(Y)$ or $A(Z)$. Which one is it when?
\end{enumerate}
(Part (c) is starred in the book.)}
\pf
For an irreducible polynomial $g\in k[x,y]$ the ideal $(g)$ is prime ($k[x,y]$ is a UFD), hence radical, so by the Nullstellensatz (Thm.~1.3A, Cor.~1.4) $I(Z(g))=(g)$ and $A(Z(g))=k[x,y]/(g)$.

(a) The polynomial $y-x^2$ is irreducible (degree~$1$ in $y$ and monic), so $A(Y)=k[x,y]/(y-x^2)$. Consider the $k$-algebra homomorphism $\alpha:k[x,y]\to k[t]$, $x\mapsto t$, $y\mapsto t^2$. It is surjective and kills $y-x^2$. Conversely, every $g(x,y)$ is congruent modulo $(y-x^2)$ to $g(x,x^2)\in k[x]$ (substitute $y=x^2+(y-x^2)$ and expand), so if $\alpha(g)=g(t,t^2)=0$ then $g(x,x^2)=0$ and $g\in(y-x^2)$. Thus $\ker\alpha=(y-x^2)$ and $A(Y)\cong k[t]$.

(b) The polynomial $xy-1$ is irreducible (degree $1$ in $x$ with coprime coefficients $y$ and $-1$), so $A(Z)=k[x,y]/(xy-1)$. In $A(Z)$ the class $\bar x$ is a unit ($\bar x\bar y=1$). It is not in $k$: if $x-c\in(xy-1)$ for some $c\in k$, say $x-c=(xy-1)h$, then $h\neq0$ would give degree $\ge 2$ on the right, and $h=0$ would give $x=c$, absurd. So $A(Z)$ has a unit not in $k$. The units of $k[t]$ are exactly the nonzero constants. Since a $k$-algebra isomorphism preserves units and fixes $k$, $A(Z)\not\cong k[t]$.

(c) Write $f=q+(\text{terms of degree}\le1)$ with $q=ax^2+bxy+cy^2\neq0$ the quadratic part. Since $k$ is algebraically closed, the binary form $q$ is a product of two linear forms, $q=\ell_1\ell_2$. Either $\ell_1,\ell_2$ are linearly independent, or proportional (then $q=\lambda\ell^2$, $\lambda\ne0$). For $\ell_i=p_ix+r_iy$ one computes $b^2-4ac=(p_1r_2-p_2r_1)^2$, so the two cases are distinguished by whether the discriminant $b^2-4ac$ vanishes (this is valid in every characteristic). An invertible affine change of coordinates of $\mathbf A^2$ induces a $k$-algebra automorphism of $k[x,y]$ carrying $(f)$ to the ideal of the transformed polynomial, hence does not change the isomorphism class of $A(W)=k[x,y]/(f)$.

\emph{Case 1: $\ell_1,\ell_2$ independent} ($b^2\neq4ac$). Take $u=\ell_1$, $v=\ell_2$ as new coordinates. Then $f=uv+\alpha u+\beta v+\gamma=(u+\beta)(v+\alpha)+\delta$ with $\delta=\gamma-\alpha\beta$. If $\delta=0$, $f$ would be reducible; so $\delta\ne0$. With $u'=-(u+\beta)/\delta$, $v'=v+\alpha$ we get $f=-\delta(u'v'-1)$, so $(f)=(u'v'-1)$ and $A(W)\cong k[u',v']/(u'v'-1)\cong A(Z)$. By (b) it is not $\cong A(Y)$.

\emph{Case 2: $q=\lambda\ell^2$} ($b^2=4ac$). Choose a linear form $m$ independent of $\ell$ and use $u=\ell$, $v=m$ as coordinates. Then $f=\lambda u^2+\alpha u+\beta v+\gamma$. If $\beta=0$, $f=\lambda(u-r_1)(u-r_2)$ would be reducible; so $\beta\ne0$ and $f=\beta(v-h(u))$ with $h(u)=-(\lambda u^2+\alpha u+\gamma)/\beta$. The map $k[u,v]\to k[u]$, $u\mapsto u$, $v\mapsto h(u)$ is surjective, and $g(u,v)\equiv g(u,h(u))\pmod{(v-h)}$, so its kernel is $(v-h)=(f)$. Hence $A(W)\cong k[u]\cong A(Y)$, and by (b) it is not $\cong A(Z)$.

\emph{Answer.} $A(W)\cong A(Y)$ (a polynomial ring in one variable, ``parabola'') exactly when the quadratic part of $f$ is a square of a linear form ($b^2-4ac=0$); $A(W)\cong A(Z)$ (``hyperbola'') exactly when the quadratic part has two distinct linear factors ($b^2-4ac\ne0$). Never both, by (b).\hfill$\square$

\exercise{1.2}{\emph{The twisted cubic curve.} Let $Y\subseteq\mathbf A^3$ be the set $Y=\{(t,t^2,t^3)\mid t\in k\}$. Show that $Y$ is an affine variety of dimension~$1$. Find generators for the ideal $I(Y)$. Show that $A(Y)$ is isomorphic to a polynomial ring in one variable over $k$. We say that $Y$ is given by the \emph{parametric representation} $x=t$, $y=t^2$, $z=t^3$.}
\pf
Put $\mathfrak a=(y-x^2,\,z-x^3)\subseteq k[x,y,z]$. A point $(a,b,c)$ lies in $Z(\mathfrak a)$ iff $b=a^2$, $c=a^3$, i.e.\ iff it equals $(t,t^2,t^3)$ with $t=a$; so $Y=Z(\mathfrak a)$ is an algebraic set.

Let $\alpha:k[x,y,z]\to k[t]$ be the $k$-algebra map $x\mapsto t,\ y\mapsto t^2,\ z\mapsto t^3$. It is surjective and $\mathfrak a\subseteq\ker\alpha$. Conversely, writing $y=x^2+(y-x^2)$, $z=x^3+(z-x^3)$ shows $g(x,y,z)\equiv g(x,x^2,x^3)\pmod{\mathfrak a}$, so $\alpha(g)=g(t,t^2,t^3)=0$ forces $g\in\mathfrak a$. Hence $\ker\alpha=\mathfrak a$ and $k[x,y,z]/\mathfrak a\cong k[t]$, a domain, so $\mathfrak a$ is prime, in particular radical. By the Nullstellensatz (Cor.~1.4), $I(Y)=I(Z(\mathfrak a))=\sqrt{\mathfrak a}=\mathfrak a=(y-x^2,\,z-x^3)$.

Since $I(Y)$ is prime, $Y$ is irreducible, i.e.\ an affine variety (Cor.~1.4), and $A(Y)=k[x,y,z]/I(Y)\cong k[t]$. By Prop.~1.7 and Thm.~1.8A(a), $\dim Y=\dim k[t]=\operatorname{tr.deg}_k k(t)=1$ (equivalently, Prop.~1.9 gives $\dim\mathbf A^1=1$).\hfill$\square$

\exercise{1.3}{Let $Y$ be the algebraic set in $\mathbf A^3$ defined by the two polynomials $x^2-yz$ and $xz-x$. Show that $Y$ is a union of three irreducible components. Describe them and find their prime ideals.}
\pf
Let $(x,y,z)\in Y$. From $x(z-1)=0$: either $x=0$ or $z=1$.
\begin{itemize}
\item If $x=0$, then $yz=0$, so $y=0$ or $z=0$.
\item If $z=1$, then $y=x^2$.
\end{itemize}
Conversely each of the three sets below lies in $Y$. Hence $Y=L_1\cup L_2\cup C$ where
\[
L_1=\{x=y=0\}\ (\text{the $z$-axis}),\qquad L_2=\{x=z=0\}\ (\text{the $y$-axis}),\qquad C=\{z=1,\ y=x^2\}\ (\text{a parabola}).
\]
Their ideals are $\mathfrak p_1=(x,y)$, $\mathfrak p_2=(x,z)$, $\mathfrak p_3=(z-1,\,y-x^2)$. These ideals are prime because $k[x,y,z]/\mathfrak p_1\cong k[z]$, $k[x,y,z]/\mathfrak p_2\cong k[y]$, $k[x,y,z]/\mathfrak p_3\cong k[x]$ are domains; being prime they are radical, so by the Nullstellensatz they are the ideals $I(L_1),I(L_2),I(C)$ and $L_1,L_2,C$ are irreducible (Cor.~1.4).

None contains another: the point $(0,0,c)\in L_1$ with $c\notin\{0,1\}$ (possible as $k$ is infinite) lies in neither $L_2$ (needs $z=0$) nor $C$ (needs $z=1$); $(0,c,0)\in L_2$ with $c\ne0$ lies in neither $L_1$ (needs $y=0$) nor $C$ ($z=0\ne1$); $(1,1,1)\in C$ lies in neither $L_1$ nor $L_2$ ($x\ne0$). So $Y=L_1\cup L_2\cup C$ is an irredundant decomposition into irreducible closed sets, and by uniqueness (Prop.~1.5, Cor.~1.6) these are the three irreducible components, with prime ideals $(x,y)$, $(x,z)$, $(z-1,\,y-x^2)$.\hfill$\square$

\exercise{1.4}{If we identify $\mathbf A^2$ with $\mathbf A^1\times\mathbf A^1$ in the natural way, show that the Zariski topology on $\mathbf A^2$ is not the product topology of the Zariski topologies on the two copies of $\mathbf A^1$.}
\pf
The closed sets of $\mathbf A^1$ are $\mathbf A^1$ and its finite subsets, so the nonempty open sets are the cofinite subsets of $k$. The product topology has as a basis the sets $U\times V$ with $U,V$ open in $\mathbf A^1$.

Let $\Delta=\{(a,a)\}=Z(x-y)\subseteq\mathbf A^2$; it is closed in the Zariski topology of $\mathbf A^2$. Suppose $\Delta$ were closed in the product topology. Then $\mathbf A^2\setminus\Delta$ would be open, and it is nonempty ($(0,1)\notin\Delta$), so it contains a nonempty basic open set $U\times V$ with $U,V$ cofinite in $k$. But $k$ is infinite, so $U\cap V\ne\emptyset$; for $a\in U\cap V$ we get $(a,a)\in(U\times V)\cap\Delta$, a contradiction. Hence $\Delta$ is Zariski closed but not closed in the product topology, and the two topologies differ. (In fact the product topology is strictly coarser.)\hfill$\square$

\exercise{1.5}{Show that a $k$-algebra $B$ is isomorphic to the affine coordinate ring of some algebraic set in some affine space $\mathbf A^n$, if and only if $B$ is a finitely generated $k$-algebra with no nilpotent elements.}
\pf
($\Rightarrow$) Let $Y\subseteq\mathbf A^n$ be an algebraic set. $A(Y)=k[x_1,\dots,x_n]/I(Y)$ is generated by the images of the $x_i$. If $g\in k[x_1,\dots,x_n]$ and $g^m\in I(Y)$, then $g(P)^m=0$ for all $P\in Y$, so $g(P)=0$, i.e.\ $g\in I(Y)$. So $A(Y)$ has no nonzero nilpotents.

($\Leftarrow$) Let $B$ be generated by $b_1,\dots,b_n$ and reduced. The $k$-algebra map $k[x_1,\dots,x_n]\to B$, $x_i\mapsto b_i$, is surjective with kernel an ideal $\mathfrak a$, and $B\cong k[x_1,\dots,x_n]/\mathfrak a$. Since $B$ is reduced, $\mathfrak a$ is radical. Let $Y=Z(\mathfrak a)$; by the Nullstellensatz (Cor.~1.4) $I(Y)=\sqrt{\mathfrak a}=\mathfrak a$, so $B\cong A(Y)$. (If $B=0$ then $\mathfrak a=(1)$ and $Y=\emptyset$.)\hfill$\square$

\exercise{1.6}{Any nonempty open subset of an irreducible topological space is dense and irreducible. If $Y$ is a subset of a topological space $X$, which is irreducible in its induced topology, then the closure $\overline Y$ is also irreducible.}
\pf
Recall that irreducible spaces are nonempty and are not the union of two proper closed subsets.

Let $X$ be irreducible and $U\subseteq X$ nonempty open.

\emph{Dense.} $X=\overline U\cup(X\setminus U)$ is a union of two closed sets, and $X\setminus U\ne X$ since $U\ne\emptyset$. By irreducibility $\overline U=X$.

\emph{Irreducible.} $U\ne\emptyset$. Suppose $U=C_1\cup C_2$ with $C_i$ closed in $U$. Let $F_i=\overline{C_i}$ (closure in $X$), so $C_i=U\cap F_i$. Then $X=\overline U=\overline{C_1\cup C_2}=F_1\cup F_2$. By irreducibility of $X$, $F_1=X$ say, so $C_1=U\cap F_1=U$. Thus $U$ is irreducible.

Now let $Y\subseteq X$ be irreducible in the induced topology. $\overline Y\ne\emptyset$. Suppose $\overline Y=F_1\cup F_2$ with $F_i$ closed in $X$. Then $Y=(Y\cap F_1)\cup(Y\cap F_2)$ is a union of two closed subsets of $Y$, so $Y\subseteq F_i$ for some $i$, say $i=1$; then $\overline Y\subseteq F_1$, i.e.\ $F_1=\overline Y$. So $\overline Y$ is irreducible.\hfill$\square$

\exercise{1.7}{
\begin{enumerate}[label=(\alph*),leftmargin=*]
\item Show that the following conditions are equivalent for a topological space $X$: (i) $X$ is noetherian; (ii) every nonempty family of closed subsets has a minimal element; (iii) $X$ satisfies the ascending chain condition for open subsets; (iv) every nonempty family of open subsets has a maximal element.
\item A noetherian topological space is quasi-compact, i.e.\ every open cover has a finite subcover.
\item Any subset of a noetherian topological space is noetherian in its induced topology.
\item A noetherian space which is also Hausdorff must be a finite set with the discrete topology.
\end{enumerate}}
\pf
(a) Recall: $X$ is noetherian if it satisfies the descending chain condition (DCC) on closed subsets. Taking complements reverses inclusions and exchanges closed and open sets; hence (i)$\Leftrightarrow$(iii) and (ii)$\Leftrightarrow$(iv). It remains to prove (i)$\Leftrightarrow$(ii).

(ii)$\Rightarrow$(i): given $Y_1\supseteq Y_2\supseteq\cdots$ closed, the family $\{Y_j\}$ has a minimal element $Y_r$; then for $j\ge r$, $Y_j\subseteq Y_r$ and minimality gives $Y_j=Y_r$.

(i)$\Rightarrow$(ii): if a nonempty family $\mathfrak S$ of closed sets had no minimal element, choose $Y_1\in\mathfrak S$; as $Y_1$ is not minimal there is $Y_2\in\mathfrak S$ with $Y_2\subsetneq Y_1$; as $Y_2$ is not minimal there is $Y_3\subsetneq Y_2$, and so on, producing an infinite strictly descending chain, contradicting (i).

(b) Let $\{U_i\}_{i\in I}$ be an open cover of the noetherian space $X$. The family of finite unions of the $U_i$ is a nonempty family of open sets, so by (iv) it has a maximal element $V=U_{i_1}\cup\cdots\cup U_{i_s}$. For every $i$, $V\cup U_i$ is again in the family, so $V\cup U_i=V$ by maximality, i.e.\ $U_i\subseteq V$. Hence $V\supseteq\bigcup U_i=X$, and $U_{i_1},\dots,U_{i_s}$ cover $X$.

(c) Let $X$ be noetherian, $Y\subseteq X$, and $Y_1\supseteq Y_2\supseteq\cdots$ closed subsets of $Y$. Their closures $\overline{Y_j}$ in $X$ form a descending chain of closed subsets of $X$, so $\overline{Y_j}=\overline{Y_r}$ for $j\ge r$. Since $Y_j=Y\cap\overline{Y_j}$ (as $Y_j$ is closed in $Y$), $Y_j=Y_r$ for $j\ge r$.

(d) Let $X$ be noetherian and Hausdorff. By (c) every subset of $X$ is noetherian, hence quasi-compact by (b). In a Hausdorff space compact subsets are closed, so every subset of $X$ is closed. Thus $X$ carries the discrete topology. Being quasi-compact (by (b)) and discrete, $X$ is finite (the cover by singletons has a finite subcover).\hfill$\square$

\exercise{1.8}{Let $Y$ be an affine variety of dimension $r$ in $\mathbf A^n$. Let $H$ be a hypersurface in $\mathbf A^n$, and assume that $Y\not\subseteq H$. Then every irreducible component of $Y\cap H$ has dimension $r-1$. (See (7.1) for a generalization.)}
\pf
Write $H=Z(f)$ with $f\in A=k[x_1,\dots,x_n]$ (the argument works for any $f$ with $Y\not\subseteq Z(f)$). Let $\mathfrak p=I(Y)$, a prime ideal, and $B=A/\mathfrak p=A(Y)$, a domain, finitely generated over $k$, with $\dim B=\dim Y=r$ (Prop.~1.7). Let $\bar f$ be the class of $f$ in $B$. Since $Y\not\subseteq H$, $f\notin\mathfrak p$, i.e.\ $\bar f\ne0$.

Now $Y\cap H=Z(\mathfrak p+(f))$. By the correspondence of Cor.~1.4, the irreducible closed subsets of $Y\cap H$ are the sets $Z(\mathfrak q)$ with $\mathfrak q$ a prime ideal of $A$ containing $\mathfrak p+(f)$ (and $I(Z(\mathfrak q))=\mathfrak q$), inclusion-reversing. So the irreducible components of $Y\cap H$ (the maximal irreducible closed subsets) are the $Z(\mathfrak q)$ with $\mathfrak q$ minimal among primes containing $\mathfrak p+(f)$. (If $Y\cap H=\emptyset$ there is nothing to prove.) Let $\bar{\mathfrak q}=\mathfrak q/\mathfrak p$: it is a prime of $B$ minimal over $(\bar f)$.

The Noetherian domain $B$ has $\bar f\neq0$, so $\bar f$ is not a zero divisor, and $\bar f$ is not a unit because it lies in the prime $\bar{\mathfrak q}$. By Krull's Hauptidealsatz (Thm.~1.11A), $\operatorname{ht}\bar{\mathfrak q}=1$. By Thm.~1.8A(b) applied to $B$,
\[
\operatorname{ht}\bar{\mathfrak q}+\dim B/\bar{\mathfrak q}=\dim B=r,\quad\text{so}\quad \dim B/\bar{\mathfrak q}=r-1.
\]
Finally $B/\bar{\mathfrak q}=A/\mathfrak q=A(Z(\mathfrak q))$, so by Prop.~1.7 the component $Z(\mathfrak q)$ has dimension $r-1$.\hfill$\square$

\exercise{1.9}{Let $\mathfrak a\subseteq A=k[x_1,\dots,x_n]$ be an ideal which can be generated by $r$ elements. Then every irreducible component of $Z(\mathfrak a)$ has dimension $\geqslant n-r$.}
\pf
Let $\mathfrak a=(f_1,\dots,f_r)$. I prove by induction on $r\ge0$ the claim

\smallskip
$P(r)$: for any $f_1,\dots,f_r\in A$, every irreducible component of $Z(f_1,\dots,f_r)$ has dimension $\ge n-r$.

\smallskip
$r=0$: $Z(\emptyset)=\mathbf A^n$ is irreducible of dimension $n$ (Prop.~1.9).

$r-1\Rightarrow r$: Let $Z'=Z(f_1,\dots,f_{r-1})$ and $H=Z(f_r)$, so $Z(f_1,\dots,f_r)=Z'\cap H$. Let $W$ be an irreducible component of $Z'\cap H$. Being irreducible and contained in $Z'$, $W$ lies in some irreducible component $Y$ of $Z'$, which is an affine variety with $\dim Y\ge n-r+1$ by $P(r-1)$. Now $W\subseteq Y\cap H\subseteq Z'\cap H$, and $W$ is maximal among irreducible closed subsets of $Z'\cap H$, hence also of $Y\cap H$: so $W$ is an irreducible component of $Y\cap H$.
\begin{itemize}
\item If $Y\not\subseteq H$, Exercise~1.8 gives $\dim W=\dim Y-1\ge n-r$.
\item If $Y\subseteq H$, then $Y\cap H=Y$, so $W=Y$ and $\dim W\ge n-r+1\ge n-r$.
\end{itemize}
This proves $P(r)$.\hfill$\square$

\exercise{1.10}{
\begin{enumerate}[label=(\alph*),leftmargin=*]
\item If $Y$ is any subset of a topological space $X$, then $\dim Y\le\dim X$.
\item If $X$ is a topological space which is covered by a family of open subsets $\{U_i\}$, then $\dim X=\sup\dim U_i$.
\item Give an example of a topological space $X$ and a dense open subset $U$ with $\dim U<\dim X$.
\item If $Y$ is a closed subset of an irreducible finite-dimensional topological space $X$, and if $\dim Y=\dim X$, then $Y=X$.
\item Give an example of a noetherian topological space of infinite dimension.
\end{enumerate}}
\pf
(a) Let $Z_0\subsetneq Z_1\subsetneq\cdots\subsetneq Z_n$ be irreducible closed subsets of $Y$. Their closures $\overline{Z_i}$ in $X$ are irreducible (Ex.~1.6) and closed, and $\overline{Z_i}\cap Y=Z_i$ (as $Z_i$ is closed in $Y$), so the $\overline{Z_i}$ are distinct and strictly increasing. Hence $\dim X\ge n$.

(b) By (a), $\dim U_i\le\dim X$. Conversely let $Z_0\subsetneq\cdots\subsetneq Z_n$ be a chain of irreducible closed subsets of $X$. Pick $P\in Z_0$ and $i$ with $P\in U=U_i$. Then $Z_j\cap U\ne\emptyset$ for all $j$ (it contains $P$), and it is open in $Z_j$, hence irreducible (Ex.~1.6) and dense in $Z_j$; it is closed in $U$. Since $Z_j\cap U$ is dense in $Z_j$, $\overline{Z_j\cap U}=Z_j$, so the sets $Z_j\cap U$ are distinct: $Z_0\cap U\subsetneq\cdots\subsetneq Z_n\cap U$ is a chain of irreducible closed subsets of $U$. Hence $\dim U\ge n$, and $\sup\dim U_i\ge\dim X$.

(c) Let $X=\{\eta,s\}$ with the topology whose open sets are $\emptyset,\{\eta\},X$ (so the closed sets are $\emptyset,\{s\},X$; this is $\operatorname{Spec}$ of a discrete valuation ring). The closed set $\{s\}$ and $X$ are irreducible (the only proper closed subsets are $\emptyset,\{s\}$), so $\{s\}\subsetneq X$ is a chain of length $1$ and $\dim X=1$ (no chain of length $2$ exists since there are only two nonempty closed sets). The open set $U=\{\eta\}$ is dense (its closure is the smallest closed set containing $\eta$, namely $X$), and $\dim U=0$ since $U$ is a single point.

(d) Let $n=\dim Y=\dim X<\infty$, and let $Z_0\subsetneq\cdots\subsetneq Z_n$ be a chain of irreducible closed subsets of $Y$. As $Y$ is closed in $X$, each $Z_i$ is closed in $X$. If $Y\ne X$, then $Z_n\subseteq Y\subsetneq X$, and since $X$ is itself irreducible and closed, $Z_0\subsetneq\cdots\subsetneq Z_n\subsetneq X$ is a chain of length $n+1$ in $X$, contradicting $\dim X=n$. So $Y=X$.

(e) Let $X=\mathbb N=\{0,1,2,\dots\}$ with closed sets $\emptyset$, $\mathbb N$, and the initial segments $I_m=\{0,1,\dots,m\}$ ($m\ge0$). These form a topology: any nonempty intersection of the sets $I_m$ is the smallest of them, and the union of two of them is the larger one (or $\mathbb N$); $\emptyset$ and $\mathbb N$ are included. It is noetherian: a strictly descending chain of closed sets can contain $\mathbb N$ at most once as the first term, followed by finitely many $I_m$'s since $m$ strictly decreases, so every descending chain is stationary. Each $I_m$ is irreducible: its closed subsets are $\emptyset$ and $I_{m'}$, $m'\le m$, and $I_a\cup I_b=I_{\max(a,b)}$, so $I_m=C_1\cup C_2$ forces $C_1=I_m$ or $C_2=I_m$. Then $I_0\subsetneq I_1\subsetneq\cdots\subsetneq I_m$ is a chain of irreducible closed subsets of length $m$ for every $m$, so $\dim X=\infty$.

(Note that Nagata's classical example gives, more subtly, a noetherian \emph{ring} of infinite Krull dimension, i.e.\ a noetherian space $\operatorname{Spec}R$ of infinite dimension; the exercise only asks for a noetherian topological space, for which the example above suffices.)\hfill$\square$

\exercise{1.11}{Let $Y\subseteq\mathbf A^3$ be the curve given parametrically by $x=t^3$, $y=t^4$, $z=t^5$. Show that $I(Y)$ is a prime ideal of height~$2$ in $k[x,y,z]$ which cannot be generated by $2$ elements. We say that $Y$ is not a \emph{local complete intersection}---cf.\ (Ex.~2.17). (Starred in the book.)}
\pf
Let $A=k[x,y,z]$ and let $\alpha:A\to k[t]$ be the $k$-algebra map $x\mapsto t^3$, $y\mapsto t^4$, $z\mapsto t^5$; put $\mathfrak p=\ker\alpha$.

\emph{$I(Y)=\mathfrak p$.} For $g\in A$: $g\in I(Y)$ iff $g(t^3,t^4,t^5)=0$ for all $t\in k$ iff the polynomial $\alpha(g)\in k[t]$ vanishes at every point of the infinite field $k$ iff $\alpha(g)=0$. So $I(Y)=\mathfrak p$.

\emph{Prime of height $2$.} $A/\mathfrak p\cong k[t^3,t^4,t^5]\subseteq k[t]$ is a domain, so $\mathfrak p$ is prime and $Y$ is an affine variety. The fraction field of $k[t^3,t^4,t^5]$ contains $t^4/t^3=t$, so it is $k(t)$, of transcendence degree $1$. By Thm.~1.8A(a), $\dim A/\mathfrak p=1$, and by Thm.~1.8A(b) and Prop.~1.9,
\[
\operatorname{ht}\mathfrak p=\dim A-\dim A/\mathfrak p=3-1=2.
\]

\emph{Not generated by $2$ elements.} Grade $A$ by weight: $\operatorname{wt}x=3$, $\operatorname{wt}y=4$, $\operatorname{wt}z=5$, so a monomial $x^ay^bz^c$ has weight $3a+4b+5c$ and $\alpha$ sends it to $t^{3a+4b+5c}$. Write $g=\sum_wg_w$ for the decomposition into weight-homogeneous components. Then $\alpha(g)=\sum_w g_w(1,1,1)\,t^w$, where $g_w(1,1,1)$ is the sum of the coefficients of $g_w$. Hence
\[
g\in\mathfrak p\iff g_w(1,1,1)=0\ \text{for all }w.
\]
In particular $\mathfrak p$ is a homogeneous ideal, and its weight-$w$ part $\mathfrak p_w$ is the space of weight-$w$ polynomials whose coefficients sum to $0$. The monomials of weight $w$ for small $w$ are: weights $3,4,5,6,7$: only $x;\ y;\ z;\ x^2;\ xy$ respectively (and nothing in weights $1,2$); weight $8$: $y^2,xz$; weight $9$: $x^3,yz$; weight $10$: $z^2,x^2y$. (Indeed solving $3a+4b+5c=w$ in nonnegative integers gives exactly these.) Hence $\mathfrak p_w=0$ for $w<8$, and
\[
f_1=y^2-xz\ (\text{wt }8),\qquad f_2=x^3-yz\ (\text{wt }9),\qquad f_3=z^2-x^2y\ (\text{wt }10)
\]
span $\mathfrak p_8,\mathfrak p_9,\mathfrak p_{10}$ respectively.

Let $\mathfrak m=(x,y,z)$. Every element of $\mathfrak p$ has all its nonzero homogeneous components of weight $\ge8$, and every element of $\mathfrak m$ has components of weight $\ge3$; hence every element of $\mathfrak m\mathfrak p$ has all components of weight $\ge11$. Consider the $k$-vector space $\mathfrak p/\mathfrak m\mathfrak p$ (a module over $A/\mathfrak m=k$). If $c_1f_1+c_2f_2+c_3f_3\in\mathfrak m\mathfrak p$ with $c_i\in k$, then comparing components of weights $8,9,10$ (where $\mathfrak m\mathfrak p$ has nothing) gives $c_1f_1=c_2f_2=c_3f_3=0$, so $c_1=c_2=c_3=0$. Thus the images of $f_1,f_2,f_3$ are linearly independent in $\mathfrak p/\mathfrak m\mathfrak p$ and $\dim_k\mathfrak p/\mathfrak m\mathfrak p\ge3$.

If $\mathfrak p=(g_1,g_2)$, then $\mathfrak p/\mathfrak m\mathfrak p$ would be spanned over $k$ by the images of $g_1,g_2$, so would have dimension $\le2$, a contradiction. Hence $\mathfrak p=I(Y)$ cannot be generated by $2$ elements (although it has height $2$). (In fact $\mathfrak p=(f_1,f_2,f_3)$, but we do not need this.)\hfill$\square$

\exercise{1.12}{Give an example of an irreducible polynomial $f\in\mathbf R[x,y]$, whose zero set $Z(f)$ in $\mathbf A^2_{\mathbf R}$ is not irreducible (cf.\ 1.4.2).}
\pf
Here $\mathbf A^2_{\mathbf R}=\mathbf R^2$ with the Zariski topology (closed sets are common zero sets of real polynomials). Let
\[
f=(x^2-1)^2+y^2\in\mathbf R[x,y].
\]

\emph{$f$ is irreducible in $\mathbf R[x,y]$.} View $f=y^2+u^2$ with $u=x^2-1\in\mathbf R[x]$, a monic polynomial of degree $2$ in $y$ over the domain $\mathbf R[x]$. If $f=gh$ with neither factor a unit, then (degrees in $y$ add to $2$, and a factor of $y$-degree $0$ would lie in $\mathbf R[x]$ and divide the leading coefficient $1$, hence be a unit) both factors have $y$-degree $1$ with leading coefficients $a,c\in\mathbf R[x]$, $ac=1$, so $a,c$ are units, and $f$ has a root $y=r(x)\in\mathbf R[x]$. Then $r(x)^2+(x^2-1)^2=0$ in $\mathbf R[x]$; evaluating at a real $x$ with $x^2\ne1$ gives a sum of squares of reals which is $>0$, a contradiction. So $f$ is irreducible.

\emph{$Z(f)$ is reducible.} For real $(a,b)$, $f(a,b)=0$ iff $a^2=1$ and $b=0$. So $Z(f)=\{(1,0),(-1,0)\}=Z(x-1,y)\cup Z(x+1,y)$ is the union of two points, each of which is closed. Thus $Z(f)$ is a union of two proper closed subsets, hence not irreducible.

(Remark: the Nullstellensatz fails over $\mathbf R$ since $\mathbf R$ is not algebraically closed; this is why the correspondence ``irreducible $\leftrightarrow$ prime'' of Cor.~1.4 breaks down.)\hfill$\square$

\section*{2. Projective varieties}

Throughout, $k$ is algebraically closed, $S=k[x_0,\dots,x_n]$ with its grading by degree, $S_+=\bigoplus_{d>0}S_d=(x_0,\dots,x_n)$, and $\theta:\mathbf A^{n+1}\setminus\{0\}\to\mathbf P^n$ is the quotient map. For a subset $T\subseteq S$ I write $\widetilde Z(T)\subseteq\mathbf A^{n+1}$ for the \emph{affine} zero set, and for $X\subseteq\mathbf A^{n+1}$ I write $\widetilde I(X)\subseteq S$ for its ordinary ideal. Recall from the text that the sum, product, intersection and radical of homogeneous ideals are homogeneous, and that a homogeneous ideal is prime iff $fg\in\mathfrak a$ ($f,g$ homogeneous) implies $f\in\mathfrak a$ or $g\in\mathfrak a$. A homogeneous element $f$ of positive degree satisfies $f(\lambda a)=\lambda^{\deg f}f(a)$, so ``$f(P)=0$'' makes sense for $P\in\mathbf P^n$, and for $a\neq0$ in $\mathbf A^{n+1}$ we have $f(a)=0$ iff $f(\theta(a))=0$. Since $I(Y)$ is generated by homogeneous elements vanishing on $Y$, a homogeneous element of $I(Y)$ vanishes on $Y$ (it is an $S$-combination of homogeneous generators). When $Y\ne\emptyset$, $I(Y)\subseteq S_+$ (a nonzero constant does not vanish on $Y$), and $I(\emptyset)=S$.

\exercise{2.1}{Prove the ``homogeneous Nullstellensatz'', which says that if $\mathfrak a\subseteq S$ is a homogeneous ideal, and if $f\in S$ is a homogeneous polynomial with $\deg f>0$, such that $f(P)=0$ for all $P\in Z(\mathfrak a)$ in $\mathbf P^n$, then $f^q\in\mathfrak a$ for some $q>0$. [Hint: Interpret the problem in terms of the affine $(n+1)$-space whose affine coordinate ring is $S$, and use the usual Nullstellensatz (1.3A).]}
\pf
Let $\widetilde Z(\mathfrak a)\subseteq\mathbf A^{n+1}$ be the affine zero set of $\mathfrak a$. Let $a\in\widetilde Z(\mathfrak a)$, $a\neq 0$. Every homogeneous element $g\in\mathfrak a$ (of positive degree) vanishes at $a$, hence at $\theta(a)$; so $\theta(a)\in Z(\mathfrak a)$ (if $\mathfrak a$ contains a homogeneous element of degree $0$, i.e.\ a nonzero constant, then $\mathfrak a=S$ and $\widetilde Z(\mathfrak a)=\emptyset$, nothing to prove). By hypothesis $f(\theta(a))=0$, hence $f(a)=0$. Also $f(0)=0$ because $\deg f>0$. So $f$ vanishes on all of $\widetilde Z(\mathfrak a)$, and the Nullstellensatz (Thm.~1.3A) in $k[x_0,\dots,x_n]$ gives $f^q\in\mathfrak a$ for some $q>0$.\hfill$\square$

\exercise{2.2}{For a homogeneous ideal $\mathfrak a\subseteq S$, show that the following conditions are equivalent:
\begin{enumerate}[label=(\roman*),leftmargin=*]
\item $Z(\mathfrak a)=\emptyset$ (the empty set);
\item $\sqrt{\mathfrak a}$ is either $S$ or the ideal $S_+=\bigoplus_{d>0}S_d$;
\item $\mathfrak a\supseteq S_d$ for some $d>0$.
\end{enumerate}}
\pf
(i)$\Rightarrow$(ii). If $a\in\widetilde Z(\mathfrak a)$, $a\ne0$, then as in 2.1 $\theta(a)\in Z(\mathfrak a)$. So $Z(\mathfrak a)=\emptyset$ forces $\widetilde Z(\mathfrak a)\subseteq\{0\}$. If $\widetilde Z(\mathfrak a)=\emptyset$, the Nullstellensatz gives $\sqrt{\mathfrak a}=\widetilde I(\emptyset)=S$. If $\widetilde Z(\mathfrak a)=\{0\}$, then $\sqrt{\mathfrak a}=\widetilde I(\{0\})=(x_0,\dots,x_n)=S_+$.

(ii)$\Rightarrow$(iii). If $\sqrt{\mathfrak a}=S$ then $1\in\mathfrak a$ and $\mathfrak a\supseteq S_d$ for all $d$. If $\sqrt{\mathfrak a}=S_+$, then for each $i$ there is $q_i\ge1$ with $x_i^{q_i}\in\mathfrak a$. Put $d=1+\sum_i(q_i-1)>0$. A monomial $x^\alpha$ of degree $d$ has some $\alpha_i\ge q_i$ (otherwise $|\alpha|\le\sum(q_i-1)<d$), so $x^\alpha\in(x_i^{q_i})\subseteq\mathfrak a$. Hence $S_d\subseteq\mathfrak a$.

(iii)$\Rightarrow$(i). If $S_d\subseteq\mathfrak a$ then $x_i^d\in\mathfrak a$ for all $i$, so a point of $Z(\mathfrak a)$ would have all coordinates $a_i$ with $a_i^d=0$, impossible for a point of $\mathbf P^n$.\hfill$\square$

\exercise{2.3}{
\begin{enumerate}[label=(\alph*),leftmargin=*]
\item If $T_1\subseteq T_2$ are subsets of $S^h$, then $Z(T_1)\supseteq Z(T_2)$.
\item If $Y_1\subseteq Y_2$ are subsets of $\mathbf P^n$, then $I(Y_1)\supseteq I(Y_2)$.
\item For any two subsets $Y_1,Y_2$ of $\mathbf P^n$, $I(Y_1\cup Y_2)=I(Y_1)\cap I(Y_2)$.
\item If $\mathfrak a\subseteq S$ is a homogeneous ideal with $Z(\mathfrak a)\neq\emptyset$, then $I(Z(\mathfrak a))=\sqrt{\mathfrak a}$.
\item For any subset $Y\subseteq\mathbf P^n$, $Z(I(Y))=\overline Y$.
\end{enumerate}}
\pf
(a) A point annihilated by all of $T_2$ is annihilated by all of $T_1$.

(b) Every homogeneous $f$ vanishing on $Y_2$ vanishes on $Y_1$, so the generators of $I(Y_2)$ lie in $I(Y_1)$.

(c) By (b), $I(Y_1\cup Y_2)\subseteq I(Y_1)\cap I(Y_2)$. Conversely $I(Y_1)\cap I(Y_2)$ is a homogeneous ideal, hence generated by its homogeneous elements; such an element $f$ is homogeneous in $I(Y_1)$ and in $I(Y_2)$, so vanishes on $Y_1$ and on $Y_2$, hence on $Y_1\cup Y_2$, i.e.\ $f\in I(Y_1\cup Y_2)$.

(d) $\supseteq$: $\sqrt{\mathfrak a}$ is homogeneous, hence generated by homogeneous $g$ with $g^q\in\mathfrak a$ for some $q$. For $P\in Z(\mathfrak a)$, $g(P)^q=g^q(P)=0$, so $g(P)=0$; thus $g\in I(Z(\mathfrak a))$ (if $\deg g=0$ then $g$ is a constant with $g^q\in\mathfrak a$ vanishing at a point, so $g=0$; fine since $Z(\mathfrak a)\ne\emptyset$). $\subseteq$: $I(Z(\mathfrak a))$ is generated by homogeneous $f$ vanishing on $Z(\mathfrak a)$. If $\deg f>0$ then $f^q\in\mathfrak a$ by 2.1. If $\deg f=0$, $f$ is a constant vanishing at a point of the nonempty set $Z(\mathfrak a)$, so $f=0\in\sqrt{\mathfrak a}$. Hence $I(Z(\mathfrak a))\subseteq\sqrt{\mathfrak a}$.

(e) Every generator of $I(Y)$ vanishes on $Y$, so $Y\subseteq Z(I(Y))$, which is closed; hence $\overline Y\subseteq Z(I(Y))$. Conversely, let $Z(T)\supseteq Y$ be any closed set containing $Y$ ($T$ a set of homogeneous elements). Every $t\in T$ vanishes on $Y$, so $t\in I(Y)$, i.e.\ $T\subseteq I(Y)$ and by (a) $Z(I(Y))\subseteq Z(T)$. So $Z(I(Y))$ is contained in every closed set containing $Y$: $Z(I(Y))\subseteq\overline Y$. (For $Y=\emptyset$: $I=S$, $Z(S)=\emptyset$.)\hfill$\square$

\exercise{2.4}{
\begin{enumerate}[label=(\alph*),leftmargin=*]
\item There is a 1-1 inclusion-reversing correspondence between algebraic sets in $\mathbf P^n$, and homogeneous radical ideals of $S$ not equal to $S_+$, given by $Y\mapsto I(Y)$ and $\mathfrak a\mapsto Z(\mathfrak a)$. (Note: Since $S_+$ does not occur in this correspondence, it is sometimes called the \emph{irrelevant maximal ideal} of $S$.)
\item An algebraic set $Y\subseteq\mathbf P^n$ is irreducible if and only if $I(Y)$ is a prime ideal.
\item Show that $\mathbf P^n$ itself is irreducible.
\end{enumerate}}
\pf
(a) \emph{$I(Y)$ is a homogeneous radical ideal $\ne S_+$.} It is homogeneous by definition. Its radical is homogeneous and generated by homogeneous $g$ with $g^q\in I(Y)$; then $g^q$ vanishes on $Y$ (homogeneous elements of $I(Y)$ vanish on $Y$), so $g(P)=0$ for $P\in Y$, and $g\in I(Y)$ unless $\deg g=0$, in which case $g$ is a constant vanishing on $Y$, so $g=0$ or $Y=\emptyset$ (and then $I(Y)=S\ni g$). Hence $\sqrt{I(Y)}=I(Y)$. If $I(Y)=S_+$, then for algebraic $Y$ we get $Y=Z(I(Y))=Z(S_+)=\emptyset$ by 2.3(e), but $I(\emptyset)=S\ne S_+$, contradiction.

\emph{The two maps are inverse.} For $Y$ algebraic, $Z(I(Y))=\overline Y=Y$ by 2.3(e). Let $\mathfrak a$ be homogeneous radical, $\ne S_+$. If $Z(\mathfrak a)\ne\emptyset$, then $I(Z(\mathfrak a))=\sqrt{\mathfrak a}=\mathfrak a$ by 2.3(d). If $Z(\mathfrak a)=\emptyset$, then by 2.2 $\mathfrak a=\sqrt{\mathfrak a}\in\{S,S_+\}$, so $\mathfrak a=S=I(\emptyset)$. Both maps reverse inclusions by 2.3(a),(b). So we have the bijection.

(b) Let $Y$ be algebraic. Suppose $Y$ irreducible (so $Y\neq\emptyset$, $I(Y)\ne S$). Let $f,g$ be homogeneous with $fg\in I(Y)$. If one of them has degree $0$, it is a nonzero constant or $0$ and the claim is trivial. Otherwise $Y\subseteq Z(fg)=Z(f)\cup Z(g)$, so $Y=(Y\cap Z(f))\cup(Y\cap Z(g))$ is a union of two closed subsets, hence $Y\subseteq Z(f)$ or $Y\subseteq Z(g)$, i.e.\ $f\in I(Y)$ or $g\in I(Y)$. So $I(Y)$ is prime. Conversely let $I(Y)$ be prime; then $I(Y)\neq S$ so $Y\neq\emptyset$. If $Y=Y_1\cup Y_2$ with $Y_i$ closed, then $I(Y)=I(Y_1)\cap I(Y_2)\supseteq I(Y_1)I(Y_2)$ (2.3(c)), so by primality $I(Y)\supseteq I(Y_1)$, say; with 2.3(b) $I(Y)=I(Y_1)$, so $Y=Z(I(Y))=Z(I(Y_1))=Y_1$. Thus $Y$ is irreducible.

(c) $\mathbf P^n=Z(0)$ is algebraic. If a homogeneous $f$ vanishes at all points of $\mathbf P^n$, then it vanishes on all of $k^{n+1}$ (nonzero points by homogeneity, and $0$ if $\deg f>0$), and since $k$ is infinite, $f=0$. So $I(\mathbf P^n)=0$, a prime ideal as $S$ is a domain. By (b) $\mathbf P^n$ is irreducible.\hfill$\square$

\exercise{2.5}{
\begin{enumerate}[label=(\alph*),leftmargin=*]
\item $\mathbf P^n$ is a noetherian topological space.
\item Every algebraic set in $\mathbf P^n$ can be written uniquely as a finite union of irreducible algebraic sets, no one containing another. These are called its \emph{irreducible components}.
\end{enumerate}}
\pf
(a) Let $Y_1\supseteq Y_2\supseteq\cdots$ be a descending chain of closed subsets of $\mathbf P^n$. Then $I(Y_1)\subseteq I(Y_2)\subseteq\cdots$ (2.3(b)) is an ascending chain of ideals of the noetherian ring $S$, so it is stationary, say $I(Y_m)=I(Y_{m+1})=\cdots$. Since $Y_j=Z(I(Y_j))$ (2.3(e)) the chain $Y_j$ is stationary too. By Ex.~1.7, $\mathbf P^n$ is noetherian.

(b) \emph{Existence.} Suppose the family of closed sets which are not finite unions of irreducible closed sets is nonempty. By (a) and Ex.~1.7 it has a minimal member $Y$. Then $Y\ne\emptyset$ (empty union) and $Y$ is not irreducible, so $Y=Y_1\cup Y_2$ with $Y_i\subsetneq Y$ closed; by minimality both $Y_i$ are finite unions of irreducible closed sets, hence so is $Y$, a contradiction. Discarding members contained in others gives a representation with no member containing another. (Closed subsets of an algebraic set are algebraic.)

\emph{Uniqueness.} Let $Y=Y_1\cup\cdots\cup Y_r=Z_1\cup\cdots\cup Z_s$ be two such representations. Then $Y_i=\bigcup_j(Y_i\cap Z_j)$ with $Y_i$ irreducible, so $Y_i\subseteq Z_j$ for some $j$; likewise $Z_j\subseteq Y_{i'}$ for some $i'$. So $Y_i\subseteq Y_{i'}$, whence $i=i'$ and $Y_i=Z_j$. Thus every $Y_i$ is some $Z_j$, and symmetrically; the two lists coincide.\hfill$\square$

\exercise{2.6}{If $Y$ is a projective variety with homogeneous coordinate ring $S(Y)$, show that $\dim S(Y)=\dim Y+1$. [Hint: Let $\varphi_i:U_i\to\mathbf A^n$ be the homeomorphism of (2.2), let $Y_i$ be the affine variety $\varphi_i(Y\cap U_i)$, and let $A(Y_i)$ be its affine coordinate ring. Show that $A(Y_i)$ can be identified with the subring of elements of degree $0$ of the localized ring $S(Y)_{x_i}$. Then show that $S(Y)_{x_i}\cong A(Y_i)[x_i,x_i^{-1}]$. Now use (1.7), (1.8A), and (Ex.~1.10), and look at transcendence degrees. Conclude also that $\dim Y=\dim Y_i$ whenever $Y_i$ is nonempty.]}
\pf
Here $\dim S(Y)$ is the Krull dimension. $Y$ is nonempty and irreducible, so $\mathfrak p=I(Y)$ is a homogeneous prime (2.4(b)) and $S(Y)=S/\mathfrak p$ is a finitely generated $k$-domain.

\emph{Set-up.} Choose $i$ with $Y\cap U_i\neq\emptyset$, i.e.\ $x_i\notin I(Y)$ (such $i$ exists since $Y\ne\emptyset$); the argument works for every such $i$, and I take $i=0$ for notation. Put $A=k[y_1,\dots,y_n]$, $y_j=x_j/x_0$. Then $Y_0=\varphi_0(Y\cap U_0)$ is closed in $\mathbf A^n$ (as $Y\cap U_0$ is closed in $U_0$ and $\varphi_0$ is a homeomorphism, Prop.~2.2), and irreducible (a nonempty open subset of the irreducible $Y$ is irreducible by Ex.~1.6, and $\varphi_0$ is a homeomorphism). So $Y_0$ is an affine variety with coordinate ring $A(Y_0)=A/I(Y_0)$ a domain.

\emph{Claim:} $S(Y)_{x_0}\cong A(Y_0)[x_0,x_0^{-1}]$, with $x_0$ transcendental over $A(Y_0)$. We have $S_{x_0}=k[x_0^{\pm1},x_1,\dots,x_n]=A[x_0,x_0^{-1}]$ (Laurent polynomial ring over $A$, $x_j=x_0y_j$), and $S(Y)_{x_0}=S_{x_0}/I(Y)S_{x_0}$. I show $I(Y)S_{x_0}=I(Y_0)S_{x_0}$. Write $\alpha(f)=f(1,y_1,\dots,y_n)$ and $\beta(g)=x_0^{e}g(x_1/x_0,\dots,x_n/x_0)$ for $g\in A$ of degree $e$ (the maps of Prop.~2.2).
\begin{itemize}
\item[$\subseteq$:] a generator $f$ of $I(Y)$ is homogeneous of some degree $d$ and $f=x_0^d\alpha(f)(y)$ in $S_{x_0}$. For $P=(1,a_1,\dots,a_n)\in Y\cap U_0$ we have $\alpha(f)(a)=f(P)=0$, so $\alpha(f)\in I(Y_0)$.
\item[$\supseteq$:] for $g\in I(Y_0)$, $\beta(g)$ is homogeneous and for $P=(a_0,\dots,a_n)\in Y\cap U_0$, $\beta(g)(P)=a_0^eg(a_1/a_0,\dots)=0$; so $Z(\beta(g))$ is closed in $\mathbf P^n$ and contains $Y\cap U_0$, which is dense in $Y$ ($Y$ irreducible, Ex.~1.6), hence $Z(\beta(g))\supseteq Y$. Thus the homogeneous element $\beta(g)$ vanishes on $Y$, so $\beta(g)\in I(Y)$ and $g=\beta(g)/x_0^e\in I(Y)S_{x_0}$.
\end{itemize}
Hence $S(Y)_{x_0}\cong A[x_0^{\pm1}]/I(Y_0)A[x_0^{\pm1}]=A(Y_0)[x_0^{\pm1}]$. Its degree-$0$ part (for the grading in which $x_0$ is a unit of degree~$1$ and $x_j$ has degree~$1$) is exactly $A(Y_0)$, which gives the identification in the hint.

\emph{Transcendence degrees.} $S(Y)$ is a domain in which $x_0\ne0$, so $\operatorname{Frac}S(Y)=\operatorname{Frac}S(Y)_{x_0}=K(x_0)$ with $K=\operatorname{Frac}A(Y_0)$ and $x_0$ transcendental over $K$. By Thm.~1.8A(a) and Prop.~1.7,
\[
\dim S(Y)=\operatorname{tr.deg}_k K(x_0)=\operatorname{tr.deg}_kK+1=\dim A(Y_0)+1=\dim Y_0+1 .
\]
So $\dim Y_i=\dim S(Y)-1$ for every $i$ with $Y_i\ne\emptyset$. Finally the sets $Y\cap U_i$ (those nonempty) form an open cover of $Y$ and are homeomorphic to the $Y_i$; by Ex.~1.10(b), $\dim Y=\sup_i\dim(Y\cap U_i)=\dim S(Y)-1$. Thus $\dim S(Y)=\dim Y+1$, and $\dim Y=\dim Y_i$ whenever $Y_i\ne\emptyset$.\hfill$\square$

\exercise{2.7}{
\begin{enumerate}[label=(\alph*),leftmargin=*]
\item $\dim\mathbf P^n=n$.
\item If $Y\subseteq\mathbf P^n$ is a quasi-projective variety, then $\dim Y=\dim\overline Y$. [Hint: Use (Ex.~2.6) to reduce to (1.10).]
\end{enumerate}}
\pf
(a) $\mathbf P^n$ is a projective variety (2.4(c)) with $I(\mathbf P^n)=0$, so $S(\mathbf P^n)=S$, and $\dim S=n+1$ by Prop.~1.9 applied to $k[x_0,\dots,x_n]$. By 2.6, $\dim\mathbf P^n=n$.

(b) By definition $Y$ is a nonempty open subset of a projective variety $W$; $Y$ is irreducible, and by Ex.~1.6 so is $\overline Y\subseteq W$, so $\overline Y$ is a projective variety; and $Y=Y\cap\overline Y$ is open in $\overline Y$. Let $i$ be such that $Y\cap U_i\ne\emptyset$. The set $Z_i=\varphi_i(\overline Y\cap U_i)$ is closed in $\mathbf A^n$ and is an affine variety (a nonempty open subset of $\overline Y$, hence irreducible). Further $\varphi_i(Y\cap U_i)$ is open in $Z_i$ and nonempty, hence dense in $Z_i$; so $\varphi_i(Y\cap U_i)$ is a quasi-affine variety whose closure is $Z_i$. By Prop.~1.10,
\[\dim (Y\cap U_i)=\dim\varphi_i(Y\cap U_i)=\dim Z_i,\]
and by 2.6 applied to the projective variety $\overline Y$, $\dim Z_i=\dim\overline Y$. The nonempty sets $Y\cap U_i$ cover $Y$, so Ex.~1.10(b) gives $\dim Y=\sup_i\dim(Y\cap U_i)=\dim\overline Y$.\hfill$\square$

\exercise{2.8}{A projective variety $Y\subseteq\mathbf P^n$ has dimension $n-1$ if and only if it is the zero set of a single irreducible homogeneous polynomial $f$ of positive degree. $Y$ is called a \emph{hypersurface} in $\mathbf P^n$.}
\pf
(Here $n\ge1$. Irreducibility of $Y$ is essential: $Z(x_0x_1)$ has dimension $n-1$ but is not the zero set of an irreducible polynomial; the statement concerns projective varieties $Y$, i.e.\ irreducible algebraic sets.)

($\Leftarrow$) Let $f$ be irreducible homogeneous of positive degree $e$, $Y=Z(f)$. $S$ is a UFD, so $(f)$ is prime. \emph{$Y\ne\emptyset$:} otherwise by 2.2(iii) $(f)\supseteq S_d$ for some $d>0$, so $f\mid x_0^d$ and $f\mid x_1^d$; unique factorization gives $f=cx_0^a=c'x_1^b$ with $a,b\ge1$, absurd. Hence by 2.3(d), $I(Y)=\sqrt{(f)}=(f)$, prime, so $Y$ is a projective variety (2.4(b)). Now $\mathfrak p=(f)$ is a prime of height $1$ (Krull's Hauptidealsatz, Thm.~1.11A: it is minimal over $(f)$), so by Thm.~1.8A(b), $\dim S/(f)=\dim S-1=n$. By 2.6, $\dim Y=\dim S(Y)-1=n-1$.

($\Rightarrow$) Let $Y$ be a projective variety of dimension $n-1$. $\mathfrak p=I(Y)$ is a homogeneous prime and by 2.6, $\dim S/\mathfrak p=n$. By Thm.~1.8A(b), $\operatorname{ht}\mathfrak p=\dim S-n=1$. As $S$ is a UFD, $\mathfrak p=(f)$ is principal (Prop.~1.12A), with $f$ irreducible. \emph{We may take $f$ homogeneous:} write $f=f_a+\dots+f_b$ with $f_b\ne0$ the top component. As $\mathfrak p$ is homogeneous, $f_b\in(f)$, say $f_b=fg$, $g\neq0$. Comparing top-degree parts, $\deg f_b=b+\deg g$, so $\deg g=0$, $g$ a constant and $f=g^{-1}f_b$ is homogeneous. Since $\mathfrak p\ne S$, $\deg f>0$. Then $Y=Z(I(Y))=Z(f)$ by 2.3(e).\hfill$\square$

\exercise{2.9}{\emph{Projective closure of an affine variety.} If $Y\subseteq\mathbf A^n$ is an affine variety, we identify $\mathbf A^n$ with an open set $U_0\subseteq\mathbf P^n$ by the homeomorphism $\varphi_0$. Then we can speak of $\overline Y$, the closure of $Y$ in $\mathbf P^n$, which is called the \emph{projective closure} of $Y$.
\begin{enumerate}[label=(\alph*),leftmargin=*]
\item Show that $I(\overline Y)$ is the ideal generated by $\beta(I(Y))$, using the notation of the proof of (2.2).
\item Let $Y\subseteq\mathbf A^3$ be the twisted cubic of (Ex.~1.2). Its projective closure $\overline Y\subseteq\mathbf P^3$ is called the \emph{twisted cubic curve} in $\mathbf P^3$. Find generators for $I(Y)$ and $I(\overline Y)$, and use this example to show that if $f_1,\dots,f_r$ generate $I(Y)$, then $\beta(f_1),\dots,\beta(f_r)$ do not necessarily generate $I(\overline Y)$.
\end{enumerate}}
\pf
Notation (Prop.~2.2): $A=k[y_1,\dots,y_n]$; for $g\in A$ of degree $e$, $\beta(g)=x_0^eg(x_1/x_0,\dots,x_n/x_0)\in S_e$; for $f\in S$, $\alpha(f)=f(1,y_1,\dots,y_n)$. Note $\beta(gh)=\beta(g)\beta(h)$.

(a) (This works for any subset $Y\subseteq\mathbf A^n=U_0$.) Let $J$ be the ideal of $S$ generated by $\beta(I(Y))$, where $I(Y)\subseteq A$.

$J\subseteq I(\overline Y)$: for $g\in I(Y)$ and $P=(a_0,\dots,a_n)\in U_0$ with image $Q=\varphi_0(P)\in Y$, $\beta(g)(P)=a_0^eg(a_1/a_0,\dots)=a_0^eg(Q)=0$. So $\beta(g)$ vanishes on $Y$ and hence on the closed set $Z(\beta(g))\supseteq Y$, that is on $\overline Y$.

$I(\overline Y)\subseteq J$: $I(\overline Y)$ is generated by homogeneous $f$ vanishing on $\overline Y$, so it suffices to see $f\in J$ for such $f\neq0$. Write $f=x_0^mf'$ with $x_0\nmid f'$, $f'$ homogeneous of degree $d'$. Since $x_0\nmid f'$, $f'$ has a monomial not involving $x_0$; the map $x^\alpha\mapsto y^{(\alpha_1,\dots,\alpha_n)}$ is injective on monomials of degree $d'$, so $\alpha(f')$ has degree exactly $d'$ and $\beta(\alpha(f'))=f'$. For $Q=(a_1,\dots,a_n)\in Y$, $\alpha(f')(Q)=f'(1,a)=0$ because $f$ vanishes at $(1,a)\in\overline Y$ and $1\ne0$. So $\alpha(f')\in I(Y)$, and $f=x_0^m\beta(\alpha(f'))\in J$.

(b) $Y=\{(t,t^2,t^3)\}\subseteq\mathbf A^3$ with coordinates $x,y,z$; by Ex.~1.2, $I(Y)=(f_1,f_2)$ with $f_1=y-x^2$, $f_2=z-x^3$. Let $x_0,\dots,x_3$ be homogeneous coordinates on $\mathbf P^3$ with $x=x_1/x_0$, $y=x_2/x_0$, $z=x_3/x_0$. Then $\beta(f_1)=x_0x_2-x_1^2$, $\beta(f_2)=x_0^2x_3-x_1^3$. Put
\[
q_1=x_0x_2-x_1^2,\qquad q_2=x_0x_3-x_1x_2,\qquad q_3=x_1x_3-x_2^2 .
\]
\emph{Claim: $I(\overline Y)=(q_1,q_2,q_3)=:J'$.} Each $q_i$ vanishes at all points $(1,t,t^2,t^3)$ of $Y$ (e.g.\ $q_2(1,t,t^2,t^3)=t^3-t\cdot t^2=0$), hence on $\overline Y$; so $J'\subseteq I(\overline Y)$. Conversely, by $q_1,q_2,q_3$ we have the congruences $x_1^2\equiv x_0x_2$, $x_1x_2\equiv x_0x_3$, $x_2^2\equiv x_1x_3$ mod $J'$. Call a monomial $x_0^ax_1^bx_2^cx_3^d$ \emph{reduced} if $b+c\le1$. If $b+c\ge2$ the monomial contains $x_1^2$, $x_1x_2$ or $x_2^2$, and replacing it by $x_0x_2$, $x_0x_3$, $x_1x_3$ respectively changes the monomial into another monomial of the same degree (modulo $J'$) with smaller value of $b+c$ (it drops by $1$, $2$, $1$). By induction on $b+c$, every homogeneous polynomial of degree $m$ is congruent mod $J'$ to a linear combination of reduced monomials of degree $m$. These are $x_0^ax_3^d$ ($a+d=m$), $x_0^ax_1x_3^d$ and $x_0^ax_2x_3^d$ ($a+d=m-1$); evaluated at $(1,t,t^2,t^3)$ they give the powers $t^{3d}$, $t^{3d+1}$, $t^{3d+2}$ respectively, which are pairwise distinct powers $t^0,\dots,t^{3m}$. Now let $f\in I(\overline Y)$ be homogeneous of degree $m$ and $f\equiv r$ mod $J'$ with $r$ a combination of reduced monomials. Then $r\in I(\overline Y)$, so $r(1,t,t^2,t^3)=0$ for all $t\in k$; this is a polynomial in $t$ with infinitely many zeros, so it is $0$, and since the powers are distinct all coefficients of $r$ vanish: $r=0$, $f\in J'$. Thus $I(\overline Y)=J'$.

\emph{$\beta(f_1),\beta(f_2)$ do not generate $I(\overline Y)$.} The ideal $(\beta(f_1),\beta(f_2))=(q_1,\,x_0^2x_3-x_1^3)$ is generated by homogeneous elements of degrees $2$ and $3$, so its homogeneous elements of degree $2$ are $c\,q_1$ ($c\in k$). But $x_1x_3-x_2^2=q_3\in I(\overline Y)$ is of degree $2$ and is not a multiple of $q_1$. Hence $(\beta(f_1),\beta(f_2))\subsetneq I(\overline Y)$.\hfill$\square$

\exercise{2.10}{\emph{The Cone Over a Projective Variety.} Let $Y\subseteq\mathbf P^n$ be a nonempty algebraic set, and let $\theta:\mathbf A^{n+1}\setminus\{(0,\dots,0)\}\to\mathbf P^n$ be the map which sends the point with affine coordinates $(a_0,\dots,a_n)$ to the point with homogeneous coordinates $(a_0,\dots,a_n)$. We define the \emph{affine cone} over $Y$ to be $C(Y)=\theta^{-1}(Y)\cup\{(0,\dots,0)\}$.
\begin{enumerate}[label=(\alph*),leftmargin=*]
\item Show that $C(Y)$ is an algebraic set in $\mathbf A^{n+1}$, whose ideal is equal to $I(Y)$, considered as an ordinary ideal in $k[x_0,\dots,x_n]$.
\item $C(Y)$ is irreducible if and only if $Y$ is.
\item $\dim C(Y)=\dim Y+1$.
\end{enumerate}
Sometimes we consider the projective closure $\overline{C(Y)}$ of $C(Y)$ in $\mathbf P^{n+1}$. This is called the \emph{projective cone} over $Y$.}
\pf
(a) Let $\mathfrak a=I(Y)$. Since $Y\ne\emptyset$, $\mathfrak a\subseteq S_+$, so $0\in\widetilde Z(\mathfrak a)$. For $a\ne0$: $a\in\widetilde Z(\mathfrak a)$ iff every homogeneous element of $\mathfrak a$ vanishes at $a$ iff $\theta(a)\in Z(\mathfrak a)=Z(I(Y))=Y$ (2.3(e)). Hence $\widetilde Z(\mathfrak a)=C(Y)$, an algebraic set. By the Nullstellensatz (Thm.~1.3A), $\widetilde I(C(Y))=\sqrt{\mathfrak a}$, which equals $\mathfrak a=I(Y)$ by 2.4(a) ($I(Y)$ is radical).

(b) By Cor.~1.4, $C(Y)$ is irreducible iff $\widetilde I(C(Y))=I(Y)$ is prime, iff $Y$ is irreducible (2.4(b)).

(c) First let $Y$ be irreducible. Then $C(Y)$ is an affine variety with $A(C(Y))=S/I(Y)=S(Y)$, so by Prop.~1.7 and 2.6, $\dim C(Y)=\dim S(Y)=\dim Y+1$. In general write $Y=Y_1\cup\dots\cup Y_s$ with $Y_j$ irreducible components (2.5(b)); then $C(Y)=\bigcup C(Y_j)$, each $C(Y_j)$ closed and irreducible by (a),(b). In a finite union $X=\bigcup X_j$ of closed sets, $\dim X=\max\dim X_j$: a chain of irreducible closed subsets $Z_0\subsetneq\dots\subsetneq Z_m$ in $X$ has $Z_m=\bigcup(Z_m\cap X_j)$, so $Z_m\subseteq X_j$ for some $j$, and then the whole chain lies in $X_j$. Hence $\dim C(Y)=\max_j(\dim Y_j+1)=\dim Y+1$.\hfill$\square$

\exercise{2.11}{\emph{Linear Varieties in $\mathbf P^n$.} A hypersurface defined by a linear polynomial is called a \emph{hyperplane}.
\begin{enumerate}[label=(\alph*),leftmargin=*]
\item Show that the following two conditions are equivalent for a variety $Y$ in $\mathbf P^n$: (i) $I(Y)$ can be generated by linear polynomials; (ii) $Y$ can be written as an intersection of hyperplanes. In this case we say that $Y$ is a \emph{linear variety} in $\mathbf P^n$.
\item If $Y$ is a linear variety of dimension $r$ in $\mathbf P^n$, show that $I(Y)$ is minimally generated by $n-r$ linear polynomials.
\item Let $Y,Z$ be linear varieties in $\mathbf P^n$, with $\dim Y=r$, $\dim Z=s$. If $r+s-n\ge0$, then $Y\cap Z\ne\emptyset$. Furthermore, if $Y\cap Z\ne\emptyset$, then $Y\cap Z$ is a linear variety of dimension $\ge r+s-n$. (Think of $\mathbf A^{n+1}$ as a vector space over $k$, and work with subspaces.)
\end{enumerate}}
\pf
``Linear polynomial'' means a nonzero linear form (homogeneous of degree~1).

\textbf{Lemma.} Let $W\subseteq k^{n+1}$ be a linear subspace of dimension $d+1\ge1$, $\mathbf P(W)=\theta(W\setminus0)$, and let $l_1,\dots,l_{m}$ ($m=n-d$) be a basis of the space of linear forms vanishing on $W$. Then $\mathbf P(W)=Z(l_1,\dots,l_m)$, $I(\mathbf P(W))=(l_1,\dots,l_m)$, $\mathbf P(W)$ is a projective variety and $\dim\mathbf P(W)=d$.

\emph{Proof.} Complete $l_1,\dots,l_m$ to a basis $l_1,\dots,l_{n+1}$ of the space of linear forms. Then $S=k[l_1,\dots,l_{n+1}]$ is a polynomial ring in these new coordinates, and the notions of homogeneous, degree, zero set are the same in the new coordinates. $W$ is the common zero set of $l_1,\dots,l_m$ in $k^{n+1}$ (dimension count), so $\mathbf P(W)=Z(l_1,\dots,l_m)\ne\emptyset$. The ideal $\mathfrak q=(l_1,\dots,l_m)$ is homogeneous and prime ($S/\mathfrak q=k[l_{m+1},\dots,l_{n+1}]$ is a domain), so by 2.3(d), $I(\mathbf P(W))=\sqrt{\mathfrak q}=\mathfrak q$, and $\mathbf P(W)$ is a variety (2.4(b)) with $S(\mathbf P(W))=k[l_{m+1},\dots,l_{n+1}]$ of dimension $n+1-m=d+1$ (Prop.~1.9). By 2.6, $\dim\mathbf P(W)=d$.\hfill$\square$

(a) (i)$\Rightarrow$(ii): if $I(Y)=(l_1,\dots,l_s)$ with the $l_i$ linear (nonzero), then $Y=Z(I(Y))=\bigcap Z(l_i)$ by 2.3(e), an intersection of hyperplanes. (ii)$\Rightarrow$(i): let $Y=\bigcap_{\lambda}Z(l_\lambda)$, and let $V$ be the span of the $l_\lambda$ in $S_1$, with basis $l_1,\dots,l_m$. Then $Y=Z(l_1,\dots,l_m)=\mathbf P(W)$ with $W=\{l_1=\dots=l_m=0\}$, and $\dim W=n+1-m\ge1$ since $Y\ne\emptyset$ (a variety). By the Lemma, $I(Y)=(l_1,\dots,l_m)$.

(b) By (a) and the Lemma, $Y=\mathbf P(W)$ with $\dim W=r+1$ and $I(Y)=(l_1,\dots,l_{n-r})$ with $l_i$ linearly independent (the Lemma with $d=r$). Suppose $I=I(Y)$ were generated by $t$ elements $g_1,\dots,g_t$ (of any degrees). Since $S=k+S_+$, $I=\sum Sg_j=\sum kg_j+S_+I$, so the images of $g_j$ span the $k$-vector space $I/S_+I$. On the other hand the images of $l_1,\dots,l_{n-r}$ in $I/S_+I$ are linearly independent: $S_+I$ lives in degrees $\ge2$, so a relation $\sum c_il_i\in S_+I$ of degree~1 forces $\sum c_il_i=0$, and $c_i=0$. Thus $t\ge n-r$, and $I(Y)$ is minimally generated by the $n-r$ linear polynomials $l_i$.

(c) By (a) and the Lemma, $Y=\mathbf P(W_Y)$, $Z=\mathbf P(W_Z)$ with subspaces $W_Y,W_Z\subseteq k^{n+1}$ of dimensions $r+1$, $s+1$. Then $Y\cap Z=\mathbf P(W_Y\cap W_Z)$ (where $\mathbf P(0)=\emptyset$). Now
\[
\dim(W_Y\cap W_Z)\ge(r+1)+(s+1)-(n+1)=r+s-n+1 .
\]
If $r+s-n\ge0$ this is $\ge1$, so $W_Y\cap W_Z$ contains a nonzero vector and $Y\cap Z\ne\emptyset$. If $Y\cap Z\ne\emptyset$, then $Y\cap Z=\mathbf P(W)$ with $W=W_Y\cap W_Z$ of dimension $d+1\ge r+s-n+1$; by the Lemma it is a variety with $I(Y\cap Z)$ generated by linear forms, i.e.\ a linear variety, of dimension $d\ge r+s-n$.\hfill$\square$

\exercise{2.12}{\emph{The $d$-Uple Embedding.} For given $n,d>0$, let $M_0,M_1,\dots,M_N$ be all the monomials of degree $d$ in the $n+1$ variables $x_0,\dots,x_n$, where $N=\binom{n+d}{n}-1$. We define a mapping $\rho_d:\mathbf P^n\to\mathbf P^N$ by sending the point $P=(a_0,\dots,a_n)$ to the point $\rho_d(P)=(M_0(a),\dots,M_N(a))$ obtained by substituting the $a_i$ in the monomials $M_j$. This is called the $d$-uple embedding of $\mathbf P^n$ in $\mathbf P^N$. For example, if $n=1$, $d=2$, then $N=2$, and the image $Y$ of the 2-uple embedding of $\mathbf P^1$ in $\mathbf P^2$ is a conic.
\begin{enumerate}[label=(\alph*),leftmargin=*]
\item Let $\theta:k[y_0,\dots,y_N]\to k[x_0,\dots,x_n]$ be the homomorphism defined by sending $y_i$ to $M_i$, and let $\mathfrak a$ be the kernel of $\theta$. Then $\mathfrak a$ is a homogeneous prime ideal, and so $Z(\mathfrak a)$ is a projective variety in $\mathbf P^N$.
\item Show that the image of $\rho_d$ is exactly $Z(\mathfrak a)$. (One inclusion is easy. The other will require some calculation.)
\item Now show that $\rho_d$ is a homeomorphism of $\mathbf P^n$ onto the projective variety $Z(\mathfrak a)$.
\item Show that the twisted cubic curve in $\mathbf P^3$ (Ex.~2.9) is equal to the 3-uple embedding of $\mathbf P^1$ in $\mathbf P^3$, for suitable choice of coordinates.
\end{enumerate}}
\pf
$\rho_d$ is well defined: $M_j(\lambda a)=\lambda^dM_j(a)$, and not all $M_j(a)$ vanish because $x_i^d$ are among the $M_j$. Index monomials by exponent vectors $\alpha\in\mathbf N^{n+1}$, $|\alpha|=d$, writing $y_\alpha\mapsto x^\alpha$ and $e_i$ for the unit vectors. (This $\theta$ is unrelated to the quotient map $\mathbf A^{n+1}\setminus0\to\mathbf P^n$, which we do not use here.)

(a) $k[y_0,\dots,y_N]/\mathfrak a\cong\operatorname{im}\theta\subseteq k[x_0,\dots,x_n]$ is a nonzero subring of a domain, so $\mathfrak a$ is a prime ideal ($\ne$ the unit ideal). It is homogeneous: if $g=\sum_eg_e$ ($g_e$ homogeneous of degree $e$ in the $y$'s), then $\theta(g_e)$ is homogeneous of degree $de$, so $\theta(g)=0$ iff every $\theta(g_e)=0$; hence the homogeneous components of elements of $\mathfrak a$ lie in $\mathfrak a$. That $Z(\mathfrak a)$ is a projective variety is seen after (b): it is nonempty, and $I(Z(\mathfrak a))=\sqrt{\mathfrak a}=\mathfrak a$ (2.3(d)) is prime, so $Z(\mathfrak a)$ is irreducible by 2.4(b).

(b) $\subseteq$: for $P=(a_0,\dots,a_n)$ and homogeneous $g\in\mathfrak a$, $g(\rho_d(P))=g(M_0(a),\dots,M_N(a))=\theta(g)(a)=0$.

$\supseteq$: let $Q=(b_\alpha)\in Z(\mathfrak a)$. For each $\alpha$, $y_\alpha^d-\prod_iy_{de_i}^{\alpha_i}\in\mathfrak a$ (both terms map to $x^{d\alpha}$). Some $b_\alpha\ne0$, so $\prod_ib_{de_i}^{\alpha_i}\ne0$, so $b_{de_i}\ne0$ for some $i$. Scale $Q$ so that $b_{de_i}=1$. Define $P=(a_0,\dots,a_n)$ by $a_i=1$ and $a_j=b_{(d-1)e_i+e_j}$ for $j\ne i$. For any $\alpha$ with $|\alpha|=d$ the binomial
\[
y_\alpha\,y_{de_i}^{d-1}-\prod_jy_{(d-1)e_i+e_j}^{\alpha_j}
\]
is homogeneous (of degree $d$ in the $y$'s) and maps under $\theta$ to $x^\alpha x_i^{d(d-1)}-\prod_j(x_i^{d-1}x_j)^{\alpha_j}=0$. So it lies in $\mathfrak a$ and vanishes at $Q$: $b_\alpha\cdot1=\prod_jb_{(d-1)e_i+e_j}^{\alpha_j}=a^\alpha$ (note $b_{(d-1)e_i+e_i}=b_{de_i}=1=a_i$). Hence $\rho_d(P)=Q$.

(c) \emph{Injective.} Suppose $\rho_d(a)=\rho_d(a')$, so $M_\alpha(a')=\lambda M_\alpha(a)$ for all $\alpha$, some $\lambda\neq0$. Choosing $\mu$ with $\mu^d=\lambda$ and replacing $a'$ by $\mu^{-1}a'$ we may assume $\lambda=1$. Pick $i$ with $a_i\ne0$. Then $a_i'^d=a_i^d\ne0$, so $a_i'=\zeta a_i$ with $\zeta^d=1$; and $a'_ja_i'^{d-1}=a_ja_i^{d-1}$ gives $a'_j=\zeta^{-(d-1)}a_j=\zeta a_j$. So $a'=\zeta a$, the same point of $\mathbf P^n$. With (b), $\rho_d:\mathbf P^n\to Y:=Z(\mathfrak a)$ is bijective.

\emph{Continuous.} If $C=Z(T')\subseteq Y$ with $T'$ homogeneous in $k[y]$, then $\rho_d^{-1}(C)=Z(\theta(T'))$ is closed (for $g\in T'$, $g(\rho_d(P))=\theta(g)(P)$).

\emph{Closed.} Let $C\subseteq\mathbf P^n$ be closed, $C=Z(f_1,\dots,f_r)$ with $f_j$ homogeneous (finite set by noetherianity). Choose $m_0\ge1$ with $dm_0\ge\deg f_j$ for all $j$, and let $T'$ be the set of homogeneous elements of degree $dm$, $m\ge m_0$, in the ideal $(f_1,\dots,f_r)$. \emph{Claim $Z(T')=C$.} $\supseteq$ is clear. If $P\in Z(T')$, pick $i$ with $a_i\ne0$; then $x_i^{dm_0-\deg f_j}f_j\in T'$ vanishes at $P$, so $f_j(P)=0$. Each monomial of degree $dm$ is a product of $m$ monomials of degree $d$ (split the exponent vector), so each $h\in T'$ can be written $h=\theta(H_h)$ with $H_h$ homogeneous of degree $m$ in the $y$'s. Then
\[
\rho_d(C)=Y\cap Z(\{H_h:h\in T'\}):
\]
$\subseteq$: $H_h(\rho_d(P))=h(P)=0$ for $P\in C=Z(T')$; $\supseteq$: a point $Q$ of the right side is $\rho_d(P)$ by (b) and $h(P)=H_h(Q)=0$ for all $h$, so $P\in Z(T')=C$. So $\rho_d(C)$ is closed in $Y$.

A continuous closed bijection is a homeomorphism: $\rho_d:\mathbf P^n\to Y$ is a homeomorphism, and $Y$ is a projective variety.

(d) For $n=1$, $d=3$: $N=3$, take $M_0,\dots,M_3=x_0^3,x_0^2x_1,x_0x_1^2,x_1^3$, so $\rho_3(s:t)=(s^3:s^2t:st^2:t^3)$. Then $\rho_3(1:t)=(1:t:t^2:t^3)$, and $\rho_3(0:1)=(0:0:0:1)$. Take the coordinates of $\mathbf P^3$ to be $(x_0:x_1:x_2:x_3)=(y_0:\dots:y_3)$ as in 2.9(b), where $\mathbf A^3=U_0$ has coordinates $x=x_1/x_0,\dots$ and the twisted cubic of Ex.~1.2 is $Y=\{(1:t:t^2:t^3)\}$. Thus $Y=\rho_3(\{s\ne0\})$ and $\overline Y\subseteq\rho_3(\mathbf P^1)$ because $\rho_3(\mathbf P^1)$ is closed (by (c)) and contains $Y$. Conversely $\rho_3^{-1}(\overline Y)$ is a closed subset of $\mathbf P^1$ containing the infinite set $\{(1:t)\}$. A proper closed subset of $\mathbf P^1$ is finite (if $f\neq0$ is homogeneous in $x_0,x_1$, $f(1,t)$ is a nonzero polynomial in $t$ -- otherwise $f=x_0^{\deg f}f(1,x_1/x_0)=0$ -- so $Z(f)$ is finite), so $\rho_3^{-1}(\overline Y)=\mathbf P^1$ and $\overline Y\supseteq\rho_3(\mathbf P^1)$. Hence $\overline Y=\rho_3(\mathbf P^1)$, the 3-uple embedding of $\mathbf P^1$.\hfill$\square$

\exercise{2.13}{Let $Y$ be the image of the 2-uple embedding of $\mathbf P^2$ in $\mathbf P^5$. This is the \emph{Veronese surface}. If $Z\subseteq Y$ is a closed curve (a \emph{curve} is a variety of dimension 1), show that there exists a hypersurface $V\subseteq\mathbf P^5$ such that $V\cap Y=Z$.}
\pf
Let $\rho=\rho_2:\mathbf P^2\to Y$, a homeomorphism by 2.12(c). Then $W=\rho^{-1}(Z)$ is a closed subset of $\mathbf P^2$, homeomorphic to $Z$, hence irreducible (so a projective variety) of dimension $1=2-1$. By Ex.~2.8, $W=Z(f)$ for an irreducible homogeneous $f\in k[x_0,x_1,x_2]$ of degree $e>0$.

Now $f^2$ is homogeneous of even degree $2e$, and every monomial of degree $2e$ is a product of $e$ monomials of degree $2$; so $f^2=\theta(g)$ for some homogeneous polynomial $g\in k[y_0,\dots,y_5]$ of degree $e>0$ ($g\ne0$ since $f^2\ne0$), where $\theta$ is as in 2.12. Let $V=Z(g)\subseteq\mathbf P^5$, a hypersurface (the zero set of one homogeneous polynomial; $g$ need not be irreducible). For $P\in\mathbf P^2$, $g(\rho(P))=\theta(g)(P)=f(P)^2$, so $\rho(P)\in V$ iff $P\in Z(f)=W$. As $\rho:\mathbf P^2\to Y$ is a bijection,
\[
V\cap Y=\rho(\rho^{-1}(V))=\rho(W)=Z .
\]\hfill$\square$

\exercise{2.14}{\emph{The Segre Embedding.} Let $\psi:\mathbf P^r\times\mathbf P^s\to\mathbf P^N$ be the map defined by sending the ordered pair $(a_0,\dots,a_r)\times(b_0,\dots,b_s)$ to $(\dots,a_ib_j,\dots)$ in lexicographic order, where $N=rs+r+s$. Note that $\psi$ is well-defined and injective. It is called the \emph{Segre embedding}. Show that the image of $\psi$ is a subvariety of $\mathbf P^N$. [Hint: Let $\mathbf P^N$ have homogeneous coordinates $\{z_{ij}\mid i=0,\dots,r,\ j=0,\dots,s\}$, and let $\mathfrak a$ be the kernel of the homomorphism $k[\{z_{ij}\}]\to k[x_0,\dots,x_r,y_0,\dots,y_s]$ which sends $z_{ij}$ to $x_iy_j$. Then $\operatorname{Im}\psi=Z(\mathfrak a)$.]}
\pf
\emph{Well defined and injective.} Replacing $(a)$ by $\lambda a$ and $(b)$ by $\mu b$ multiplies all $a_ib_j$ by $\lambda\mu\neq0$; and not all $a_ib_j$ vanish since some $a_i\ne0$, some $b_j\ne0$. If $a_ib_j=c\,a_i'b_j'$ for all $i,j$ ($c\ne0$): choose $i_0,j_0$ with $a_{i_0}b_{j_0}\ne0$, so $a'_{i_0}b'_{j_0}\ne0$. Taking $j=j_0$: $a_i=(c\,b'_{j_0}/b_{j_0})\,a'_i$ for all $i$, so $(a)=(a')$ in $\mathbf P^r$; symmetrically $(b)=(b')$.

\emph{Image.} Let $\theta:k[z_{ij}]\to k[x,y]$, $z_{ij}\mapsto x_iy_j$, $\mathfrak a=\ker\theta$. As in 2.12(a), $\mathfrak a$ is prime (image in a domain) and homogeneous (a form of degree $e$ in the $z$ goes to a form of degree $2e$). If $g\in\mathfrak a$ is homogeneous, $g(\psi(a,b))=\theta(g)(a,b)=0$; so $\operatorname{Im}\psi\subseteq Z(\mathfrak a)$. Conversely let $Q=(c_{ij})\in Z(\mathfrak a)$ and choose $(i_0,j_0)$ with $c_{i_0j_0}\ne0$; scale so $c_{i_0j_0}=1$. For all $i,j$ the binomial $z_{ij_0}z_{i_0j}-z_{ij}z_{i_0j_0}$ lies in $\mathfrak a$ (both terms map to $x_ix_{i_0}y_jy_{j_0}$). So $c_{ij_0}c_{i_0j}=c_{ij}$. Put $a_i=c_{ij_0}$, $b_j=c_{i_0j}$; then $a_{i_0}=b_{j_0}=1$, so $(a)\in\mathbf P^r$, $(b)\in\mathbf P^s$ and $\psi((a),(b))=(a_ib_j)=(c_{ij})=Q$. Hence $\operatorname{Im}\psi=Z(\mathfrak a)$, a nonempty closed set with $I(Z(\mathfrak a))=\sqrt{\mathfrak a}=\mathfrak a$ prime (2.3(d)); by 2.4(b) it is irreducible, i.e.\ a (projective) subvariety of $\mathbf P^N$.\hfill$\square$

\exercise{2.15}{\emph{The Quadric Surface in $\mathbf P^3$} (Fig.~2). Let $Q$ be the surface (we call a surface a variety of dimension 2) in $\mathbf P^3$ defined by the equation $xy-zw=0$.
\begin{enumerate}[label=(\alph*),leftmargin=*]
\item Show that $Q$ is equal to the Segre embedding of $\mathbf P^1\times\mathbf P^1$ in $\mathbf P^3$, for suitable choice of coordinates.
\item Show that $Q$ contains two families of lines (a \emph{line} is a linear variety of dimension 1) $\{L_t\},\{M_t\}$, each parametrized by $t\in\mathbf P^1$, with the properties that if $L_t\ne L_u$, then $L_t\cap L_u=\emptyset$; if $M_t\ne M_u$, then $M_t\cap M_u=\emptyset$; and for all $t,u$, $L_t\cap M_u=$ one point.
\item Show that $Q$ contains other curves besides these lines, and deduce that the Zariski topology on $Q$ is not homeomorphic via $\psi$ to the product topology on $\mathbf P^1\times\mathbf P^1$ (where each $\mathbf P^1$ has its Zariski topology).
\end{enumerate}}
\pf
(a) Take $r=s=1$, $N=3$ in 2.14, with $\psi((a_0,a_1),(b_0,b_1))=(a_0b_0:a_0b_1:a_1b_0:a_1b_1)=(z_{00}:z_{01}:z_{10}:z_{11})$. Choose coordinates $(x,y,z,w)=(z_{00},z_{11},z_{01},z_{10})$; then $xy-zw=z_{00}z_{11}-z_{01}z_{10}$. Clearly $\operatorname{Im}\psi\subseteq Q$ (it satisfies $a_0b_0a_1b_1=a_0b_1a_1b_0$). Conversely, the proof of 2.14 shows that every point $(c_{ij})$ at which the binomials $z_{ij_0}z_{i_0j}-z_{ij}z_{i_0j_0}$ ($i,j,i_0,j_0\in\{0,1\}$) vanish lies in $\operatorname{Im}\psi$; for $r=s=1$ each such binomial is either identically $0$ or $\pm(z_{00}z_{11}-z_{01}z_{10})$, which vanishes on $Q$. So $Q=\operatorname{Im}\psi\cong$ the Segre embedding of $\mathbf P^1\times\mathbf P^1$.

(b) For $t=(t_0:t_1)\in\mathbf P^1$ put $L_t=\psi(\{t\}\times\mathbf P^1)$, $M_t=\psi(\mathbf P^1\times\{t\})$. In $z$-coordinates,
\[
\psi(t,(s_0:s_1))=s_0(t_0,0,t_1,0)+s_1(0,t_0,0,t_1),\qquad
\psi((s_0:s_1),t)=s_0(t_0,t_1,0,0)+s_1(0,0,t_0,t_1),
\]
and the two vectors in each sum are linearly independent. So $L_t=\mathbf P(W)$ with $W$ the plane (2-dimensional subspace of $k^4$) spanned by the first two vectors, and likewise $M_t$; by the Lemma in 2.11 these are lines (linear varieties of dimension $1$), contained in $Q$. Since $\psi$ is injective, $\psi(A)\cap\psi(B)=\psi(A\cap B)$. Therefore for $t\ne u$: $L_t\cap L_u=\psi((\{t\}\cap\{u\})\times\mathbf P^1)=\emptyset$, similarly $M_t\cap M_u=\emptyset$; and $L_t\cap M_u=\psi(\{t\}\times\mathbf P^1\cap\mathbf P^1\times\{u\})=\{\psi(t,u)\}$, a single point, for all $t,u$.

(c) Let $\Delta\subseteq\mathbf P^1\times\mathbf P^1$ be the diagonal and $C=\psi(\Delta)=\{(a_0^2:a_0a_1:a_0a_1:a_1^2)\}$ (in $z_{00},z_{01},z_{10},z_{11}$). Claim: $C=Q\cap H$ with $H=Z(z_{01}-z_{10})$. Indeed $\psi(a,b)\in H$ iff $a_0b_1=a_1b_0$ iff $(a)=(b)$. So $C$ is closed in $Q$ (and in $\mathbf P^3$). It is irreducible and one-dimensional: in the plane $H$ with coordinates $(u_0:u_1:u_2)=(z_{00}:z_{01}:z_{11})$ (a linear change of coordinates on $\mathbf P^3$) $C$ is the image of $\mathbf P^1$ under the 2-uple embedding $(a_0:a_1)\mapsto(a_0^2:a_0a_1:a_1^2)$, the conic $u_0u_2=u_1^2$, which by 2.12(c) is homeomorphic to $\mathbf P^1$, so is a curve ($\dim\mathbf P^1=1$ by 2.7(a)). $C$ is not a line: it contains $(1:0:0:0),(0:0:0:1),(1:1:1:1)$, three points whose vectors in $k^4$ are linearly independent, so no 2-dimensional subspace contains them. Hence $C$ is a curve in $Q$ different from all $L_t$, $M_u$.

Now suppose the Zariski topology on $Q$ corresponded to the product topology under $\psi$. Since $C$ is closed in $Q$, $\psi^{-1}(C)=\Delta$ would be closed in the product topology on $\mathbf P^1\times\mathbf P^1$. But in the product topology every proper closed set $F$ is contained in a finite union of ``fibers'' $\{p\}\times\mathbf P^1$, $\mathbf P^1\times\{q\}$: the complement of $F$ contains a basic open set $U\times V$ with $U,V$ nonempty open, hence cofinite (proper closed subsets of $\mathbf P^1$ are finite), and then $F\subseteq((\mathbf P^1\setminus U)\times\mathbf P^1)\cup(\mathbf P^1\times(\mathbf P^1\setminus V))$. But each fiber meets $\Delta$ in at most one point, so a finite union of fibers meets the infinite set $\Delta$ in a finite set, and $\Delta\ne\mathbf P^1\times\mathbf P^1$ is not contained in it. So $\Delta$ is not closed in the product topology, a contradiction. Thus $\psi$ is not a homeomorphism for the product topology (it is of course a bijection onto $Q$).\hfill$\square$

\exercise{2.16}{
\begin{enumerate}[label=(\alph*),leftmargin=*]
\item The intersection of two varieties need not be a variety. For example, let $Q_1$ and $Q_2$ be the quadric surfaces in $\mathbf P^3$ given by the equations $x^2-yw=0$ and $xy-zw=0$, respectively. Show that $Q_1\cap Q_2$ is the union of a twisted cubic curve and a line.
\item Even if the intersection of two varieties is a variety, the ideal of the intersection may not be the sum of the ideals. For example, let $C$ be the conic in $\mathbf P^2$ given by the equation $x^2-yz=0$. Let $L$ be the line given by $y=0$. Show that $C\cap L$ consists of one point $P$, but that $I(C)+I(L)\ne I(P)$.
\end{enumerate}}
\pf
(a) Let $(x:y:z:w)$ be the coordinates. $Q_1=Z(x^2-yw)$ is a variety: $x^2-yw$ is of degree $1$ in $y$ with coprime coefficients $-w,x^2$, hence irreducible, so 2.8 applies. $Q_2$ is a variety by 2.15(a). I show $Q_1\cap Q_2=C\cup L$ with
\[
C=\{(ab^2:a^2b:a^3:b^3)\},\qquad L=Z(x,w)=\{(0:y:z:0)\}.
\]
Let $P=(x:y:z:w)\in Q_1\cap Q_2$, so $x^2=yw$ and $xy=zw$. If $w\ne0$, scale $w=1$: $y=x^2$, $z=xy=x^3$, so $P=(x:x^2:x^3:1)=(ab^2:a^2b:a^3:b^3)$ with $(a:b)=(x:1)$, so $P\in C$. If $w=0$ then $x^2=0$, $x=0$ and $P\in L$ (the second equation then holds automatically). Conversely $L\subseteq Q_1\cap Q_2$ trivially, and $C\subseteq Q_1\cap Q_2$ because $(ab^2)^2=a^2b^4=(a^2b)(b^3)$ and $(ab^2)(a^2b)=a^3b^3=(a^3)(b^3)$.

$L$ is a line (linear variety of dimension $1$, by the Lemma in 2.11). $C$ is the image of the 3-uple embedding $\rho_3(a:b)=(a^3:a^2b:ab^2:b^3)$ of $\mathbf P^1$ after the permutation of coordinates $(x,y,z,w)=(X_2,X_1,X_0,X_3)$; so by 2.12(c),(d) it is a twisted cubic curve (closed, irreducible, homeomorphic to $\mathbf P^1$). $C\ne\emptyset\ne L$, $C\not\subseteq L$ (e.g.\ $(1:1:1:1)\in C\setminus L$), $L\not\subseteq C$ ($C\cap L=\{(0:0:1:0)\}$, the point $b=0$, while $L$ is infinite). So $Q_1\cap Q_2$ is the union of a twisted cubic and a line, neither containing the other; hence it is reducible, not a variety.

(b) $f=x^2-yz$ is irreducible (degree $1$ in $y$ with coprime coefficients $-z,x^2$); $Z(f)\ne\emptyset$ (it contains $(0:1:0)$). So $(f)$ is prime and $I(C)=\sqrt{(f)}=(f)$ by 2.3(d). Likewise $I(L)=(y)$. A point of $C\cap L$ has $y=0$, so $x^2=0$, $x=0$; so $C\cap L=\{P\}$ with $P=(0:0:1)$. Now
\[
I(C)+I(L)=(x^2-yz,y)=(x^2,y),
\]
while $x\in I(P)$ ($x$ vanishes at $P$). But $x\notin(x^2,y)$: the degree-$1$ part of $(x^2,y)$ is $ky$. So $I(C)+I(L)\ne I(P)$. (In fact $I(P)=(x,y)$.)\hfill$\square$

\exercise{2.17}{\emph{Complete intersections.} A variety $Y$ of dimension $r$ in $\mathbf P^n$ is a (strict) \emph{complete intersection} if $I(Y)$ can be generated by $n-r$ elements. $Y$ is a \emph{set-theoretic complete intersection} if $Y$ can be written as the intersection of $n-r$ hypersurfaces.
\begin{enumerate}[label=(\alph*),leftmargin=*]
\item Let $Y$ be a variety in $\mathbf P^n$, let $Y=Z(\mathfrak a)$; and suppose that $\mathfrak a$ can be generated by $q$ elements. Then show that $\dim Y\ge n-q$.
\item Show that a strict complete intersection is a set-theoretic complete intersection.
\item The converse of (b) is false. For example let $Y$ be the twisted cubic curve in $\mathbf P^3$ (Ex.~2.9). Show that $I(Y)$ cannot be generated by two elements. On the other hand, find hypersurfaces $H_1,H_2$ of degrees $2,3$ respectively, such that $Y=H_1\cap H_2$.
\item It is an unsolved problem whether every closed irreducible curve in $\mathbf P^3$ is a set-theoretic intersection of two surfaces. See Hartshorne [1] and Hartshorne [5, Ch.~III, Section 5] for commentary.
\end{enumerate}
(Parts (c) and (d) are starred in the book, (d) with two stars; (d) is a remark on an open problem and calls for no proof.)}
\pf
(a) Let $\mathfrak a$ be a homogeneous ideal with $Z(\mathfrak a)=Y$ (e.g.\ $\mathfrak a=I(Y)$), generated by $g_1,\dots,g_q$ (not necessarily homogeneous). As in 2.10(a), $\widetilde Z(\mathfrak a)=C(Y)$: for $a\neq 0$, $a\in\widetilde Z(\mathfrak a)$ iff $\theta(a)\in Z(\mathfrak a)=Y$, and $0\in\widetilde Z(\mathfrak a)$ because $\mathfrak a\subseteq S_+$ (as $Z(\mathfrak a)\ne\emptyset$, $\mathfrak a$ contains no nonzero constant). Also $\widetilde Z(\mathfrak a)=\widetilde Z(g_1,\dots,g_q)$. By Ex.~1.9 (applied in $\mathbf A^{n+1}$) every irreducible component of $C(Y)$ has dimension $\ge(n+1)-q$. But $C(Y)$ is irreducible by 2.10(b), so $\dim C(Y)\ge n+1-q$, and by 2.10(c) $\dim Y=\dim C(Y)-1\ge n-q$.

(b) \emph{Lemma (graded generators).} Let $I\subseteq S_+$ be a homogeneous ideal generated by $m$ elements. Then $I$ is generated by $\mu\le m$ homogeneous elements. \emph{Proof.} As in 2.11(b), the images of $m$ generators span $I/S_+I$, so $\mu:=\dim_kI/S_+I\le m$. The vector space $I/S_+I$ is graded and spanned by classes of homogeneous elements, so choose homogeneous $h_1,\dots,h_\mu\in I$ whose classes form a basis. Let $M=I/(h_1,\dots,h_\mu)$, a graded module with $S_+M=M$ (since $I=\sum kh_j+S_+I$). If $M\ne0$, take the lowest degree $e$ with $M_e\ne0$ (degrees are bounded below, $I\subseteq S_+$); then $M_e=(S_+M)_e=\sum_{j\ge1}S_jM_{e-j}=0$, a contradiction. So $I=(h_1,\dots,h_\mu)$.\hfill$\square$

Now let $Y$ be a strict complete intersection of dimension $r$: $I(Y)$ is generated by $n-r$ elements. By the Lemma ($I(Y)\subseteq S_+$) it is generated by $\mu\le n-r$ homogeneous elements $h_1,\dots,h_\mu$, each of positive degree. By (a) with $q=\mu$, $r=\dim Y\ge n-\mu$, so $\mu\ge n-r$, hence $\mu=n-r$. Then $Y=Z(I(Y))=Z(h_1,\dots,h_{n-r})=\bigcap_jZ(h_j)$ (2.3(e)), an intersection of $n-r$ hypersurfaces (zero sets of single homogeneous polynomials, as in the problem). (If $r=n$, then $Y=\mathbf P^n$ is the empty intersection.)

(c) Let $Y=\overline Y$ be the twisted cubic of 2.9, with homogeneous coordinates $x_0,\dots,x_3$, $Y=\{(1:t:t^2:t^3)\}\cup\{(0:0:0:1)\}$ (2.12(d)).

\emph{$I(Y)$ is not generated by two elements.} By 2.9(b), $I:=I(Y)=(q_1,q_2,q_3)$ with $q_1=x_0x_2-x_1^2$, $q_2=x_0x_3-x_1x_2$, $q_3=x_1x_3-x_2^2$. The elements $q_1,q_2,q_3$ are linearly independent over $k$ ($q_1$ contains the monomial $x_0x_2$, $q_2$ contains $x_0x_3$, $q_3$ contains $x_1x_3$, each missing from the other two). $I$ has no elements of degree $\le1$ (as $I$ is generated in degree $2$), so $S_+I$ lives in degrees $\ge3$, and $I/S_+I$ contains the classes of $q_1,q_2,q_3$ in degree $2$, linearly independent. So $\dim_kI/S_+I\ge3$. If $I$ were generated by two elements, their classes would span $I/S_+I$ (as in 2.11(b)), giving $\dim\le2$. Contradiction.

\emph{$Y=H_1\cap H_2$.} Let $H_1=Z(q_1)$ (a quadric) and $H_2=Z(F)$ with the cubic $F=x_0x_3^2-2x_1x_2x_3+x_2^3$. Both $q_1$ and $F$ are irreducible (each is of degree $1$ in $x_0$ with coprime coefficients: $-x_1^2$ and $x_2$ for $q_1$, and $x_3^2$ and $x_2(x_2^2-2x_1x_3)$ for $F$), so $H_1,H_2$ are hypersurfaces in the strict sense of 2.8. Both vanish on $Y$: at $(1,t,t^2,t^3)$, $q_1=t^2-t^2=0$ and $F=t^6-2t^6+t^6=0$; at $(0,0,0,1)$ clearly $q_1=F=0$. Conversely let $P=(x_0:x_1:x_2:x_3)\in H_1\cap H_2$.
\begin{itemize}
\item If $x_0\ne0$, scale $x_0=1$. Then $q_1=0$ gives $x_2=x_1^2$, and $F=x_3^2-2x_1^3x_3+x_1^6=(x_3-x_1^3)^2=0$ gives $x_3=x_1^3$. So $P=(1:t:t^2:t^3)$ with $t=x_1$, which lies in $Y$.
\item If $x_0=0$, then $q_1=0$ gives $x_1^2=0$, $x_1=0$, and $F=0$ gives $x_2^3=0$, $x_2=0$. So $P=(0:0:0:1)\in Y$.
\end{itemize}
Hence $Y=H_1\cap H_2$ with $\deg H_1=2$, $\deg H_2=3$. (So $Y$ is a set-theoretic, but not a strict, complete intersection.)

(d) This is a remark about an open problem; there is nothing to prove.\hfill$\square$

\section*{3. Morphisms}

\exercise{3.1}{
\begin{enumerate}[label=(\alph*),leftmargin=*]
\item Show that any conic in $\mathbf A^2$ is isomorphic either to $\mathbf A^1$ or to $\mathbf A^1-\{0\}$ (cf.\ Ex.~1.1).
\item Show that $\mathbf A^1$ is \emph{not} isomorphic to any proper open subset of itself. (This result is generalized by Ex.~6.7 below.)
\item Any conic in $\mathbf P^2$ is isomorphic to $\mathbf P^1$.
\item We will see later (Ex.~4.8) that any two curves are homeomorphic. But show now that $\mathbf A^2$ is not even homeomorphic to $\mathbf P^2$.
\item If an affine variety is isomorphic to a projective variety, then it consists of only one point.
\end{enumerate}}
\pf
\textbf{Conventions.} ``Conic'' means the zero set of an irreducible polynomial (resp.\ irreducible homogeneous polynomial) of degree~$2$; a reducible quadric is a union of two lines, which is not a variety, so it must be excluded. We use freely: $k$ is algebraically closed; regularity of a function is a local property, hence so is being a morphism (continuity and regularity of pull-backs can be checked on an open cover of the source). We also use two elementary facts.

\emph{Lemma A.} Let $X\subseteq\mathbf P^m$ be quasi-projective and let $F_0,\dots,F_n\in k[x_0,\dots,x_m]$ be homogeneous of the same degree with no common zero on $X$. Then $P\mapsto(F_0(P):\dots:F_n(P))$ is a morphism $X\to\mathbf P^n$. (Similarly for quasi-affine $X$ and polynomial maps into $\mathbf A^n$; for an affine target this is Lemma~3.6.)

\emph{Proof.} Call the map $\varphi$. A closed set of $\mathbf P^n$ is an intersection of sets $Z(g)$ with $g$ homogeneous, and $\varphi^{-1}Z(g)=X\cap Z(g(F_0,\dots,F_n))$ is closed, so $\varphi$ is continuous. If $f$ is regular on an open $V\subseteq\mathbf P^n$ and $P\in\varphi^{-1}(V)$, write $f=a/b$ near $\varphi(P)$ on an open $U\ni\varphi(P)$, with $a,b$ homogeneous of the same degree and $b\ne0$ on $U$. Then on $\varphi^{-1}(U)$ we have $f\circ\varphi=a(F)/b(F)$, a quotient of homogeneous polynomials of the same degree with nowhere vanishing denominator. So $f\circ\varphi$ is regular. $\square$

\emph{Lemma B} (Ex.~3.10 below). If $\varphi:X\to Y$ is a morphism and $X',Y'$ are irreducible locally closed subsets with $\varphi(X')\subseteq Y'$, then $\varphi|_{X'}:X'\to Y'$ is a morphism. In particular: a morphism into $Y$ whose image lies in a subvariety $Y'$ is a morphism into $Y'$; and an isomorphism restricts to isomorphisms between corresponding subvarieties. (This is proved in 3.10 using only the definitions; there is no circularity.) Also, the restriction of a linear change of coordinates is an isomorphism, because projective (affine) linear automorphisms and their inverses are given by invertible matrices and Lemma~A applies.

\textbf{(a)} Let $C=Z(f)\subset\mathbf A^2$ with $f$ irreducible of degree $2$. Affine-linear changes of coordinates of $\mathbf A^2$ are automorphisms of $\mathbf A^2$ preserving the degree and irreducibility of $f$, so we may change coordinates freely. Let $q$ be the quadratic part of $f$. Since $k$ is algebraically closed, a nonzero binary quadratic form is a product of two linear forms (in any characteristic), $q=\ell_1\ell_2$.

\emph{Case 1: $\ell_1,\ell_2$ proportional.} Choose coordinates $(x,y)$ with $y=\ell_1$. Then $f=cy^2+ax+by+e$ with $c\neq0$. If $a=0$, then $f$ is a quadratic polynomial in $y$ alone, hence a product of two linear factors, contradicting irreducibility. So $a\ne0$ and $C=\{(g(y),y)\}$ where $g(y)=-(cy^2+by+e)/a$. The maps $\mathbf A^1\to C$, $t\mapsto(g(t),t)$, and $C\to\mathbf A^1$, $(x,y)\mapsto y$, are morphisms (Lemma~3.6, coordinates are polynomials/regular), and they are mutually inverse. So $C\cong\mathbf A^1$.

\emph{Case 2: $\ell_1,\ell_2$ independent.} Take $u=\ell_1,v=\ell_2$ as coordinates. Then $f=c\,uv+au+bv+e$ with $c\neq0$, so $f/c=(u+b')(v+a')+e''$ for suitable constants $a',b',e''$. Put $u'=u+b'$, $v'=v+a'$. If $e''=0$, then $f$ is reducible; so $e''\ne0$ and, after replacing $v'$ by $-v'/e''$, $C=Z(u'v'-1)$. The map $\mathbf A^1-\{0\}\to C$, $t\mapsto(t,1/t)$ is a morphism by Lemma~3.6 ($t$ and $1/t$ are regular on $\mathbf A^1-\{0\}$), and $C\to\mathbf A^1-\{0\}$, $(u',v')\mapsto u'$ (nonzero on $C$) is a morphism; they are mutually inverse. (The second map is a morphism into the open subvariety $\mathbf A^1-\{0\}$ by Lemma~B.) So $C\cong\mathbf A^1-\{0\}$.

\textbf{(b)} Let $U\subsetneq\mathbf A^1$ be open. If $U=\emptyset$ there is no isomorphism, since $\mathbf A^1\ne\emptyset$. Otherwise $U=\mathbf A^1-\{a_1,\dots,a_r\}$ with $r\ge1$ (the closed proper subsets of $\mathbf A^1$ are finite). Put $a=a_1$ and $g=(t-a)|_U$. Then $g$ is regular on $U$ and nowhere zero, so $1/g$ is regular on $U$: $g$ is a unit of $\mathcal O(U)$. Suppose $\varphi:\mathbf A^1\to U$ is an isomorphism. Then $\varphi^*:\mathcal O(U)\to\mathcal O(\mathbf A^1)$ (pull-back of functions) is a ring isomorphism (its inverse is the pull-back by $\varphi^{-1}$). By Thm.~3.2(a), $\mathcal O(\mathbf A^1)=k[t]$, whose units are the nonzero constants. So $g\circ\varphi=c$ is a constant, and since $\varphi$ is surjective, $g\equiv c$ on $U$, i.e.\ $t-a=c$ on the infinite set $U$, which is absurd. Hence $\mathbf A^1\not\cong U$.

\textbf{(c)} Let $C=Z(f)\subset\mathbf P^2$, $f$ irreducible homogeneous of degree $2$ in $x_0,x_1,x_2$. $C$ has a point: the restriction of $f$ to a line is a binary quadratic form (or zero), which has a zero since $k$ is algebraically closed. After a projective linear change of coordinates we may assume $P=(0:0:1)\in C$. As $f(P)=0$, there is no $x_2^2$ term: $f=x_2\ell(x_0,x_1)+q(x_0,x_1)$ with $\ell$ linear and $q$ quadratic. If $\ell=0$, $f=q$ is a product of linear forms: reducible. So $\ell\ne0$; change coordinates in $x_0,x_1$ so that $\ell=x_0$. Write $q=ax_0^2+bx_0x_1+cx_1^2$. If $c=0$ then $f=x_0(x_2+ax_0+bx_1)$ is reducible; so $c\ne0$. Replace $x_2$ by $x_2+ax_0+bx_1$ and $x_1$ by $\sqrt c\,x_1$, and then $x_2$ by $-x_2$: then $f=-(x_0x_2-x_1^2)$. So we may assume $C=Z(x_0x_2-x_1^2)$.

Define $\alpha:\mathbf P^1\to C$ by $(s:t)\mapsto(s^2:st:t^2)$ (a morphism to $\mathbf P^2$ by Lemma~A, with image in $C$, hence a morphism to $C$ by Lemma~B). Define $\beta:C\to\mathbf P^1$ on $C\cap\{x_0\ne0\}$ by $(x_0:x_1)$ and on $C\cap\{x_2\ne0\}$ by $(x_1:x_2)$. These are morphisms (Lemma~A), they agree on the overlap since $x_0x_2=x_1^2$ gives $(x_0:x_1)=(x_1:x_2)$ whenever $x_0x_2\neq0$, and the two open sets cover $C$ (if $x_0=x_2=0$ then $x_1=0$). So $\beta$ is a morphism. Now $\beta\alpha(s:t)=(s^2:st)=(s:t)$ if $s\neq0$ and $=(st:t^2)=(s:t)$ if $t\ne0$. And for $(x_0:x_1:x_2)\in C$ with $x_0\neq0$: $\alpha\beta=(x_0^2:x_0x_1:x_1^2)=(x_0^2:x_0x_1:x_0x_2)=(x_0:x_1:x_2)$; with $x_2\neq0$: $\alpha\beta=(x_1^2:x_1x_2:x_2^2)=(x_0x_2:x_1x_2:x_2^2)=(x_0:x_1:x_2)$. Thus $C\cong\mathbf P^1$.

\textbf{(d)} Dimension is a topological notion (supremum of lengths of chains of irreducible closed subsets), and irreducibility and closedness are preserved by homeomorphisms. In $\mathbf A^2$ the sets $Z(x)$ and $Z(x-1)$ are disjoint irreducible closed subsets, each isomorphic to $\mathbf A^1$, hence of dimension $1$ (curves). A homeomorphism $h:\mathbf A^2\to\mathbf P^2$ would carry them to two disjoint irreducible closed subsets of $\mathbf P^2$ of dimension $1$, i.e.\ two disjoint curves in $\mathbf P^2$. This contradicts Ex.~3.7(a) below (any two curves in $\mathbf P^2$ meet), whose proof uses only 3.5 and 3.1(e), not (d).

\textbf{(e)} Let $X\subseteq\mathbf A^n$ be an affine variety and $\theta:X\to Y$ an isomorphism onto a projective variety. Then $\theta^*$ gives a ring isomorphism $\mathcal O(Y)\to\mathcal O(X)$. By Thm.~3.4(a), $\mathcal O(Y)=k$, and by Thm.~3.2(a), $\mathcal O(X)=A(X)=k[x_1,\dots,x_n]/I(X)$. So the natural map $k\to A(X)$ is onto: each $x_i$ is congruent mod $I(X)$ to a constant $c_i$. Then $x_i-c_i\in I(X)$ for all $i$, so $X\subseteq\{(c_1,\dots,c_n)\}$, and $X\ne\emptyset$ gives that $X$ is a single point. \hfill$\square$


\exercise{3.2}{A morphism whose underlying map on the topological spaces is a homeomorphism need not be an isomorphism.
\begin{enumerate}[label=(\alph*),leftmargin=*]
\item For example, let $\varphi:\mathbf A^1\to\mathbf A^2$ be defined by $t\mapsto(t^2,t^3)$. Show that $\varphi$ defines a bijective bicontinuous morphism of $\mathbf A^1$ onto the curve $y^2=x^3$, but that $\varphi$ is not an isomorphism.
\item For another example, let the characteristic of the base field $k$ be $p>0$, and define a map $\varphi:\mathbf A^1\to\mathbf A^1$ by $t\mapsto t^p$. Show that $\varphi$ is bijective and bicontinuous but not an isomorphism. This is called the \emph{Frobenius morphism}.
\end{enumerate}}
\pf
\textbf{(a)} Put $Y=Z(y^2-x^3)$. Clearly $\varphi(t)=(t^2,t^3)\in Y$.

\emph{$\varphi:\mathbf A^1\to Y$ is bijective.} Injective: if $\varphi(t)=\varphi(s)$ with $t\ne0$, then $t=t^3/t^2=s^3/s^2=s$ (note $s^2=t^2\neq0$); if $t=0$ then $s^2=0$ so $s=0$. Surjective: let $(x,y)\in Y$. If $x=0$ then $y^2=0$, and $(0,0)=\varphi(0)$. If $x\ne0$ put $t=y/x$; then $t^2=y^2/x^2=x^3/x^2=x$ and $t^3=tx=y$.

\emph{$Y$ is an affine variety and $\varphi$ is a morphism.} $Y=\varphi(\mathbf A^1)$ is the continuous image of an irreducible space, hence irreducible, and it is closed by definition, so $Y$ is an affine variety. The coordinates $t^2,t^3$ of $\varphi$ are regular, so $\varphi$ is a morphism by Lemma~3.6.

\emph{Bicontinuity.} $\varphi$ is continuous. It is also closed: the closed subsets of $\mathbf A^1$ are $\mathbf A^1$ (with image $Y$, closed) and finite sets (with finite images, which are closed). A closed continuous bijection is a homeomorphism.

\emph{Not an isomorphism.} Otherwise $\varphi^*:\mathcal O(Y)\to\mathcal O(\mathbf A^1)$ would be an isomorphism. By Thm.~3.2(a), $\mathcal O(Y)=A(Y)$, which is generated as a $k$-algebra by the coordinate functions $x,y$, and $\mathcal O(\mathbf A^1)=k[t]$. But $\varphi^*(x)=t^2$, $\varphi^*(y)=t^3$, so the image of $\varphi^*$ is $k[t^2,t^3]$, the span of monomials $t^{2a+3b}$; the exponent $1$ does not occur, so $t\notin k[t^2,t^3]$ and $\varphi^*$ is not surjective.

\textbf{(b)} $\varphi$ is a morphism by Lemma~3.6. \emph{Injective:} $t^p=s^p$ implies $(t-s)^p=t^p-s^p=0$ (characteristic $p$), so $t=s$. \emph{Surjective:} $k$ is algebraically closed, so $T^p-a$ has a root. As in (a), $\varphi$ is closed (finite sets go to finite sets, $\mathbf A^1$ to $\mathbf A^1$), hence a homeomorphism. If $\varphi$ were an isomorphism, $\varphi^*:k[t]\to k[t]$, $f(t)\mapsto f(t^p)$, would be onto. Its image is $k[t^p]\not\ni t$ as $p>1$. So $\varphi$ is not an isomorphism. \hfill$\square$


\exercise{3.3}{
\begin{enumerate}[label=(\alph*),leftmargin=*]
\item Let $\varphi:X\to Y$ be a morphism. Then for each $P\in X$, $\varphi$ induces a homomorphism of local rings $\varphi^*_P:\mathcal O_{\varphi(P),Y}\to\mathcal O_{P,X}$.
\item Show that a morphism $\varphi$ is an isomorphism if and only if it is a homeomorphism, and the induced map $\varphi^*_P$ on local rings is an isomorphism, for all $P\in X$.
\item Show that if $\varphi(X)$ is dense in $Y$, then the map $\varphi^*_P$ is \emph{injective} for all $P\in X$.
\end{enumerate}}
\pf
\textbf{(a)} Let $Q=\varphi(P)$. Define $\varphi^*_P(\langle V,f\rangle)=\langle\varphi^{-1}(V),f\circ\varphi\rangle$. Here $V\ni Q$ is open in $Y$ and $f$ is regular on $V$; then $\varphi^{-1}(V)$ is an open neighbourhood of $P$ (continuity) and $f\circ\varphi$ is regular on it by the definition of a morphism. \emph{Well defined:} if $\langle V,f\rangle\sim\langle V',f'\rangle$ then $f=f'$ on $V\cap V'$, hence $f\circ\varphi=f'\circ\varphi$ on $\varphi^{-1}(V\cap V')=\varphi^{-1}(V)\cap\varphi^{-1}(V')$, so the two germs at $P$ agree. \emph{Ring homomorphism:} $(f+g)\circ\varphi=f\circ\varphi+g\circ\varphi$, $(fg)\circ\varphi=(f\circ\varphi)(g\circ\varphi)$, and $1\circ\varphi=1$. It is moreover a \emph{local} homomorphism: if $f(Q)=0$ then $(f\circ\varphi)(P)=0$, so $\varphi^*_P(\mathfrak m_Q)\subseteq\mathfrak m_P$. Also, if $\psi:Y\to Z$ is another morphism, $(\psi\circ\varphi)^*_P=\varphi^*_P\circ\psi^*_{\varphi(P)}$, and $(\mathrm{id}_X)^*_P=\mathrm{id}$.

\textbf{(b)} ($\Rightarrow$) If $\varphi$ is an isomorphism with inverse $\psi$, then $\varphi$ is a homeomorphism, and by functoriality, $\varphi^*_P\circ\psi^*_{\varphi(P)}=(\psi\varphi)^*_P=\mathrm{id}$, and, with $Q=\varphi(P)$, $\psi^*_{Q}\circ\varphi^*_{\psi(Q)}=(\varphi\psi)^*_Q=\mathrm{id}$, i.e.\ $\psi^*_{\varphi(P)}\circ\varphi^*_P=\mathrm{id}$. So $\varphi^*_P$ is an isomorphism with inverse $\psi^*_{\varphi(P)}$.

($\Leftarrow$) Suppose $\varphi$ is a homeomorphism and all $\varphi^*_P$ are isomorphisms. Let $\psi=\varphi^{-1}:Y\to X$, continuous. We must show: if $U\subseteq X$ is open and $f$ is regular on $U$, then $f\circ\psi$ is regular on $\psi^{-1}(U)=\varphi(U)$. Let $Q\in\varphi(U)$, $Q=\varphi(P)$, $P\in U$. Since $\varphi^*_P$ is surjective there is $\langle V,g\rangle$ with $V\ni Q$ open in $Y$ and $g$ regular on $V$, such that the germ of $g\circ\varphi$ equals the germ of $f$ at $P$: there is an open $W$ with $P\in W\subseteq U\cap\varphi^{-1}(V)$ and $f=g\circ\varphi$ on $W$. Then $\varphi(W)$ is open in $Y$ (as $\varphi$ is a homeomorphism), contains $Q$, and lies in $V$. For $y\in\varphi(W)$, $y=\varphi(w)$ with $w\in W$, and $f(\psi(y))=f(w)=g(\varphi(w))=g(y)$. So $f\circ\psi=g$ on the neighbourhood $\varphi(W)$ of $Q$, hence is regular at $Q$. Thus $\psi$ is a morphism, and it is inverse to $\varphi$.

\textbf{(c)} Suppose $\varphi(X)$ is dense in $Y$ and $\varphi^*_P\langle V,f\rangle=0$. Then there is an open $W\ni P$, $W\subseteq\varphi^{-1}(V)$ with $f\circ\varphi=0$ on $W$, i.e.\ $f$ vanishes on $\varphi(W)\subseteq V$. $W$ is a nonempty open subset of the irreducible space $X$, hence dense in $X$, so by continuity $\varphi(X)\subseteq\overline{\varphi(W)}$. Hence $\varphi(W)$ is dense in $Y$ (its closure contains the dense set $\varphi(X)$), so $\varphi(W)$ is dense in the open set $V$ as well (the closure of $\varphi(W)$ in $V$ is $\overline{\varphi(W)}\cap V=V$). A regular function is continuous (Lemma~3.1), so $\{y\in V:f(y)=0\}$ is closed in $V$; it contains the dense subset $\varphi(W)$, hence equals $V$. So $f=0$ on $V$ and the germ $\langle V,f\rangle$ is $0$. \hfill$\square$


\exercise{3.4}{Show that the $d$-uple embedding of $\mathbf P^n$ (Ex.~2.12) is an isomorphism onto its image.}
\pf
Let $M_0,\dots,M_N$ be the monomials of degree $d$ in $x_0,\dots,x_n$, with $y_I$ the coordinate on $\mathbf P^N$ corresponding to the monomial $I=x^I$. Let $v=\rho_d:\mathbf P^n\to\mathbf P^N$ and $Z=Z(\mathfrak a)$ the image, a projective variety (Ex.~2.12(a),(b)). By Ex.~2.12(c), $v:\mathbf P^n\to Z$ is a homeomorphism; in particular $v$ is a bijection onto $Z$.

\emph{$v$ is a morphism.} The $M_I$ are homogeneous of degree $d$ without common zero on $\mathbf P^n$ (at a point $x$, some $x_j\ne0$ and $M=x_j^d\ne0$). By Lemma~A (see 3.1), $v:\mathbf P^n\to\mathbf P^N$ is a morphism; it has image in the subvariety $Z$, so it is a morphism $\mathbf P^n\to Z$ (Lemma~B / Ex.~3.10).

\emph{The inverse is a morphism.} For $j=0,\dots,n$ let $Z_j=Z\cap\{y_{de_j}\ne0\}$, where $y_{de_j}$ is the coordinate corresponding to $x_j^d$. These are open in $Z$ and cover $Z$ (every point of $\mathbf P^n$ has some $x_j\ne0$, and $y_{de_j}(v(x))=x_j^d$). On $Z_j$ define $\psi_j=(y_{(d-1)e_j+e_0}:\dots:y_{(d-1)e_j+e_n})$, where $y_{(d-1)e_j+e_i}$ is the coordinate for $x_j^{d-1}x_i$. These are linear forms in the $y_I$, and for $i=j$ the form is $y_{de_j}$, which is nonzero on $Z_j$; so by Lemma~A, $\psi_j:Z_j\to\mathbf P^n$ is a morphism. If $y=v(x)\in Z_j$, then $x_j\ne0$ and $\psi_j(y)=(x_j^{d-1}x_0:\dots:x_j^{d-1}x_n)=x$. Hence $\psi_j=v^{-1}|_{Z_j}$ (here we use that $v$ is bijective onto $Z$). Since $v^{-1}$ is a morphism on each member of an open cover of $Z$, it is a morphism $Z\to\mathbf P^n$ (continuity and regularity of pull-backs are local on the source). Therefore $v$ is an isomorphism of $\mathbf P^n$ onto $Z$. \hfill$\square$


\exercise{3.5}{By abuse of language, we will say that a variety ``is affine'' if it is isomorphic to an affine variety. If $H\subseteq\mathbf P^n$ is any hypersurface, show that $\mathbf P^n-H$ is affine. [\emph{Hint}: Let $H$ have degree $d$. Then consider the $d$-uple embedding of $\mathbf P^n$ in $\mathbf P^N$ and use the fact that $\mathbf P^N$ minus a hyperplane is affine.]}
\pf
Let $H=Z(f)$ with $f$ homogeneous of degree $d\ge1$ (this is all we use; $f$ need not be irreducible). Write $f=\sum_Ic_Ix^I$, sum over monomials of degree $d$, and put $L=\sum_Ic_Iy_I$, a nonzero linear form on $\mathbf P^N$. Then $f(x)=L(v(x))$ for the $d$-uple embedding $v:\mathbf P^n\to Z\subseteq\mathbf P^N$; hence
\[
v(\mathbf P^n-H)=Z-Z(L)=Z\cap(\mathbf P^N-H_L),\qquad H_L=Z(L).
\]
By Ex.~3.4, $v$ is an isomorphism $\mathbf P^n\to Z$, so its restriction is an isomorphism of the open subvariety $\mathbf P^n-H$ onto the open subset $Z\cap(\mathbf P^N-H_L)$ of $Z$ (nonempty, since $f\ne0$ has some non-root).

Choose a projective linear automorphism $\sigma$ of $\mathbf P^N$ carrying $H_L$ to $\{y_0=0\}$ (it is an isomorphism of varieties). Then $\sigma$ maps $Z\cap(\mathbf P^N-H_L)$ isomorphically onto $Z'\cap U_0$ where $Z'=\sigma(Z)$ and $U_0=\{y_0\ne0\}$ (Ex.~3.10 / Lemma~B). By Prop.~3.3, $\varphi_0:U_0\to\mathbf A^N$ is an isomorphism. The set $Z'\cap U_0$ is closed in $U_0$ and irreducible (nonempty open in the irreducible $Z'$), so $W:=\varphi_0(Z'\cap U_0)$ is a closed irreducible subset of $\mathbf A^N$, i.e.\ an affine variety, and $\varphi_0$ restricts to an isomorphism $Z'\cap U_0\to W$ (Lemma~B applied to $\varphi_0$ and $\varphi_0^{-1}$). Composing, $\mathbf P^n-H\cong W$ is affine. \hfill$\square$


\exercise{3.6}{There are quasi-affine varieties which are not affine. For example, show that $X=\mathbf A^2-\{(0,0)\}$ is not affine. [\emph{Hint}: Show that $\mathcal O(X)\cong k[x,y]$ and use (3.5). See (III, Ex.~4.3) for another proof.]}
\pf
\emph{Step 1: $\mathcal O(X)=k[x,y]$, i.e.\ every regular function on $X$ is the restriction of a polynomial.} Let $f\in\mathcal O(X)$. We have $X=D_x\cup D_y$ where $D_x=\mathbf A^2-Z(x)$, $D_y=\mathbf A^2-Z(y)$. The set $D_x$ is isomorphic to the closed subset $W=Z(xz-1)\subseteq\mathbf A^3$ via $(x,y)\mapsto(x,y,1/x)$, with inverse $(x,y,z)\mapsto(x,y)$ (both are morphisms by Lemma~3.6 and Lemma~A). $W$ is irreducible, being homeomorphic to $D_x$, and $xz-1$ is irreducible, so $I(W)=(xz-1)$ and $\mathcal O(W)=A(W)=k[x,y,z]/(xz-1)$ (Thm.~3.2(a)). Therefore $f|_{D_x}=g/x^a$ for some $g\in k[x,y]$, $a\geq0$, and we may choose $a$ minimal. Similarly $f|_{D_y}=h/y^b$ with $h\in k[x,y]$, $b\ge0$ minimal. On $D_x\cap D_y$ (a nonempty open subset of $\mathbf A^2$, dense) we have $g/x^a=h/y^b$, so $gy^b=hx^a$ in the domain $k[x,y]$ (equality of functions on a dense set of $\mathbf A^2$ means equality of polynomials). If $a\ge1$, then $x^a\mid gy^b$ and $\gcd(x,y)=1$ gives $x^a\mid g$, so $g/x^a$ is a polynomial, contradicting minimality of $a$ unless $a=0$; hence $a=0$, and likewise $b=0$. So $f|_{D_x}=g$ and $f|_{D_y}=h$ with $g,h$ polynomials, and $g=h$ on $D_x\cap D_y$, hence $g=h$. Thus $f=g|_X$. The restriction map $i^*:k[x,y]\to\mathcal O(X)$ is injective since $X$ is dense in $\mathbf A^2$; hence $i^*:\mathcal O(\mathbf A^2)\to\mathcal O(X)$ is an isomorphism, where $i:X\hookrightarrow\mathbf A^2$ is the inclusion.

\emph{Step 2.} Suppose $X$ were affine: there is an isomorphism $\theta:Y\to X$ with $Y$ an affine variety. Put $j=i\circ\theta:Y\to\mathbf A^2$, a morphism. Then $j^*=\theta^*\circ i^*:\mathcal O(\mathbf A^2)\to\mathcal O(Y)$ is an isomorphism of $k$-algebras (product of two isomorphisms). Let $h=(j^*)^{-1}:A(Y)=\mathcal O(Y)\to A(\mathbf A^2)=\mathcal O(\mathbf A^2)$. By Prop.~3.5 (with source $\mathbf A^2$ and affine target $Y$), the bijection $\operatorname{Hom}(\mathbf A^2,Y)\to\operatorname{Hom}(A(Y),\mathcal O(\mathbf A^2))$ gives a morphism $g:\mathbf A^2\to Y$ with $g^*=h$. Then $(j\circ g)^*=g^*\circ j^*=h\circ j^*=\mathrm{id}_{\mathcal O(\mathbf A^2)}=(\mathrm{id}_{\mathbf A^2})^*$. By Prop.~3.5 again (source and target $\mathbf A^2$) morphisms $\mathbf A^2\to\mathbf A^2$ are determined by the induced maps on $A(\mathbf A^2)$, so $j\circ g=\mathrm{id}_{\mathbf A^2}$. Thus $j$ is surjective onto $\mathbf A^2$. But $j(Y)=i(\theta(Y))=X\not\ni(0,0)$, a contradiction. So $X$ is quasi-affine but not affine. \hfill$\square$


\exercise{3.7}{
\begin{enumerate}[label=(\alph*),leftmargin=*]
\item Show that any two curves in $\mathbf P^2$ have a nonempty intersection.
\item More generally, show that if $Y\subseteq\mathbf P^n$ is a projective variety of dimension $\ge1$, and if $H$ is a hypersurface, then $Y\cap H\ne\emptyset$. [\emph{Hint}: Use (Ex.~3.5) and (Ex.~3.1e).] See (7.2) for a generalization.
\end{enumerate}}
\pf
\textbf{(b)} Suppose $Y\cap H=\emptyset$, so $Y\subseteq\mathbf P^n-H$. By Ex.~3.5 there is an isomorphism $\theta:\mathbf P^n-H\to W$ onto an affine variety $W\subseteq\mathbf A^N$ (a closed irreducible set). Since $Y$ is closed in $\mathbf P^n$, it is closed in $\mathbf P^n-H$, so $\theta(Y)$ is closed in $W$, hence in $\mathbf A^N$, and irreducible: $\theta(Y)$ is an affine variety. By Lemma~B (Ex.~3.10), $\theta$ restricts to an isomorphism $Y\to\theta(Y)$ (the structure on $Y$ induced from $\mathbf P^n$ is the one it has as a subvariety of $\mathbf P^n-H$, since it is the same locally closed subset). So the projective variety $Y$ is isomorphic to an affine variety. By Ex.~3.1(e), $Y$ is a point, which has dimension $0$, contradicting $\dim Y\ge1$. Hence $Y\cap H\neq\emptyset$.

\textbf{(a)} Let $Y$ and $C$ be curves in $\mathbf P^2$, i.e.\ projective varieties of dimension $1$. By Ex.~2.8, $C=Z(g)$ is a hypersurface (a projective variety of dimension $n-1=1$ in $\mathbf P^2$ is the zero set of one irreducible homogeneous polynomial). Apply (b) with $n=2$, $H=C$: $Y\cap C\neq\emptyset$. \hfill$\square$


\exercise{3.8}{Let $H_i$ and $H_j$ be the hyperplanes in $\mathbf P^n$ defined by $x_i=0$ and $x_j=0$, with $i\ne j$. Show that any regular function on $\mathbf P^n-(H_i\cap H_j)$ is constant. (This gives an alternate proof of (3.4a) in the case $Y=\mathbf P^n$.)}
\pf
We may assume $i=0$, $j=1$, so $n\ge1$. Let $U=\mathbf P^n-(H_0\cap H_1)=U_0\cup U_1$, $U_l=\{x_l\ne0\}$, and let $f\in\mathcal O(U)$. By Prop.~3.3, $U_0\cong\mathbf A^n$ with coordinates $t_k=x_k/x_0$ ($1\le k\le n$), and $\mathcal O(\mathbf A^n)=k[t_1,\dots,t_n]$ (Thm.~3.2(a)); so $f|_{U_0}=p(t_1,\dots,t_n)$ for a polynomial $p$. Similarly $U_1\cong\mathbf A^n$ with coordinates $s_0=x_0/x_1$, $s_k=x_k/x_1$ ($2\le k\le n$), and $f|_{U_1}=q(s_0,s_2,\dots,s_n)$ for a polynomial $q$.

The set $U_0\cap U_1=\{x_0x_1\ne0\}$ is a nonempty open subset of $U_0\cong\mathbf A^n$, on which $s_0=1/t_1$, $s_k=t_k/t_1$. Equality of $f$ computed both ways gives, as functions on a dense open subset of $\mathbf A^n$, hence as elements of $k(t_1,\dots,t_n)$:
\[
p(t_1,\dots,t_n)=q\bigl(t_1^{-1},t_2t_1^{-1},\dots,t_nt_1^{-1}\bigr).
\]
Write $q=\sum c_{a}\,s_0^{a_0}s_2^{a_2}\cdots s_n^{a_n}$. The monomial $s_0^{a_0}s_2^{a_2}\cdots s_n^{a_n}$ goes to the Laurent monomial $t_1^{-(a_0+a_2+\dots+a_n)}t_2^{a_2}\cdots t_n^{a_n}$, and $(a_0,a_2,\dots,a_n)\mapsto(-(a_0+a_2+\dots+a_n),a_2,\dots,a_n)$ is injective. Distinct Laurent monomials are linearly independent in $k[t_1^{\pm1},t_2,\dots,t_n]\subseteq k(t_1,\dots,t_n)$. The left side $p$ contains only monomials with nonnegative exponent of $t_1$, while every monomial on the right has $t_1$-exponent $-(a_0+a_2+\dots+a_n)\le0$. By linear independence only monomials with all $a_l=0$ can have nonzero coefficient: $q=c_0$ is constant, hence $p=c_0$ is constant. So $f$ is constant on $U_0$ and, since $U_0$ is dense in the irreducible $U$ and regular functions that agree on a dense open set agree (Remark~3.1.1), $f\equiv c_0$ on $U$.

For $Y=\mathbf P^n$: a regular function on $\mathbf P^n$ restricts to a regular function on $U$ ($n\ge1$), hence is constant on the dense open $U$, hence constant (Remark~3.1.1). (For $n=0$ the claim is trivial.) \hfill$\square$


\exercise{3.9}{The homogeneous coordinate ring of a projective variety is not invariant under isomorphism. For example, let $X=\mathbf P^1$, and let $Y$ be the $2$-uple embedding of $\mathbf P^1$ in $\mathbf P^2$. Then $X\cong Y$ (Ex.~3.4). But show that $S(X)\not\cong S(Y)$.}
\pf
By Ex.~3.4 with $n=1$, $d=2$, $v:\mathbf P^1\to\mathbf P^2$, $(s:t)\mapsto(s^2:st:t^2)$ is an isomorphism onto its image $Y$. Using the coordinates $y_0,y_1,y_2$ on $\mathbf P^2$, we showed in Ex.~3.1(c) that the image is $Y=Z(g)$ with $g=y_0y_2-y_1^2$ (directly: a point $(a:b:c)$ with $ac=b^2$ equals $(a^2:ab:b^2)$ if $a\ne0$, and equals $(0:0:1)=v(0:1)$ if $a=0$). The polynomial $g$ is irreducible (it has degree $1$ in $y_0$ with coefficients $y_2$ and $-y_1^2$, which have no common factor), so $I(Y)=\sqrt{(g)}=(g)$, so $S(Y)=k[y_0,y_1,y_2]/(g)$, while $S(X)=k[s,t]$.

\emph{Claim: $S(Y)\cong R:=k[s^2,st,t^2]\subseteq k[s,t]$.} Consider $\theta:k[y_0,y_1,y_2]\to k[s,t]$, $y_0\mapsto s^2,y_1\mapsto st,y_2\mapsto t^2$. Its image is $R$ and $g\in\ker\theta$. Conversely, modulo $g$ every polynomial can be written as $A(y_0,y_2)+y_1B(y_0,y_2)$ (replace $y_1^2$ by $y_0y_2$). Then $\theta(A+y_1B)=A(s^2,t^2)+st\,B(s^2,t^2)$; the first summand involves only monomials $s^{2a}t^{2b}$ and the second only $s^{2a+1}t^{2b+1}$, so if this is $0$ both vanish; as $s^2,t^2$ are algebraically independent, $A=B=0$. So $\ker\theta=(g)$ and $S(Y)\cong R$.

\emph{$R$ is not a UFD.} $R$ is a domain, being a subring of $k[s,t]$. $R$ consists of the polynomials all of whose monomials have even total degree. Its units are the nonzero constants (units of $k[s,t]$). Consider $s^2,\ st,\ t^2\in R$ (the elements $y_0,y_1,y_2$). Each is irreducible in $R$: if a degree-$2$ element equals $ab$ with $a,b\in R$ nonzero, then $\deg a+\deg b=2$; a nonconstant element of $R$ has degree $\ge2$ (all its monomials have even degree and it has at least one nonconstant one), so one of $a,b$ is constant, i.e.\ a unit. Moreover $st$ is not an associate of $s^2$ or $t^2$ (associates differ by a nonzero constant). But $(st)\cdot(st)=s^2\cdot t^2$ gives two factorizations of one element into irreducibles which are not associate. So $R$ is not a UFD.

Since $S(X)=k[s,t]$ is a UFD and $S(Y)\cong R$ is not, $S(X)\not\cong S(Y)$, even as abstract rings (a fortiori as graded $k$-algebras). \hfill$\square$


\exercise{3.10}{\emph{Subvarieties.} A subset of a topological space is \emph{locally closed} if it is an open subset of its closure, or, equivalently, if it is the intersection of an open set with a closed set.

If $X$ is a quasi-affine or quasi-projective variety and $Y$ is an irreducible locally closed subset, then $Y$ is also a quasi-affine (respectively, quasi-projective) variety, by virtue of being a locally closed subset of the same affine or projective space. We call this the \emph{induced structure} on $Y$, and we call $Y$ a \emph{subvariety} of $X$.

Now let $\varphi:X\to Y$ be a morphism, let $X'\subseteq X$ and $Y'\subseteq Y$ be irreducible locally closed subsets such that $\varphi(X')\subseteq Y'$. Show that $\varphi|_{X'}:X'\to Y'$ is a morphism.}
\pf
Let $\varphi'=\varphi|_{X'}:X'\to Y'$. It is continuous, since $X',Y'$ carry the subspace topologies. Let $V'\subseteq Y'$ be open and $f:V'\to k$ regular on $V'$ (with respect to the induced structure of $Y'$, i.e.\ locally a quotient of polynomials $g/h$, resp.\ homogeneous polynomials of the same degree, with $h\ne0$). We must show $f\circ\varphi'$ is regular on $\varphi'^{-1}(V')$.

\emph{Step 1: locally $f$ is the restriction of a regular function on $Y$.} Let $Q\in V'$. There are an open set $W'\subseteq V'$ containing $Q$ and polynomials $g,h$ (homogeneous of the same degree in the projective case) with $h\ne0$ on $W'$ and $f=g/h$ on $W'$. As $W'$ is open in $Y'$, $W'=W\cap Y'$ for some open $W\subseteq Y$. Put $V=W-Z(h)$, open in $Y$ and containing $Q$, and $F=g/h$, regular on $V$. Then $V\cap Y'\subseteq W'$ and $f=F$ on $V\cap Y'$.

\emph{Step 2: pull back.} Let $P\in\varphi'^{-1}(V')$ and $Q=\varphi(P)\in V'$; choose $V,F$ as in Step~1. Since $\varphi$ is a morphism, $G:=F\circ\varphi$ is regular on the open set $T=\varphi^{-1}(V)\subseteq X$, which contains $P$. For $x\in X'\cap T$ we have $\varphi(x)\in Y'\cap V$, so $f(\varphi(x))=F(\varphi(x))=G(x)$. Thus $f\circ\varphi'=G|_{X'\cap T}$ on the open neighbourhood $X'\cap T$ of $P$ in $X'$.

\emph{Step 3: restrictions of regular functions are regular on $X'$.} Since $G$ is regular at $P\in T$, there are an open $N\ni P$, $N\subseteq T$, and polynomials $a,b$ (homogeneous of the same degree in the projective case) with $b\neq0$ on $N$ and $G=a/b$ on $N$. Then on the open neighbourhood $X'\cap N$ of $P$ in $X'$, $f\circ\varphi'=a/b$ is a quotient of polynomials with nowhere vanishing denominator, i.e.\ regular at $P$ for the induced structure of $X'$.

Hence $f\circ\varphi'$ is regular at every point of $\varphi'^{-1}(V')$, and $\varphi'$ is a morphism. \hfill$\square$


\exercise{3.11}{Let $X$ be any variety and let $P\in X$. Show there is a 1-1 correspondence between the prime ideals of the local ring $\mathcal O_P$ and the closed subvarieties of $X$ containing $P$.}
\pf
Here ``closed subvariety'' means an irreducible closed subset of $X$ (with its induced structure), and prime ideals are all prime ideals, including $(0)$ and $\mathfrak m_P$.

\emph{Step 1: reduction to an affine neighbourhood.} We first find an open neighbourhood $U$ of $P$ in $X$ which is an affine variety. If $X\subseteq\mathbf P^n$, choose $i$ with $P\in U_i$; then $X\cap U_i$ is an open subset of $X$ and, via $\varphi_i$ (Prop.~3.3), a quasi-affine variety $X_1\subseteq\mathbf A^n$ containing $P$. (If $X$ is quasi-affine take $X_1=X$.) Let $Y_1=\overline{X_1}\subseteq\mathbf A^n$ (affine variety), so $X_1=Y_1-C$ with $C$ closed in $Y_1$, $P\notin C$. Then $I(C)\not\subseteq\mathfrak m_P$, so there is $g\in A(Y_1)$ with $g|_C=0$ and $g(P)\ne0$. The set $D(g)=Y_1-Z(g)\subseteq X_1$ is an open neighbourhood of $P$ and is isomorphic to $\{(y,z)\in Y_1\times\mathbf A^1:g(y)z=1\}$, a closed irreducible subset of $\mathbf A^{n+1}$ (via $y\mapsto(y,1/g(y))$ and projection, as in 3.6). So $D(g)$ is affine; call it $U$ (open in $X$, being open in $X_1\cong X\cap U_i$).

Closed irreducible subsets $Z\ni P$ of $X$ correspond bijectively to closed irreducible subsets of $U$ containing $P$: $Z\mapsto Z\cap U$ (nonempty open in irreducible $Z$, hence irreducible and dense in $Z$; closed in $U$), with inverse $T\mapsto\overline T$ (closure in $X$), because $\overline{Z\cap U}=Z$ and $\overline T\cap U=T$. Also $\mathcal O_{P,X}=\mathcal O_{P,U}$, since the local ring is defined by germs of functions on neighbourhoods of $P$. Both sides of the asserted correspondence are invariant under isomorphism, so we may replace $X$ by the affine variety $U$ (identified with its isomorphic copy in $\mathbf A^{n+1}$).

\emph{Step 2: the affine case.} Let $X\subseteq\mathbf A^m$ be affine with coordinate ring $A=A(X)$ and let $\mathfrak m=\mathfrak m_P\subseteq A$ be the maximal ideal of $P$. By Thm.~3.2(c), $\mathcal O_P\cong A_{\mathfrak m}$. We have the standard correspondences (Ch.~I, Section~1): $Z\mapsto I(Z)/I(X)$ is an inclusion-reversing bijection between closed subsets of $X$ and radical ideals of $A$, and under it $Z$ is irreducible iff the ideal is prime. Moreover $P\in Z$ iff $I(Z)\subseteq\mathfrak m_P$ (as $I(Z)\subseteq I(\{P\})$ means $P\in Z(I(Z))=Z$). So
\[
\{\text{closed irreducible }Z\subseteq X,\ P\in Z\}\ \longleftrightarrow\ \{\mathfrak p\subset A\text{ prime},\ \mathfrak p\subseteq\mathfrak m\}.
\]
By the standard correspondence for localization, $\mathfrak p\mapsto\mathfrak pA_{\mathfrak m}$ is a bijection from the primes of $A$ contained in $\mathfrak m$ (i.e.\ disjoint from $A-\mathfrak m$) onto the primes of $A_{\mathfrak m}$, with inverse $\mathfrak q\mapsto\mathfrak q\cap A$. Composing, $Z\mapsto I(Z)\mathcal O_P$ is a bijection (inclusion-reversing) from closed subvarieties of $X$ through $P$ onto prime ideals of $\mathcal O_P$. In terms of germs, $I(Z)\mathcal O_P$ consists of germs of the form $a/s$ with $a\in I(Z)$ and $s(P)\ne0$, i.e.\ germs at $P$ of regular functions vanishing on $Z$ near $P$. Under this bijection $Z=\{P\}$ corresponds to $\mathfrak m_P$ and $Z=X$ to $(0)$. \hfill$\square$

% Hartshorne Chapter I, Section 3 (Morphisms), Exercises 3.12--3.21

\exercise{3.12}{If $P$ is a point on a variety $X$, then $\dim\mathcal{O}_P=\dim X$. [Hint: Reduce to the affine case and use (3.2c).]}
\pf
\emph{Step 1: $\dim V=\dim X$ for every nonempty affine open $V\subseteq X$.} Every nonempty open subset of the irreducible space $X$ is dense, and $K(X)$ is the set of equivalence classes of pairs $\langle U,f\rangle$ with $U\subseteq X$ nonempty open; since $U\cap V\neq\emptyset$ for each such $U$, restriction to $V$ gives $K(V)=K(X)$. By Theorem 3.2(d) and Theorem 1.8A(a) applied to $V$ (isomorphic to an affine variety), $\dim V=\operatorname{tr.deg}_k K(V)=\operatorname{tr.deg}_kK(X)$. So all nonempty affine open subsets have the same dimension $d$. Now $X$ is covered by affine open sets (a point of $X$ lies in an open subset of a closed set of some $\mathbf A^n$ or $\mathbf P^n$; principal open sets $D(g)$ of an affine chart form a base of the topology, and $D(g)\cong Z(\ldots,\,gt-1)\subseteq\mathbf A^{n+1}$ is affine). By Ex.~1.10(b), $\dim X=\sup\dim V=d$.

\emph{Step 2.} Let $P\in X$ and choose an affine open neighbourhood $V$ of $P$, with an isomorphism $V\cong W$ onto an affine variety $W$ with coordinate ring $A=A(W)$. The local ring $\mathcal{O}_{P,X}$ consists of germs of regular functions at $P$ on open sets $U\ni P$; shrinking $U$ into $V$ shows $\mathcal{O}_{P,X}=\mathcal{O}_{P,V}$, and the isomorphism gives $\mathcal{O}_{P,V}\cong\mathcal{O}_{P',W}$ ($P'$ the image of $P$). By Theorem 3.2(c), $\mathcal{O}_{P',W}\cong A_{\fm_{P'}}$, and by Theorem 1.8A(b) $\dim A_{\fm_{P'}}=\operatorname{ht}\fm_{P'}=\dim A-\dim A/\fm_{P'}=\dim A=\dim W$. Hence $\dim\mathcal{O}_P=\dim W=\dim V=\dim X$ by Step 1.\hfill$\square$

\exercise{3.13}{\textbf{The Local Ring of a Subvariety.} Let $Y\subseteq X$ be a subvariety. Let $\mathcal{O}_{Y,X}$ be the set of equivalence classes $\langle U,f\rangle$ where $U\subseteq X$ is open, $U\cap Y\neq\emptyset$, and $f$ is a regular function on $U$. We say $\langle U,f\rangle$ is equivalent to $\langle V,g\rangle$ if $f=g$ on $U\cap V$. Show that $\mathcal{O}_{Y,X}$ is a local ring, with residue field $K(Y)$ and dimension $=\dim X-\dim Y$. It is the \emph{local ring of $Y$ on $X$}. Note if $Y=P$ is a point we get $\mathcal{O}_P$, and if $Y=X$ we get $K(X)$. Note also that if $Y$ is not a point, then $K(Y)$ is not algebraically closed, so in this way we get local rings whose residue fields are not algebraically closed.}
\pf
Recall that a subvariety is an irreducible locally closed subset, which is a variety by Ex.~3.10.

\emph{Realisation inside $K(X)$.} If $U,V$ are open sets meeting $Y$, then $U\cap V\cap Y\ne\emptyset$ because $Y$ is irreducible, so $U\cap V\neq\emptyset$. Since $X$ is irreducible, $U\cap V$ is dense in $U$, and two regular functions on $U$ that agree on a dense subset agree (the locus where they agree is closed, being the preimage of $0$ under the continuous map $f-g\colon U\to\mathbf A^1$). Hence the relation is an equivalence relation, a class $\langle U,f\rangle$ is the same as the element of $K(X)$ given by $f$, and
\[
\mathcal{O}_{Y,X}=\{\varphi\in K(X): \varphi\text{ is represented by a regular function on some open set meeting }Y\}\subseteq K(X)
\]
is a subring (add and multiply on $U\cap V$).

\emph{Reduction to the affine case.} Write $Y=\overline Y\cap O$ with $O\subseteq X$ open. Pick $P\in Y$ and an affine open neighbourhood $V$ of $P$ inside $O$ (as in 3.12, principal open sets of affine charts are affine and form a base). Then $Y'=V\cap Y=V\cap\overline Y$ is a closed irreducible subset of $V$, hence an affine variety, and it is a nonempty open subset of $Y$. Every element of $\mathcal{O}_{Y,X}$ has a representative on some $U\cap V$, which still meets $Y$; thus $\mathcal{O}_{Y,X}=\mathcal{O}_{Y',V}$ inside $K(X)=K(V)$, and $K(Y)=K(Y')$ (as in 3.12). So we may assume $X=V$ is affine with $A=A(V)$, $K=K(V)=\operatorname{Frac}A$, and $Y=Y'$ is closed in $V$ with prime ideal $\fq=I(Y)\subseteq A$.

\emph{Claim: $\mathcal{O}_{Y,X}=A_\fq$.} If $f/g\in A_\fq$ with $g\notin\fq$, then $g$ does not vanish identically on $Y$, so $D(g)$ meets $Y$, and $f/g$ is regular on $D(g)$; thus $A_\fq\subseteq \mathcal{O}_{Y,X}$. Conversely let $\varphi$ be regular on an open set $U$ meeting $Y$. Choose a principal open set $D(h)\subseteq U$ meeting $Y$, i.e.\ $h\notin\fq$. The set $D(h)$ is isomorphic to the closed subset $Z(1-ht)$ of $\mathbf A^{n+1}$, whose coordinate ring is $A[t]/(ht-1)\cong A_h$; by Theorem 3.2(a), $\mathcal{O}(D(h))=A_h$. So $\varphi=f/h^m$ with $f\in A$, whence $\varphi\in A_\fq$.

\emph{Conclusion.} $A_\fq$ is a local ring with maximal ideal $\fq A_\fq$ and residue field $\operatorname{Frac}(A/\fq)=\operatorname{Frac}A(Y)=K(Y)$ by Theorem 3.2(d). (Equivalently the residue map is $\langle U,f\rangle\mapsto f|_{U\cap Y}\in K(Y)$.) Finally $\dim A_\fq=\operatorname{ht}\fq=\dim A-\dim A/\fq$ by Theorem 1.8A(b), and $\dim A=\dim V=\dim X$, $\dim A/\fq=\dim Y'=\dim Y$ by Step 1 of 3.12 (applied to $X$ and $Y$). Hence $\dim\mathcal{O}_{Y,X}=\dim X-\dim Y$.

\emph{Special cases.} For $Y=\{P\}$ we get $\mathcal{O}_{P}$ (no condition other than $P\in U$), and for $Y=X$ every nonempty open $U$ is allowed, giving $K(X)$.

\emph{$K(Y)$ is not algebraically closed if $\dim Y=d\ge1$.} $K(Y)$ is a finitely generated extension of $k$ of transcendence degree $d$ (Thm 3.2(d) and Thm 1.8A(a)), so it is algebraic, hence finite, over $F=k(t_1,\dots,t_d)$ for a transcendence basis $t_i$. If $K(Y)$ were algebraically closed it would be an algebraic closure of $F$. But $F$ has algebraic extensions of arbitrarily large degree ($T^n-t_1$ is irreducible over $F$ by Eisenstein at the prime $t_1$ of $k[t_2,\dots,t_d][t_1]$, and Gauss' lemma), so $\overline F$ is not finite over $F$, a contradiction.\hfill$\square$

\exercise{3.14}{\textbf{Projection from a Point.} Let $\mathbf P^n$ be a hyperplane in $\mathbf P^{n+1}$ and let $P\in\mathbf P^{n+1}-\mathbf P^n$. Define a mapping $\varphi\colon\mathbf P^{n+1}-\{P\}\to\mathbf P^n$ by $\varphi(Q)=$ the intersection of the unique line containing $P$ and $Q$ with $\mathbf P^n$.
\begin{enumerate}[label=(\alph*),leftmargin=*]
\item Show that $\varphi$ is a morphism.
\item Let $Y\subseteq\mathbf P^3$ be the twisted cubic curve which is the image of the $3$-uple embedding of $\mathbf P^1$ (Ex.~2.12). If $t,u$ are the homogeneous coordinates on $\mathbf P^1$, we say that $Y$ is the curve given parametrically by $(x,y,z,w)=(t^3,t^2u,tu^2,u^3)$. Let $P=(0,0,1,0)$, and let $\mathbf P^2$ be the hyperplane $z=0$. Show that the projection of $Y$ from $P$ is a cuspidal cubic curve in the plane, and find its equation.
\end{enumerate}}
\pf
\textbf{(a)} Let $H=\mathbf P^n$ be defined by a linear form $\ell$ and let $P=[v]$; since $P\notin H$, $\ell(v)\ne0$ and $v\notin\ker\ell$. Choose a basis of $\ker \ell$ (dimension $n+1$) and add $v$; in the corresponding linear coordinates $x_0,\dots,x_{n+1}$ we have $H=Z(x_{n+1})$ and $P=(0,\dots,0,1)$. A linear change of coordinates is an automorphism of $\mathbf P^{n+1}$ (the map and its inverse are given by linear forms, hence are morphisms, see (b) below for why such maps are morphisms), and it carries lines to lines, so we may assume these coordinates.

Let $Q=(a_0:\dots:a_{n+1})\ne P$; then $(a_0,\dots,a_n)\ne 0$. The line through $P$ and $Q$ consists of the points $s(0,\dots,0,1)+\lambda(a_0,\dots,a_{n+1})$; it meets $x_{n+1}=0$ where $s=-\lambda a_{n+1}$, i.e.\ at $(a_0:\dots:a_n:0)$. So
\[
\varphi(a_0:\dots:a_{n+1})=(a_0:\dots:a_n).
\]
This is well defined exactly off $P$. It is a morphism: let $F$ be a homogeneous polynomial in $x_0,\dots,x_n$; then $\varphi^{-1}(Z(F))=Z(F(x_0,\dots,x_n))-\{P\}$ is closed, so $\varphi$ is continuous. If $g$ is regular on an open $W\subseteq\mathbf P^n$, then near a point of $\varphi^{-1}(W)$ we have $g=G/H'$ with $G,H'$ homogeneous of the same degree, $H'$ nonvanishing there, and $g\circ\varphi=G(x_0,\dots,x_n)/H'(x_0,\dots,x_n)$, which is again a quotient of homogeneous polynomials of the same degree with nonvanishing denominator, hence regular. (This is the standard fact that $(F_0:\dots:F_n)$ with the $F_i$ homogeneous of the same degree and without common zero on an open set is a morphism.)

\textbf{(b)} In the coordinates $(x,y,z,w)$ with $P=(0,0,1,0)$ and $H=Z(z)$ the projection is $\varphi(x:y:z:w)=(x:y:w)\in\mathbf P^2$ (coordinates $a,b,c$ on $\mathbf P^2$). Note $P\notin Y$, since $t^3=u^3=0$ forces $t=u=0$. So $\varphi|_Y$ is defined and is a morphism (Ex.~3.10), and on a point of $Y$ it is
\[
(t:u)\longmapsto (a:b:c)=(t^3:t^2u:u^3).
\]
Each such point satisfies $a^2c=b^3$, since $t^6u^3=t^6u^3$. Conversely we show every point of $C=Z(a^2c-b^3)\subseteq\mathbf P^2$ is in the image. If $c=0$ then $b=0$ and the point is $(1:0:0)$, the image of $(t:u)=(1:0)$. If $c=1$ and $b=0$, then $a=0$ and $(0:0:1)$ is the image of $(0:1)$. If $c=1$ and $b\ne0$, then $a^2=b^3\neq 0$; put $t=a/b$, $u=1$: $t^2=a^2/b^2=b$ and $t^3=a^3/b^3=a^3/a^2=a$, so $(t^3:t^2u:u^3)=(a:b:1)$. Hence $\varphi(Y)=C$.

The polynomial $a^2c-b^3$ is irreducible (it is of degree $1$ in $c$ with coprime coefficients $a^2$ and $b^3$), so $C$ is an irreducible plane cubic. In the affine chart $c=1$ it is $a^2=b^3$, the cuspidal cubic (with its cusp at the origin, i.e.\ at $(0:0:1)$), and the point at infinity $(1:0:0)$ is a smooth point. Thus the projection of $Y$ from $P$ is the cuspidal cubic curve
\[
x^2w=y^3\quad\text{in the plane } z=0,
\]
that is, $a^2c=b^3$ in the coordinates $(a:b:c)=(x:y:w)$.\hfill$\square$

\exercise{3.15}{\textbf{Products of Affine Varieties.} Let $X\subseteq\mathbf A^n$ and $Y\subseteq\mathbf A^m$ be affine varieties.
\begin{enumerate}[label=(\alph*),leftmargin=*]
\item Show that $X\times Y\subseteq\mathbf A^{n+m}$ with its induced topology is irreducible. [Hint: Suppose that $X\times Y$ is a union of two closed subsets $Z_1\cup Z_2$. Let $X_i=\{x\in X\mid x\times Y\subseteq Z_i\}$, $i=1,2$. Then $X=X_1\cup X_2$, and $X_1,X_2$ are closed. Hence $X=X_1$ or $X_2$, so $X\times Y=Z_1$ or $Z_2$.] The affine variety $X\times Y$ is called the \emph{product} of $X$ and $Y$. Note that its topology is in general not equal to the product topology (Ex.~1.4).
\item Show that $A(X\times Y)\cong A(X)\otimes_kA(Y)$.
\item Show that $X\times Y$ is a product in the category of varieties, i.e., show (i) the projections $X\times Y\to X$ and $X\times Y\to Y$ are morphisms, and (ii) given a variety $Z$, and the morphisms $Z\to X$, $Z\to Y$, there is a unique morphism $Z\to X\times Y$ making a commutative diagram with the projections.
\item Show that $\dim X\times Y=\dim X+\dim Y$.
\end{enumerate}}
\pf
Let $x=(x_1,\dots,x_n)$ and $y=(y_1,\dots,y_m)$ be the coordinates on $\mathbf A^n\times\mathbf A^m=\mathbf A^{n+m}$.

\textbf{(a)} First, $X\times Y$ is closed in $\mathbf A^{n+m}$: if $X=Z(S)$ and $Y=Z(T)$ with $S\subseteq k[x]$, $T\subseteq k[y]$, then $X\times Y=Z(S\cup T)$.

For $x_0\in X$ the map $j_{x_0}\colon Y\to \mathbf A^{n+m}$, $y\mapsto(x_0,y)$, is given by polynomials, hence continuous, with image $\{x_0\}\times Y=Z(x_i-x_{0,i},\ T)$, and its inverse (the projection to the $y$-coordinates) is also polynomial. So $j_{x_0}$ is a homeomorphism of $Y$ onto $\{x_0\}\times Y$ and $\{x_0\}\times Y$ is irreducible. Likewise for $y_0\in Y$, $i^{y_0}\colon X\to X\times Y$, $x\mapsto(x,y_0)$, is continuous.

Suppose $X\times Y=Z_1\cup Z_2$ with $Z_i$ closed. For each $x_0\in X$, the irreducible set $\{x_0\}\times Y$ is the union of its closed subsets $(\{x_0\}\times Y)\cap Z_i$, so it lies in $Z_1$ or in $Z_2$. Put $X_i=\{x\in X:\{x\}\times Y\subseteq Z_i\}$; then $X=X_1\cup X_2$. Moreover
\[
X_i=\bigcap_{y_0\in Y}(i^{y_0})^{-1}(Z_i)
\]
is closed in $X$ as an intersection of closed sets. Since $X$ is irreducible, $X=X_1$ or $X=X_2$, i.e.\ $X\times Y=Z_1$ or $Z_2$. Also $X\times Y\ne\emptyset$. So $X\times Y$ is irreducible, hence an affine variety.

The topology is not the product topology: for $X=Y=\mathbf A^1$ the diagonal $Z(x_1-y_1)$ is closed in $\mathbf A^2$, but the proper closed sets of the product topology are finite unions of points, vertical lines $\{a\}\times\mathbf A^1$ and horizontal lines $\mathbf A^1\times\{b\}$ (Ex.~1.4), and the diagonal is not of this form.

\textbf{(b)} Put $I=I(X)\subseteq k[x]$, $J=I(Y)\subseteq k[y]$. Then
\[
A(X)\otimes_kA(Y)=k[x]/I\otimes_kk[y]/J\cong k[x,y]/(I+J),
\]
where $I+J$ denotes the ideal of $k[x,y]$ generated by $I$ and $J$. It therefore suffices to prove $I(X\times Y)=I+J$. Clearly $I+J\subseteq I(X\times Y)$. Conversely let $f\in k[x,y]$ vanish on $X\times Y$. In $A(X)\otimes_kA(Y)$ the image $\bar f$ of $f$ can be written $\bar f=\sum_{j=1}^rg_j\otimes h_j$ with $g_j\in A(X)$ and with $h_1,\dots,h_r\in A(Y)$ linearly independent over $k$ (choose a basis of the finite-dimensional span of the second factors that occur). For $x_0\in X$ and $y_0\in Y$ we have $0=f(x_0,y_0)=\sum_jg_j(x_0)h_j(y_0)$. Thus the element $\sum_jg_j(x_0)h_j\in A(Y)$, viewed as a function on $Y$ (since $k$ is algebraically closed, $A(Y)$ is the ring of polynomial functions on $Y$), vanishes identically, hence is $0$ in $A(Y)$; by linear independence $g_j(x_0)=0$ for all $j$. As $x_0\in X$ was arbitrary, $g_j=0$ in $A(X)$. Hence $\bar f=0$, i.e.\ $f\in I+J$. So $I(X\times Y)=I+J$ and $A(X\times Y)\cong A(X)\otimes_kA(Y)$.

\textbf{(c)} (i) The projection $\pi_1\colon X\times Y\to X$ is the restriction of the polynomial map $(x,y)\mapsto x$; by Lemma 3.6 (coordinate functions $x_i\circ\pi_1$ are regular) it is a morphism, and likewise $\pi_2$.

(ii) Let $Z$ be a variety and $\varphi\colon Z\to X$, $\psi\colon Z\to Y$ morphisms. Define $\theta\colon Z\to X\times Y\subseteq\mathbf A^{n+m}$ by $\theta(z)=(\varphi(z),\psi(z))$. Its coordinate functions are $x_i\circ\varphi$ and $y_j\circ\psi$, which are regular on $Z$, so $\theta$ is a morphism by Lemma 3.6 (applied to the affine variety $X\times Y$). Clearly $\pi_1\theta=\varphi$ and $\pi_2\theta=\psi$. If $\theta'$ is another morphism with $\pi_1\theta'=\varphi$, $\pi_2\theta'=\psi$, then as a map of sets $\theta'(z)=(\varphi(z),\psi(z))=\theta(z)$; a morphism is determined by its underlying map, so $\theta'=\theta$.

\textbf{(d)} Let $r=\dim X$ and $s=\dim Y$. By Noether normalisation (Thm 1.4A) there are algebraically independent $a_1,\dots,a_r\in A(X)$ and $b_1,\dots,b_s\in A(Y)$ such that $A(X)$ is integral over $k[a]=k[a_1,\dots,a_r]$ and $A(Y)$ is integral over $k[b]$ (where $r=\operatorname{tr.deg}K(X)$ and $s=\operatorname{tr.deg}K(Y)$ by Thm 1.8A). Tensoring over the field $k$ preserves injections, so $k[a]\otimes_kk[b]=k[a,b]$, a polynomial ring in $r+s$ variables, embeds in $A(X)\otimes_kA(Y)=A(X\times Y)$. The latter is generated as a ring by the elements $g\otimes1$ and $1\otimes h$, which are integral over $k[a,b]$; the integral elements form a subring, so $A(X\times Y)$ is integral over $k[a,b]$. Thus $K(X\times Y)$ is algebraic over $k(a,b)$, which has transcendence degree $r+s$, so $\dim X\times Y=\operatorname{tr.deg}_kK(X\times Y)=r+s$ (Thm 1.8A).\hfill$\square$

\exercise{3.16}{\textbf{Products of Quasi-Projective Varieties.} Use the Segre embedding (Ex.~2.14) to identify $\mathbf P^n\times\mathbf P^m$ with its image and hence give it a structure of projective variety. Now for any two quasi-projective varieties $X\subseteq\mathbf P^n$ and $Y\subseteq\mathbf P^m$, consider $X\times Y\subseteq\mathbf P^n\times\mathbf P^m$.
\begin{enumerate}[label=(\alph*),leftmargin=*]
\item Show that $X\times Y$ is a quasi-projective variety.
\item If $X,Y$ are both projective, show that $X\times Y$ is projective.
\item Show that $X\times Y$ is a product in the category of varieties. (Starred in the book.)
\end{enumerate}}
\pf
Let $\psi\colon\mathbf P^n\times\mathbf P^m\to\mathbf P^N$, $N=nm+n+m$, be the Segre map $((x_i),(y_j))\mapsto(x_iy_j)$, with image $W=Z(z_{ij}z_{kl}-z_{il}z_{kj})$ (Ex.~2.14), and give $W$ and all its subsets the induced structure. Throughout, $x_i,y_j,z_{ij}$ are the coordinates.

\emph{Lemma A. If $A\subseteq\mathbf P^n$ and $B\subseteq\mathbf P^m$ are closed, then $A\times B$ is closed in $W$.} Let $A=Z(F_\alpha)$, $B=Z(G_\beta)$ with $F_\alpha,G_\beta$ homogeneous of degrees $d_\alpha$, $e_\beta$. For each $j$ the polynomial $F_\alpha(x)y_j^{d_\alpha}$ is bihomogeneous of bidegree $(d_\alpha,d_\alpha)$; each monomial $x^{a}y_j^{d}$ is a product of $d$ monomials $x_{i}y_j=z_{ij}$, so $F_\alpha(x)y_j^{d_\alpha}=\tilde F_{\alpha,j}(z)$ for a homogeneous polynomial $\tilde F_{\alpha,j}$ of degree $d_\alpha$ in the $z_{ij}$. Similarly $G_\beta(y)x_i^{e_\beta}=\tilde G_{\beta,i}(z)$. For a point $\psi(x,y)\in W$, all $\tilde F_{\alpha,j}$ vanish iff $F_\alpha(x)y_j^{d_\alpha}=0$ for all $j$ iff $F_\alpha(x)=0$ (some $y_j\ne0$); similarly for $G_\beta$. So $A\times B=W\cap Z(\tilde F_{\alpha,j},\tilde G_{\beta,i})$ is closed.

\emph{Lemma B (slices).} For $a\in\mathbf P^n$ the map $j_a\colon\mathbf P^m\to\mathbf P^N$, $y\mapsto(a_iy_j)$, is the restriction of an injective linear map $k^{m+1}\to k^{N+1}$, i.e.\ an isomorphism of $\mathbf P^m$ onto a linear subspace $L_a\subseteq\mathbf P^N$; it lands in $W$, with image $\{a\}\times\mathbf P^m$. So $j_a$ is a homeomorphism onto $\{a\}\times\mathbf P^m$, and $\{a\}\times Y\cong Y$ is irreducible for $Y\subseteq\mathbf P^m$ irreducible. Similarly for $b\in\mathbf P^m$ the map $i^b\colon\mathbf P^n\to W$, $x\mapsto(x_ib_j)$, is continuous.

\textbf{(a)} Write $X=\overline X\cap(\mathbf P^n-Z_1)$ and $Y=\overline Y\cap(\mathbf P^m-Z_2)$ with $Z_1,Z_2$ closed. Then
\[
X\times Y=(\overline X\times\overline Y)\cap\bigl(W-(Z_1\times\mathbf P^m)\bigr)\cap\bigl(W-(\mathbf P^n\times Z_2)\bigr)
\]
is the intersection of a closed subset of $W$ (Lemma A) with two open subsets, hence locally closed in $\mathbf P^N$.

It remains to see that $X\times Y$ is irreducible; we repeat the argument of Ex.~3.15(a). Suppose $X\times Y=C_1\cup C_2$ with $C_i$ closed in $X\times Y$. For $x\in X$, the set $\{x\}\times Y$ is irreducible (Lemma B) and is covered by the closed sets $(\{x\}\times Y)\cap C_i$, so lies in $C_1$ or $C_2$. Put $X_i=\{x\in X:\{x\}\times Y\subseteq C_i\}$. Then $X=X_1\cup X_2$ and $X_i=\bigcap_{y\in Y}(i^y)^{-1}(C_i)$ is closed in $X$ (the maps $i^y$ are continuous). Since $X$ is irreducible, $X=X_i$ for some $i$, so $X\times Y=C_i$. Thus $X\times Y$ is an irreducible locally closed subset of $\mathbf P^N$, i.e.\ a quasi-projective variety.

\textbf{(b)} If $X,Y$ are closed in $\mathbf P^n,\mathbf P^m$, then $X\times Y$ is closed in $W$ by Lemma A, hence closed in $\mathbf P^N$, and it is irreducible by (a). So it is a projective variety.

\textbf{(c)} \emph{(i) The projections are morphisms.} Let $W_j=\{y_j\neq0\}\subseteq W$; this is the open set where $(z_{0j},\dots,z_{nj})\ne0$, and on it $\pi_1(x,y)=(x_0y_j:\dots:x_ny_j)=(z_{0j}:\dots:z_{nj})$. The same argument as in 3.14(a) (homogeneous polynomials of the same degree without common zero) shows $\pi_1|_{W_j}$ is a morphism to $\mathbf P^n$. Since the $W_j$ ($j=0,\dots,m$) cover $W$ and being a morphism is local on the source, $\pi_1$ is a morphism; similarly $\pi_2=(z_{i0}:\dots:z_{im})$ on $\{x_i\ne0\}$. By restriction (Ex.~3.10) the projections $X\times Y\to X$, $X\times Y\to Y$ are morphisms.

\emph{(ii) Universal property.} Let $Z$ be a variety, and $\varphi\colon Z\to X$, $\rho\colon Z\to Y$ morphisms. Define $\theta(z)=(\varphi(z),\rho(z))\in X\times Y$. We show $\theta\colon Z\to\mathbf P^N$ is a morphism. Fix $z_0\in Z$ and indices $i,j$ with $x_i(\varphi(z_0))\ne0$, $y_j(\rho(z_0))\neq0$. Let $Z'=\varphi^{-1}(\{x_i\ne0\})\cap\rho^{-1}(\{y_j\neq0\})$, an open neighbourhood of $z_0$ (the maps are continuous). On $Z'$ the functions $f_k=(x_k/x_i)\circ\varphi$ and $g_l=(y_l/y_j)\circ\rho$ are regular, and
\[
\theta(z)=\bigl(x_k(\varphi(z))\,y_l(\rho(z))\bigr)_{k,l}=\bigl(f_k(z)g_l(z)\bigr)_{k,l}\quad\text{(after dividing by }x_iy_j).
\]
So $\theta|_{Z'}$ lands in the standard open set $\{z_{ij}\neq0\}\cong\mathbf A^N$ and is given there by the regular functions $f_kg_l$; by Lemma 3.6 it is a morphism. Since the $Z'$ cover $Z$, $\theta$ is a morphism $Z\to\mathbf P^N$ with image in the locally closed subvariety $X\times Y$, hence a morphism $Z\to X\times Y$ by Ex.~3.10. Clearly $\pi_1\theta=\varphi$, $\pi_2\theta=\rho$. Uniqueness: any $\theta'$ with these properties has the same underlying map of sets, $z\mapsto(\varphi(z),\rho(z))$, hence $\theta'=\theta$.\hfill$\square$

\exercise{3.17}{\textbf{Normal Varieties.} A variety $Y$ is \emph{normal at a point $P\in Y$} if $\mathcal{O}_P$ is an integrally closed ring. $Y$ is \emph{normal} if it is normal at every point.
\begin{enumerate}[label=(\alph*),leftmargin=*]
\item Show that every conic in $\mathbf P^2$ is normal.
\item Show that the quadric surfaces $Q_1,Q_2$ in $\mathbf P^3$ given by equations $Q_1\colon xy=zw$; $Q_2\colon xy=z^2$ are normal (cf.\ II, Ex.~6.4 for the latter).
\item Show that the cuspidal cubic $y^2=x^3$ in $\mathbf A^2$ is not normal.
\item If $Y$ is affine, then $Y$ is normal $\Leftrightarrow$ $A(Y)$ is integrally closed.
\item Let $Y$ be an affine variety. Show that there is a normal affine variety $\tilde Y$, and a morphism $\pi\colon\tilde Y\to Y$, with the property that whenever $Z$ is a normal variety, and $\varphi\colon Z\to Y$ is a dominant morphism (i.e., $\varphi(Z)$ is dense in $Y$), then there is a unique morphism $\theta\colon Z\to\tilde Y$ such that $\varphi=\pi\circ\theta$. $\tilde Y$ is called the \emph{normalization} of $Y$. You will need (3.9A) above.
\end{enumerate}}
\pf
Throughout, ``integrally closed'' for $\mathcal{O}_P$ means integrally closed in its field of fractions, which is $K(Y)$ (Thm 3.2(d) and the same for any variety, via an affine neighbourhood). Two standard facts are used: a localisation of an integrally closed domain is integrally closed, and a UFD is integrally closed. Normality at $P$ depends only on the isomorphism class of the local ring, hence is invariant under isomorphism of varieties and may be tested on any open neighbourhood of $P$.

\textbf{(d)} (We do this first.) Let $A=A(Y)$, $K=\operatorname{Frac}A$. ($\Leftarrow$) If $A$ is integrally closed, so is each $A_{\fm_P}\cong\mathcal{O}_P$ (Thm 3.2(c)). ($\Rightarrow$) Suppose all $\mathcal{O}_P=A_{\fm_P}$ are integrally closed, and let $f\in K$ be integral over $A$. Then $f$ is integral over each $A_{\fm_P}$, so $f\in A_{\fm_P}$ for all $P$, hence $f\in\bigcap_{\fm}A_\fm=A$ (the last equality is shown in the proof of Thm 3.2(a)). So $A$ is integrally closed.

\textbf{(a)} By Ex.~3.1(c), every (irreducible) conic in $\mathbf P^2$ is isomorphic to $\mathbf P^1$, so it suffices to show $\mathbf P^1$ is normal. $\mathbf P^1$ is covered by the two open sets $\{t\ne0\}\cong\mathbf A^1$ and $\{u\ne0\}\cong\mathbf A^1$ (Prop.~3.3). By (d), $\mathbf A^1$ is normal since $A(\mathbf A^1)=k[T]$ is a UFD, hence integrally closed. So $\mathbf P^1$, and thus every conic, is normal.

\textbf{(b)} $Q_1$ and $Q_2$ are varieties (the defining polynomials are irreducible), and every point lies in one of the standard open sets, so we test normality on affine charts.

$Q_1$: on $\{x\ne0\}$ put $x=1$: the chart is $Z(y-zw)\subseteq\mathbf A^3$ with coordinates $y,z,w$, isomorphic to $\mathbf A^2$ via $(z,w)$ (inverse $(z,w)\mapsto(zw,z,w)$). Likewise $\{y\ne0\}$ ($x=zw$), $\{z\neq0\}$ ($w=xy$), $\{w\ne0\}$ ($z=xy$) are isomorphic to $\mathbf A^2$. These four sets cover $Q_1$, and $A(\mathbf A^2)=k[s,t]$ is a UFD, so $Q_1$ is normal by (d).

$Q_2$: the chart $\{x\neq0\}$ is $y=z^2$ in $\mathbf A^3$, isomorphic to $\mathbf A^2$ with coordinates $(z,w)$; similarly $\{y\ne0\}$. A point of $Q_2$ with $x=y=0$ has $z=0$, so equals $(0:0:0:1)$, and is in the chart $\{w\ne0\}$, which is the affine surface $Z(xy-z^2)\subseteq\mathbf A^3$ (the polynomial $xy-z^2$ is irreducible, so by the Nullstellensatz $A=k[x,y,z]/(xy-z^2)$). It remains to show that $A$ is integrally closed, and then use (d). Define $\alpha\colon A\to k[s,t]$ by $x\mapsto s^2$, $y\mapsto t^2$, $z\mapsto st$ (it kills $xy-z^2$). Using $z^2=xy$, $A$ is spanned by the monomials $x^ay^bz^e$ with $e\in\{0,1\}$, and $\alpha$ sends these to $s^{2a+e}t^{2b+e}$; these are pairwise distinct monomials (the parity of the exponent of $s$ gives $e$, then $a,b$), so they are linearly independent and $\alpha$ is injective. The image is $S=\operatorname{span}\{s^pt^q:p+q\text{ even}\}=\bigoplus_{d\text{ even}}k[s,t]_d$. Let $f\in\operatorname{Frac}S$ be integral over $S$. Then $f\in\operatorname{Frac}(k[s,t])$ is integral over $k[s,t]$, hence $f\in k[s,t]$. Write $f=f_{ev}+f_{odd}$ for its parts of even and odd total degree, and $f=a/b$ with $a,b\in S$, $b\neq0$. Then $a=fb=f_{ev}b+f_{odd}b$, where $f_{ev}b$ has only even-degree components and $f_{odd}b$ only odd-degree ones; since $a$ has only even-degree components, $f_{odd}b=0$, so $f_{odd}=0$ and $f\in S$. Hence $S\cong A$ is integrally closed, and $Q_2$ is normal.

\textbf{(c)} Let $Y=Z(y^2-x^3)\subseteq\mathbf A^2$; the polynomial is irreducible, so $A(Y)=k[x,y]/(y^2-x^3)$. As above, $A(Y)$ is spanned by $x^a$, $x^ay$, which the map $x\mapsto T^2$, $y\mapsto T^3$ sends to $T^{2a}$, $T^{2a+3}$: distinct monomials, so $A(Y)\cong R=k[T^2,T^3]=\operatorname{span}\{T^i:i\ne1\}\subseteq k[T]$. The element $T\in\operatorname{Frac}R$ (as $T=T^3/T^2$) is integral over $R$ (it is a root of $X^2-T^2\in R[X]$). We show $T\notin R_{\fm}$, $\fm=(T^2,T^3)$, the local ring $\mathcal{O}_{(0,0)}$: if $T=a/b$ with $a,b\in R$, $b\notin\fm$, then $b$ has nonzero constant term $b_0$ and $a=Tb$ has a nonzero coefficient $b_0$ in degree $1$, impossible for $a\in R$. So $\mathcal{O}_{(0,0)}$ is not integrally closed, and $Y$ is not normal at the origin (equivalently, $A(Y)$ is not integrally closed, by (d); in the coordinates of $Y$ the element is $y/x$).

\textbf{(e)} Let $A=A(Y)$, $K=\operatorname{Frac}A$ and let $A'$ be the integral closure of $A$ in $K$. By Theorem 3.9A (with $L=K$), $A'$ is a finitely generated $A$-module, hence a finitely generated $k$-algebra; it is a domain with $A\subseteq A'\subseteq K$, so $\operatorname{Frac}A'=K$, and $A'$ is integrally closed in $K$ (integral over integral over $A$ is integral over $A$). By Cor.~3.8 (a finitely generated domain over $k=\bar k$ is $k[x_1,\dots,x_m]/\mathfrak p$ with $\mathfrak p$ prime, i.e.\ the coordinate ring of an affine variety) there is an affine variety $\tilde Y$ with $A(\tilde Y)\cong A'$, and $\tilde Y$ is normal by (d). The inclusion $\iota\colon A\to A'$ is a homomorphism of $k$-algebras, so by Prop.~3.5 (Cor.~3.8) it corresponds to a unique morphism $\pi\colon\tilde Y\to Y$ with $\pi^*=\iota$.

\emph{Universal property.} Let $Z$ be a normal variety and $\varphi\colon Z\to Y$ a dominant morphism. The map $\varphi^*\colon A=\mathcal{O}(Y)\to\mathcal{O}(Z)$, $f\mapsto f\circ\varphi$ is an injective $k$-algebra homomorphism (if $f\circ\varphi=0$ then $f$ vanishes on the dense set $\varphi(Z)$, so $f=0$). Also $\mathcal{O}(Z)\subseteq\mathcal{O}_{Q,Z}\subseteq K(Z)$ for each $Q\in Z$, since $Z$ is irreducible. Hence $\varphi^*$ extends to a field embedding $\varphi^*\colon K\to K(Z)$. Let $g\in A'$, say $g^m+a_1g^{m-1}+\dots+a_m=0$ with $a_i\in A$. Then $\tilde g=\varphi^*(g)\in K(Z)$ is integral over $\varphi^*(A)\subseteq\mathcal{O}_{Q,Z}$ for every $Q\in Z$. Since $Z$ is normal, $\mathcal{O}_{Q,Z}$ is integrally closed in $K(Z)$, so $\tilde g\in\mathcal{O}_{Q,Z}$: $\tilde g$ is represented by a regular function $h_Q$ on an open neighbourhood $U_Q$ of $Q$. Any two representatives agree on the (nonempty, dense) intersection $U_Q\cap U_{Q'}$ since they define the same element of $K(Z)$, so the $h_Q$ glue to a regular function on $Z$, i.e.\ $\tilde g\in\mathcal{O}(Z)$. Thus $\varphi^*$ restricts to a $k$-algebra homomorphism $\Psi\colon A'\to\mathcal{O}(Z)$ with $\Psi\circ\iota=\varphi^*$. By Prop.~3.5 there is a morphism $\theta\colon Z\to\tilde Y$ with $\theta^*=\Psi$, and then $(\pi\circ\theta)^*=\theta^*\circ\pi^*=\Psi\circ\iota=\varphi^*$; since $\alpha\colon\varphi\mapsto\varphi^*$ is a bijection $\operatorname{Hom}(Z,Y)\to\operatorname{Hom}(A,\mathcal{O}(Z))$ (Prop.~3.5), $\pi\circ\theta=\varphi$.

\emph{Uniqueness.} Suppose $\theta'\colon Z\to\tilde Y$ also satisfies $\pi\circ\theta'=\varphi$. Then $\theta'^*\colon A'\to\mathcal{O}(Z)\subseteq K(Z)$ satisfies $\theta'^*|_A=\varphi^*$. For $g\in A'\subseteq K$ write $g=a/b$ with $a,b\in A$, $b\ne0$; then $bg=a$ in $A'$, so $\theta'^*(b)\theta'^*(g)=\theta'^*(a)$, i.e.\ $\varphi^*(b)\theta'^*(g)=\varphi^*(a)$ in the field $K(Z)$, where $\varphi^*(b)\neq0$. Hence $\theta'^*(g)=\varphi^*(a)/\varphi^*(b)=\Psi(g)$. So $\theta'^*=\Psi=\theta^*$ and $\theta'=\theta$ by the bijectivity in Prop.~3.5.\hfill$\square$

\exercise{3.18}{\textbf{Projectively Normal Varieties.} A projective variety $Y\subseteq\mathbf P^n$ is \emph{projectively normal} (with respect to the given embedding) if its homogeneous coordinate ring $S(Y)$ is integrally closed.
\begin{enumerate}[label=(\alph*),leftmargin=*]
\item If $Y$ is projectively normal, then $Y$ is normal.
\item There are normal varieties in projective space which are not projectively normal. For example, let $Y$ be the twisted quartic curve in $\mathbf P^3$ given parametrically by $(x,y,z,w)=(t^4,t^3u,tu^3,u^4)$. Then $Y$ is normal but not projectively normal. See (III, Ex.~5.6) for more examples.
\item Show that the twisted quartic curve $Y$ above is isomorphic to $\mathbf P^1$, which is projectively normal. Thus projective normality depends on the embedding.
\end{enumerate}}
\pf
\textbf{(a)} Let $S=S(Y)$, a graded domain, integrally closed by hypothesis, and let $Y_i=Y\cap\{x_i\neq0\}$. By the proof of Thm 3.4, $A(Y_i)\cong S_{(x_i)}$, the degree-$0$ part of $S_{x_i}$. Since every point of $Y$ lies in some $Y_i$ and the $Y_i$ are affine, by 3.17(d) it suffices to show that $S_{(x_i)}$ is integrally closed. The ring $S_{x_i}$ is a localisation of the integrally closed domain $S$, hence integrally closed in $\operatorname{Frac}S$. Let $f\in\operatorname{Frac}S_{(x_i)}$ be integral over $S_{(x_i)}$. Then $f\in\operatorname{Frac}S$ is integral over $S_{x_i}$, so $f\in S_{x_i}$. Write $f=a/b$ with $a,b\in S_{(x_i)}$, $b\neq0$, and let $f=\sum_df_d$ be the decomposition into homogeneous components of $S_{x_i}$ (a $\mathbf Z$-graded domain). From $bf=a$ and $\deg a=\deg b=0$ we get $bf_d=0$ for $d\neq0$, so $f_d=0$ for $d\ne0$ and $f=f_0\in S_{(x_i)}$. Hence $A(Y_i)$ is integrally closed, each $Y_i$ is normal, and so is $Y$.

\textbf{(c)} Define $\varphi\colon\mathbf P^1\to\mathbf P^3$, $\varphi(t:u)=(t^4:t^3u:tu^3:u^4)$. It is a morphism (the four forms have degree $4$ and no common zero: $t^4=u^4=0$ forces $t=u=0$), and $Y=\varphi(\mathbf P^1)$.

\emph{$Y$ is closed.} Points of $Y$ satisfy $xw=yz$, $y^3=x^2z$, $z^3=yw^2$ (e.g.\ $y^3=t^9u^3=x^2z$ and $z^3=t^3u^9=yw^2$). Conversely let $(x:y:z:w)$ satisfy these equations. If $x\neq0$, take $x=1$: then $z=y^3$, $w=yz=y^4$, so the point is $(1:y:y^3:y^4)=\varphi(1:y)$. If $x=0$, then $y^3=0$, so $y=0$, then $z^3=0$, so the point is $(0:0:0:1)=\varphi(0:1)$. Hence $Y=Z(xw-yz,\,y^3-x^2z,\,z^3-yw^2)$ is closed and, as the image of $\mathbf P^1$, irreducible: a projective variety.

\emph{Inverse.} On $Y$ we have $x=y=0$ only if $t=0$, so $Y=(Y\cap\{x\neq0\})\cup(Y\cap\{w\ne0\})$. Define $\rho\colon Y\to\mathbf P^1$ by $(x:y)$ on $Y\cap\{x\ne0\}$ and by $(z:w)$ on $Y\cap\{w\neq0\}$. On the overlap the two agree since $xw=yz$ on $Y$. Each piece is given by linear forms without common zero on the open set in question, so is a morphism (as in 3.14(a)), hence $\rho$ is a morphism. For $(t:u)\in\mathbf P^1$: if $t\ne0$, $\rho\varphi(t:u)=(t^4:t^3u)=(t:u)$; if $t=0$, $\rho\varphi(0:1)=(0:1)$. And $\varphi\rho=\mathrm{id}_Y$ because each $Q\in Y$ is $\varphi(t:u)$ and $\rho(Q)=(t:u)$. So $\varphi\colon\mathbf P^1\to Y$ is an isomorphism. Finally, $S(\mathbf P^1)=k[t,u]$ for the identity embedding is a UFD, hence integrally closed, so $\mathbf P^1$ is projectively normal.

\textbf{(b)} \emph{$Y$ is normal:} $Y\cong\mathbf P^1$ by (c), and $\mathbf P^1$ is normal by 3.17(a) (its local rings are the local rings of $\mathbf A^1$, localisations of $k[T]$).

\emph{$Y$ is not projectively normal.} A polynomial $F\in k[x,y,z,w]$, homogeneous of degree $d$, lies in $I(Y)$ iff $F(\lambda t^4,\lambda t^3u,\lambda tu^3,\lambda u^4)=0$ for all $\lambda,t,u$, iff $F(t^4,t^3u,tu^3,u^4)=0$ in $k[t,u]$ ($k$ is infinite). Hence $x,y,z,w\mapsto t^4,t^3u,tu^3,u^4$ induces an isomorphism of graded rings
\[
S(Y)\cong R=k[t^4,t^3u,tu^3,u^4]\subseteq k[t,u].
\]
Consider $f=t^2u^2\in k[t,u]$. It lies in $\operatorname{Frac}R$, since $f=(t^3u)^2/t^4$. It is integral over $R$: $f^2=t^4u^4=(t^4)(u^4)\in R$, so $f$ is a root of $X^2-xw\in R[X]$. But $f\notin R$: $R$ is graded by degree in $t,u$ (it is generated by elements of degree $4$), and its degree-$4$ component is spanned by the four generators $t^4,t^3u,tu^3,u^4$, which do not include $t^2u^2$; so $t^2u^2$, which is homogeneous of degree $4$, is not in $R$. So $S(Y)$ is not integrally closed and $Y$ is not projectively normal in this embedding.\hfill$\square$

\exercise{3.19}{\textbf{Automorphisms of $\mathbf A^n$.} Let $\varphi\colon\mathbf A^n\to\mathbf A^n$ be a morphism of $\mathbf A^n$ to $\mathbf A^n$ given by $n$ polynomials $f_1,\dots,f_n$ of $n$ variables $x_1,\dots,x_n$. Let $J=\det|\partial f_i/\partial x_j|$ be the \emph{Jacobian polynomial} of $\varphi$.
\begin{enumerate}[label=(\alph*),leftmargin=*]
\item If $\varphi$ is an isomorphism (in which case we call $\varphi$ an \emph{automorphism} of $\mathbf A^n$) show that $J$ is a nonzero constant polynomial.
\item The converse of (a) is an unsolved problem, even for $n=2$. See, for example, Vitushkin [1]. (Doubly starred in the book.)
\end{enumerate}}
\pf
\textbf{(a)} Suppose $\varphi$ has an inverse morphism $\psi$. A morphism to $\mathbf A^n$ is given by $n$ regular functions on $\mathbf A^n$, which are polynomials by Theorem 3.2(a); so $\psi$ is given by polynomials $g_1,\dots,g_n$, and $\psi\circ\varphi=\mathrm{id}$ says
\[
g_i\bigl(f_1(x),\dots,f_n(x)\bigr)=x_i\quad\text{in }k[x_1,\dots,x_n]
\]
(two polynomials agreeing as functions on $\mathbf A^n$, $k$ infinite, are equal). Differentiate with respect to $x_j$ using the chain rule for formal derivatives, valid over any field:
\[
\sum_{l=1}^n\frac{\partial g_i}{\partial x_l}\bigl(f(x)\bigr)\cdot\frac{\partial f_l}{\partial x_j}(x)=\delta_{ij}.
\]
In matrix form, $(Dg)(f(x))\cdot(Df)(x)=\mathrm{Id}_n$ over $k[x]$. Taking determinants, $\det(Dg)(f(x))\cdot J(x)=1$ in $k[x]$. So $J$ is a unit of $k[x_1,\dots,x_n]$, i.e.\ a nonzero constant.

\textbf{(b)} This is the (as of the book's writing, and still) open \emph{Jacobian Conjecture} (in characteristic $0$): if the Jacobian of $\varphi$ is a nonzero constant then $\varphi$ is an automorphism. No solution can be given here. We only record two remarks. For $n=1$ in characteristic $0$ it is easy: $J=f'$ a nonzero constant forces $f$ to have degree $1$, so $\varphi$ is an affine-linear isomorphism. In characteristic $p>0$ the converse is simply false: $f(x)=x-x^p$ has $f'=1$ but is not injective, since $f(a)=0$ for all $a\in\mathbf F_p\subseteq k$.\hfill$\square$

\exercise{3.20}{Let $Y$ be a variety of dimension $\ge2$, and let $P\in Y$ be a normal point. Let $f$ be a regular function on $Y-P$.
\begin{enumerate}[label=(\alph*),leftmargin=*]
\item Show that $f$ extends to a regular function on $Y$.
\item Show this would be false for $\dim Y=1$.
\end{enumerate}
See (III, Ex.~3.5) for generalization.}
\pf
\textbf{(a)} Let $d=\dim Y\ge2$. Choose an affine open neighbourhood $V$ of $P$, and let $A=A(V)$ (after identifying $V$ with an affine variety), $K=\operatorname{Frac}A=K(Y)$. By 3.12 $\dim V=d$. Let $\fm=\fm_P$ and $B=\mathcal{O}_P=A_\fm$; by hypothesis $B$ is an integrally closed Noetherian domain with fraction field $K$. The function $f|_{V-P}$ defines an element of $K$, which is regular at every point $Q\ne P$ of $V$, i.e.\ $f\in A_{\fm_Q}$ for all $Q\in V-\{P\}$.

\emph{Claim: $f\in B$.} We use the standard theorem from commutative algebra (e.g.\ Matsumura, \emph{Commutative Ring Theory}, Thm 11.5) that a Noetherian integrally closed domain $B$ (cf.\ Hartshorne II, Prop.~6.3A) satisfies $B=\bigcap_{\operatorname{ht}\fq'=1}B_{\fq'}$ in $K$, the intersection over height-one primes. The height-one primes of $B=A_\fm$ are $\fp B$ with $\fp\subseteq\fm$ a height-one prime of $A$, and $B_{\fp B}=A_\fp$. So it suffices to show $f\in A_\fp$ for every height-one prime $\fp\subseteq\fm$ of $A$. Now $\fp$ is prime, so $\fp=I(Z)$ with $Z=Z(\fp)\subseteq V$ an irreducible closed subset, and by Theorem 1.8A, $\dim Z=\dim A/\fp=d-\operatorname{ht}\fp=d-1\ge1$. An irreducible set with at most one point has dimension $0$, so $Z$ is infinite and contains a point $Q\ne P$. Then $\fp=I(Z)\subseteq\fm_Q$, so $A_\fp$ is a localisation of $A_{\fm_Q}$, and $f\in A_{\fm_Q}\subseteq A_\fp$. This proves the claim.

Thus there is an open neighbourhood $U\subseteq V$ of $P$ and a regular function $g$ on $U$ with $g=f$ in $K$, i.e.\ $g=f$ on the dense open subset $U-\{P\}$ of $U$ (regular functions on an open subset of an irreducible variety that agree on a nonempty open subset agree). Define $F=g$ on $U$ and $F=f$ on $Y-P$; these agree on the overlap $U-P$, and regularity is local, so $F$ is a regular function on $Y=U\cup(Y-P)$ extending $f$.

\textbf{(b)} Take $Y=\mathbf A^1$, which is normal ($k[T]$ is integrally closed) and of dimension $1$, $P=0$, and $f=1/T$, regular on $\mathbf A^1-\{0\}$. A regular function on $\mathbf A^1$ is a polynomial $g\in k[T]$ (Thm 3.2(a)); if $g$ extended $f$, then $Tg(T)=1$ for all $T\neq0$, hence as polynomials (since $k$ is infinite), but evaluating at $T=0$ gives $0=1$, a contradiction. So the extension property fails in dimension $1$.\hfill$\square$

\exercise{3.21}{\textbf{Group Varieties.} A \emph{group variety} consists of a variety $Y$ together with a morphism $\mu\colon Y\times Y\to Y$, such that the set of points of $Y$ with the operation given by $\mu$ is a group, and such that the inverse map $y\to y^{-1}$ is also a morphism of $Y\to Y$.
\begin{enumerate}[label=(\alph*),leftmargin=*]
\item The \emph{additive group} $\mathbf G_a$ is given by the variety $\mathbf A^1$ and the morphism $\mu\colon\mathbf A^2\to\mathbf A^1$ defined by $\mu(a,b)=a+b$. Show it is a group variety.
\item The \emph{multiplicative group} $\mathbf G_m$ is given by the variety $\mathbf A^1-\{0\}$ and the morphism $\mu(a,b)=ab$. Show it is a group variety.
\item If $G$ is a group variety, and $X$ is any variety, show that the set $\operatorname{Hom}(X,G)$ has a natural group structure.
\item For any variety $X$, show that $\operatorname{Hom}(X,\mathbf G_a)$ is isomorphic to $\mathcal{O}(X)$ as a group under addition.
\item For any variety $X$, show that $\operatorname{Hom}(X,\mathbf G_m)$ is isomorphic to the group of units in $\mathcal{O}(X)$, under multiplication.
\end{enumerate}}
\pf
By Ex.~3.15(c), 3.16(c), the product $Y\times Y$ in the category of varieties is unique up to unique isomorphism; for $Y=\mathbf A^1$, $\mathbf A^1\times\mathbf A^1=\mathbf A^2$ with the usual projections (Ex.~3.15 with $X=Y=\mathbf A^1$). We use this identification.

\textbf{(a)} $\mu(a,b)=a+b$ is a polynomial map $\mathbf A^2\to\mathbf A^1$, hence a morphism (Lemma 3.6). The set $k$ with $a+b$ is a group, with identity $0$ and inverse $a\mapsto-a$, which is a polynomial map $\mathbf A^1\to\mathbf A^1$, hence a morphism.

\textbf{(b)} Let $Y=\mathbf A^1-\{0\}$ (an open subvariety of $\mathbf A^1$). \emph{The product $Y\times Y$ is the open subset $U=\{(a,b)\in\mathbf A^2:ab\ne0\}$ of $\mathbf A^2$ with the usual projections.} Indeed the projections $U\to Y$ are restrictions of morphisms, hence morphisms; and given morphisms $\varphi,\rho\colon Z\to Y$, the map $(\varphi,\rho)\colon Z\to\mathbf A^2$ is a morphism by 3.15(c) with image in the open set $U$, hence is a morphism $Z\to U$ (Ex.~3.10; or directly, since $U$ carries the induced structure). Uniqueness is clear as maps of sets. So $\mu\colon U\to Y$, $(a,b)\mapsto ab$, is the restriction of the polynomial map $\mathbf A^2\to\mathbf A^1$ with image in $Y$, hence a morphism $U\to Y$ (Ex.~3.10). The set $k^*$ with multiplication is a group, and the inverse $a\mapsto1/a$ is a morphism $Y\to Y$ because $1/x$ is a regular function on $Y$ (a quotient of polynomials with nonvanishing denominator) with values in $Y$. So $\mathbf G_m$ is a group variety.

\textbf{(c)} Let $G$ have multiplication $\mu$ and inverse $\iota$ (a morphism). For $\varphi,\rho\in\operatorname{Hom}(X,G)$ let $(\varphi,\rho)\colon X\to G\times G$ be the morphism of 3.15(c)/3.16(c), $x\mapsto(\varphi(x),\rho(x))$, and define
\[
\varphi\cdot\rho=\mu\circ(\varphi,\rho),\qquad \varphi^{-1}=\iota\circ\varphi,\qquad e_X=\text{the constant map }X\to\{e\}\subseteq G.
\]
These are morphisms (the constant map is a morphism: it is continuous, and pulls a regular function $g$ near $e$ back to the constant $g(e)$). As maps of sets, $(\varphi\cdot\rho)(x)=\varphi(x)\rho(x)$, $\varphi^{-1}(x)=\varphi(x)^{-1}$, $e_X(x)=e$; and a morphism is determined by its underlying map of sets. So associativity, the identity and inverse laws in $\operatorname{Hom}(X,G)$ hold because they hold pointwise in the group $G$: $\operatorname{Hom}(X,G)$ is a group. This structure is natural: for a morphism $h\colon X'\to X$ the map $h^*\colon\operatorname{Hom}(X,G)\to\operatorname{Hom}(X',G)$, $\varphi\mapsto\varphi\circ h$, is a group homomorphism (since $(\varphi\rho)\circ h$ and $(\varphi\circ h)(\rho\circ h)$ agree pointwise), and a homomorphism of group varieties $G\to G'$ induces a group homomorphism $\operatorname{Hom}(X,G)\to\operatorname{Hom}(X,G')$.

\textbf{(d)} A morphism $X\to\mathbf A^1$ is the same as a regular function on $X$: given a morphism $\varphi$, $\varphi=t\circ\varphi$ is the pullback of the coordinate function $t$ (regular on $\mathbf A^1$), hence regular; conversely a regular function $f$ is a morphism $X\to\mathbf A^1$ by Lemma 3.6. So $\operatorname{Hom}(X,\mathbf G_a)=\mathcal{O}(X)$ as sets, via $\varphi\mapsto\varphi$ (viewed as function). The group law is $(\varphi+\rho)(x)=\mu(\varphi(x),\rho(x))=\varphi(x)+\rho(x)$, which is the sum in the ring $\mathcal{O}(X)$. Hence $\operatorname{Hom}(X,\mathbf G_a)\cong(\mathcal{O}(X),+)$.

\textbf{(e)} A morphism $\varphi\colon X\to\mathbf G_m=\mathbf A^1-\{0\}$ composed with the inclusion into $\mathbf A^1$ gives a regular function $f$ on $X$ that vanishes nowhere. The function $1/t$ is regular on $\mathbf A^1-\{0\}$, so $1/f=\varphi^*(1/t)$ is regular on $X$; thus $f$ is a unit of $\mathcal{O}(X)$. Conversely if $u\in\mathcal{O}(X)^\times$, then $u\colon X\to\mathbf A^1$ is a morphism whose values are nonzero ($u\cdot u^{-1}=1$), so it is a morphism $X\to\mathbf A^1-\{0\}$ by Ex.~3.10. These two operations are inverse, so $\operatorname{Hom}(X,\mathbf G_m)\leftrightarrow\mathcal{O}(X)^\times$ is a bijection. The group law is $(\varphi\cdot\rho)(x)=\mu(\varphi(x),\rho(x))=\varphi(x)\rho(x)$, the product in $\mathcal{O}(X)$. Hence $\operatorname{Hom}(X,\mathbf G_m)\cong\mathcal{O}(X)^\times$ as groups.\hfill$\square$

\section*{4. Rational maps}

Conventions as in Hartshorne: a rational map $\varphi:X\to Y$ is an equivalence class of pairs $\langle U,\varphi_U\rangle$ ($U$ nonempty open in $X$, $\varphi_U:U\to Y$ a morphism), where $\langle U,\varphi_U\rangle\sim\langle V,\varphi_V\rangle$ if $\varphi_U,\varphi_V$ agree on $U\cap V$; a rational function is a rational map to $\mathbf A^1$ (equivalently a class of pairs $\langle U,f\rangle$, $f$ regular on $U$). I use Lemma~4.1 (two morphisms $X\to Y$ of varieties agreeing on a nonempty open set are equal), Thm.~4.4 and Cor.~4.5 (birational equivalence $\Leftrightarrow$ isomorphic function fields $\Leftrightarrow$ isomorphic nonempty open subsets), and that regular functions are continuous maps to $\mathbf A^1$ and that being a regular function or a morphism is a local property.

\exercise{4.1}{If $f$ and $g$ are regular functions on open subsets $U$ and $V$ of a variety $X$, and if $f=g$ on $U\cap V$, show that the function which is $f$ on $U$ and $g$ on $V$ is a regular function on $U\cup V$. Conclude that if $f$ is a rational function on $X$, then there is a largest open subset $U$ of $X$ on which $f$ is represented by a regular function. We say that $f$ is \emph{defined} at the points of $U$.}
\pf
Let $h:U\cup V\to k$ be $f$ on $U$ and $g$ on $V$ (well defined because $f=g$ on $U\cap V$). Let $P\in U\cup V$, say $P\in U$. Since $f$ is regular at $P$, there are an open $N$ with $P\in N\subseteq U$ and polynomials $G,H$ (homogeneous of the same degree in the projective case) with $H$ nowhere zero on $N$ and $f=G/H$ on $N$. $N$ is open in $U\cup V$ and $h=G/H$ on $N$, so $h$ is regular at $P$. The case $P\in V$ is identical. So $h$ is regular on $U\cup V$.

Now let $\varphi$ be a rational function on $X$, i.e.\ an equivalence class of pairs. Let $\Omega$ be the union of all open sets $U$ such that $\langle U,f_U\rangle$ is a representative of $\varphi$. If $\langle U,f_U\rangle,\langle V,f_V\rangle$ are two representatives and $\langle W,f_W\rangle$ is any fixed one, then $f_U=f_W=f_V$ on $U\cap W\cap V$. The set $U\cap V$ is open and (if nonempty) irreducible, and $U\cap V\cap W$ is a nonempty open subset of it, hence dense in $U\cap V$ (Ex.~1.6; $X$ is irreducible so any two nonempty opens meet). The set $\{f_U=f_V\}\subseteq U\cap V$ is closed in $U\cap V$ (it is the zero set of the regular, hence continuous, function $f_U-f_V$) and contains the dense set $U\cap V\cap W$, so $f_U=f_V$ on $U\cap V$. Therefore the $f_U$ glue to a function $f_\Omega$ on $\Omega$, regular at each point by the same local argument as above. Then $\langle\Omega,f_\Omega\rangle$ represents $\varphi$ (it agrees with $f_W$ on $\Omega\cap W$ by construction) and $\Omega$ contains every open set on which $\varphi$ is represented by a regular function, so $\Omega$ is the largest such open set.\hfill$\square$

\exercise{4.2}{Same problem for rational maps. If $\varphi$ is a rational map of $X$ to $Y$, show there is a largest open set on which $\varphi$ is represented by a morphism. We say the rational map is \emph{defined} at the points of this open set.}
\pf
Let $\varphi$ be represented by $\langle U,\varphi_U\rangle$ and $\langle V,\varphi_V\rangle$; both represent $\varphi$, so they agree on $U\cap V$ by definition of the equivalence. Define $\psi$ on $U\cup V$ as $\varphi_U$ on $U$ and $\varphi_V$ on $V$. Then $\psi$ is a morphism: it is continuous, since for $Z\subseteq Y$ closed, $\psi^{-1}(Z)\cap U=\varphi_U^{-1}(Z)$ is closed in $U$ and $\psi^{-1}(Z)\cap V=\varphi_V^{-1}(Z)$ is closed in $V$, so the complement of $\psi^{-1}(Z)$ is open in $U$ and in $V$, hence open in $U\cup V$; and for $W\subseteq Y$ open and $g$ regular on $W$, $g\circ\psi$ is regular on $\psi^{-1}(W)$ because it is regular on $\psi^{-1}(W)\cap U$ and on $\psi^{-1}(W)\cap V$ (regularity is local).

For an arbitrary family: let $\Omega$ be the union of all open sets $U$ such that some representative $\langle U,\varphi_U\rangle$ of $\varphi$ exists. Any two representatives agree on the overlap of their domains, so they glue (the preceding argument works verbatim for an arbitrary open cover) to a morphism $\varphi_\Omega:\Omega\to Y$ with $\langle\Omega,\varphi_\Omega\rangle\sim\varphi$. Clearly $\Omega$ is the largest open set on which $\varphi$ is represented by a morphism.\hfill$\square$

\exercise{4.3}{
\begin{enumerate}[label=(\alph*),leftmargin=*]
\item Let $f$ be the rational function on $\mathbf P^2$ given by $f=x_1/x_0$. Find the set of points where $f$ is defined and describe the corresponding regular function.
\item Now think of this function as a rational map from $\mathbf P^2$ to $\mathbf A^1$. Embed $\mathbf A^1$ in $\mathbf P^1$, and let $\varphi:\mathbf P^2\to\mathbf P^1$ be the resulting rational map. Find the set of points where $\varphi$ is defined, and describe the corresponding morphism.
\end{enumerate}}
\pf
Let $U_i=\{x_i\ne0\}\subseteq\mathbf P^2$.

(a) \emph{Answer: $f$ is defined exactly on $U_0=\{x_0\neq0\}$, where it is the regular function $(a_0:a_1:a_2)\mapsto a_1/a_0$.} On $U_0$, $f=x_1/x_0$ is a quotient of homogeneous polynomials of the same degree with denominator nowhere zero, so it is regular there. Conversely suppose $f$ is defined at some $P$ with $x_0(P)=0$: there are an open $V\ni P$ and a regular $g$ on $V$ with $g=f$ on $V\cap U_0$. Since $g$ is regular at $P$, there are an open $N\ni P$, $N\subseteq V$, and homogeneous $G,H$ of the same degree with $H$ nowhere zero on $N$ and $g=G/H$ on $N$. On $N\cap U_0$ (nonempty open) we get $G/H=x_1/x_0$, so $x_0G-x_1H$ vanishes on $N\cap U_0$, which is dense in $\mathbf P^2$ (Ex.~1.6); as $Z(x_0G-x_1H)$ is closed it is all of $\mathbf P^2$, and since $k$ is infinite the homogeneous polynomial $x_0G-x_1H$ is zero. Thus $x_0G=x_1H$ in the UFD $k[x_0,x_1,x_2]$, and as the prime $x_0$ does not divide $x_1$, $x_0\mid H$. Then $H(P)=0$, contradicting $H\ne0$ on $N\ni P$. So the domain of definition is exactly $U_0$.

(b) Embed $\mathbf A^1\hookrightarrow\mathbf P^1$ by $t\mapsto(1:t)$ (with homogeneous coordinates $(y_0:y_1)$ on $\mathbf P^1$, so $\mathbf A^1=V_0:=\{y_0\ne0\}$). Then $\varphi(x_0:x_1:x_2)=(1:x_1/x_0)=(x_0:x_1)$ on $U_0$. \emph{Answer: $\varphi$ is defined exactly on $\mathbf P^2\setminus\{Q\}$, $Q=(0:0:1)$, and there it is the morphism $(a_0:a_1:a_2)\mapsto(a_0:a_1)$} (projection from $Q$ to the line $x_2=0$).

\emph{Defined away from $Q$.} The map $\pi(a_0:a_1:a_2)=(a_0:a_1)$ is well defined on $\mathbf P^2\setminus\{Q\}=U_0\cup U_1$. With $V_1=\{y_1\ne0\}\subseteq\mathbf P^1$ we have $\pi^{-1}(V_0)=U_0$, $\pi^{-1}(V_1)=U_1$, and the coordinate functions $y_1/y_0\circ\pi=x_1/x_0$ on $U_0$, $y_0/y_1\circ\pi=x_0/x_1$ on $U_1$ are regular. So $\pi$ is a morphism, and it agrees with $\varphi$ on $U_0$.

\emph{Not defined at $Q$.} Suppose $\tilde\pi:U'\to\mathbf P^1$ is a morphism on an open $U'\ni Q$ agreeing with $\varphi$ (i.e.\ with $(x_0:x_1)$) on $U'\cap U_0$. Then $\tilde\pi(Q)\in V_0$ or $V_1$. Let $N=\tilde\pi^{-1}(V_j)\ni Q$, open.
\begin{itemize}
\item If $\tilde\pi(Q)\in V_0$: $g=(y_1/y_0)\circ\tilde\pi$ is regular on $N$ and equals $x_1/x_0$ on $N\cap U_0$. As in (a), near $Q$ write $g=G/H$ with $H(Q)\ne0$; then $x_0G=x_1H$ on a dense open set, hence as polynomials, so $x_0\mid H$ and $H(Q)=0$ (as $x_0(Q)=0$), a contradiction.
\item If $\tilde\pi(Q)\in V_1$: $g=(y_0/y_1)\circ\tilde\pi$ is regular on $N$ and equals $x_0/x_1$ on $N\cap U_0\cap U_1$ (nonempty open). Near $Q$ write $g=G/H$, $H(Q)\ne0$; then $x_1G=x_0H$ as polynomials, so $x_1\mid H$, and $H(Q)=0$ since $x_1(Q)=0$, a contradiction.
\end{itemize}
Hence $\varphi$ is not defined at $Q$, and the domain of definition is $\mathbf P^2\setminus\{Q\}$.\hfill$\square$

\exercise{4.4}{A variety $Y$ is \emph{rational} if it is birationally equivalent to $\mathbf P^n$ for some $n$ (or, equivalently by (4.5), if $K(Y)$ is a pure transcendental extension of $k$).
\begin{enumerate}[label=(\alph*),leftmargin=*]
\item Any conic in $\mathbf P^2$ is a rational curve.
\item The cuspidal cubic $y^2=x^3$ is a rational curve.
\item Let $Y$ be the nodal cubic curve $y^2z=x^2(x+z)$ in $\mathbf P^2$. Show that the projection $\varphi$ from the point $P=(0,0,1)$ to the line $z=0$ (Ex.~3.14) induces a birational map from $Y$ to $\mathbf P^1$. Thus $Y$ is a rational curve.
\end{enumerate}}
\pf
(a) A conic is $Y=Z(F)\subseteq\mathbf P^2$ with $F$ an irreducible quadratic form, so $Y$ is a projective variety of dimension $1$. It is nonempty: restrict $F$ to a line; this is a binary form with a zero in $\mathbf P^1$ ($k$ algebraically closed). Choose a point $P\in Y$ and (after a projective linear change of coordinates) $P=(0:0:1)$. Since $F(P)=0$, $F$ has no $x_2^2$ term, so $F=x_2L(x_0,x_1)+Q(x_0,x_1)$ with $L$ linear, $Q$ quadratic. If $L=0$ then $F=Q(x_0,x_1)$ would be a product of linear forms, reducible; so $L\ne0$.

Let $\pi(x_0:x_1:x_2)=(x_0:x_1)$, a morphism on $Y\setminus\{P\}$ (Ex.~3.14), and define
\[
\psi:\mathbf P^1\supseteq W\to Y,\qquad \psi(s:t)=\bigl(sL(s,t):tL(s,t):-Q(s,t)\bigr),\qquad W=\{L(s,t)\ne0\}.
\]
$W$ is the complement of one point of $\mathbf P^1$, so nonempty open; the three coordinates are not all $0$ on $W$ and $\psi$ is a morphism. Moreover $\psi(W)\subseteq Y$: using homogeneity, $F(sL,tL,-Q)=(-Q)\cdot L(sL,tL)+Q(sL,tL)=-QL^2+L^2Q=0$. We have $\pi\circ\psi=\mathrm{id}_W$, since $\pi\psi(s:t)=(sL:tL)=(s:t)$.

Let $Y'=\{x\in Y\setminus\{P\}: L(x_0,x_1)\neq0\}$, open in $Y$. It is nonempty: otherwise $Y\subseteq Z(L)$, so $L\in I(Y)=(F)$ (projective Nullstellensatz, $F$ irreducible), impossible for degree reasons. For $x\in Y'$ we have $x_2L+Q=0$, so $\psi\pi(x)=(x_0L:x_1L:-Q)=(x_0L:x_1L:x_2L)=x$ (where $L=L(x_0,x_1)\ne0$). Thus $\psi\circ\pi=\mathrm{id}$ on $Y'$ and $\pi\circ\psi=\mathrm{id}$ on $W$. In particular $\pi(\psi(W))=W$ and $\psi(\pi(Y'))=Y'$, so the images of $\pi$ and $\psi$ contain the dense open sets $W$ resp.\ $Y'$ (Ex.~1.6): both are dominant rational maps, and they are inverse to each other. So $Y$ is birational to $\mathbf P^1$, i.e.\ rational.

(b) Let $Y=Z(y^2-x^3)\subseteq\mathbf A^2$ (its projective closure $y^2z=x^3$ contains it as an open subset, so has the same function field and the same conclusion). $y^2-x^3$ is irreducible: it is monic of degree $2$ in $y$, so a factorization would give a root $y=r(x)\in k[x]$ with $r^2=x^3$, impossible by degrees (odd vs even). Define
\[
\varphi:\mathbf A^1\to Y,\ t\mapsto(t^2,t^3),\qquad \psi:Y\setminus\{O\}\to\mathbf A^1,\ (x,y)\mapsto y/x,
\]
$O=(0,0)$. On $Y\setminus\{O\}$ we have $x\ne0$ (if $x=0$ then $y^2=0$), so $\psi$ is regular. For $t\neq0$: $\psi\varphi(t)=t^3/t^2=t$; for $(x,y)\in Y\setminus\{O\}$: $\varphi\psi(x,y)=(y^2/x^2,\,y^3/x^3)=(x,y)$ because $y^2=x^3$. Thus $\varphi$ restricts to an isomorphism $\mathbf A^1\setminus\{0\}\to Y\setminus\{O\}$ (inverse $\psi$). Hence $Y$ is birational to $\mathbf A^1$, which is an open subset of $\mathbf P^1$, so $Y$ is birational to $\mathbf P^1$. (Equivalently $K(Y)=k(x,y)=k(t)$ with $t=y/x$.)

(c) Let $F=y^2z-x^3-x^2z$ (so $Y=Z(F)$; $P=(0:0:1)\in Y$). $F$ is irreducible: as a polynomial in $y$ over $k[x,z]$ it is $zy^2-x^2(x+z)$, with content $\gcd(z,x^2(x+z))=1$; if it were reducible in $k(x,z)[y]$ it would have a root $y=r\in k(x,z)$, i.e.\ $(zr)^2=z\,x^2(x+z)$, so $z(x+z)$ would be a square in $k(x,z)$, hence in the UFD $k[x,z]$, impossible since $z$ and $x+z$ are non-associate primes. By Gauss's lemma $F$ is irreducible, and $Y$ is a variety.

The points of $Y$ with $x=y=0$: then $F=0$ holds automatically, giving only $P=(0:0:1)$. So $\pi(x:y:z)=(x:y)$ is a morphism on $Y\setminus\{P\}$ (Ex.~3.14). On $Y$, $F=0$ reads $z(y^2-x^2)=x^3$. For $(x:y:z)\in Y\setminus\{P\}$ we have $y^2-x^2\ne0$: if $y^2=x^2$ then $x^3=0$, so $x=y=0$ and the point is $P$. So $\pi$ maps $Y\setminus\{P\}$ into $W=\{(s:t)\in\mathbf P^1:t^2\ne s^2\}$, the complement of the two points $(1:\pm1)$. Define, with $D(s,t)=t^2-s^2$,
\[
\psi:W\to Y\setminus\{P\},\qquad \psi(s:t)=\bigl(sD:tD:s^3\bigr).
\]
The coordinates are not all zero ($D\ne0$ and $(s,t)\ne0$), so $\psi$ is a morphism. It lands in $Y$: $F(sD,tD,s^3)=t^2D^2s^3-s^3D^3-s^2D^2s^3=s^3D^2(t^2-D-s^2)=0$. And $\psi(s:t)\neq P$ since $sD$ or $tD\ne0$. Now $\pi\psi(s:t)=(sD:tD)=(s:t)$, and for $(x:y:z)\in Y\setminus\{P\}$, with $D=y^2-x^2\ne0$, we get $\psi\pi(x:y:z)=(xD:yD:x^3)=(xD:yD:zD)=(x:y:z)$ because $x^3=z(y^2-x^2)$. Hence $\pi:Y\setminus\{P\}\to W$ and $\psi$ are mutually inverse isomorphisms. So $Y$ is birational to $\mathbf P^1$ via the projection from $P$, and $Y$ is rational.\hfill$\square$

\exercise{4.5}{Show that the quadric surface $Q:xy=zw$ in $\mathbf P^3$ is birational to $\mathbf P^2$, but not isomorphic to $\mathbf P^2$ (cf.\ Ex.~2.15).}
\pf
\emph{$Q$ is a variety.} $xy-zw$ is irreducible: it has degree $1$ in $x$, with coefficients $y$ and $-zw$ which are coprime, so it is primitive of degree $1$ in $x$ over $k[y,z,w]$ and therefore irreducible. So $Q$ is a projective variety of dimension $2$.

\emph{Birational.} Let $U=Q\cap\{w\ne0\}$, a nonempty open subset of $Q$. Define $\alpha:\mathbf A^2\to U$, $\alpha(a,b)=(a:b:ab:1)$ (it lies in $Q$ as $ab=ab\cdot1$) and $\beta:U\to\mathbf A^2$, $\beta(x:y:z:w)=(x/w,\,y/w)$. Both are morphisms, $\beta\alpha(a,b)=(a,b)$, and $\alpha\beta(x:y:z:w)=(x/w:y/w:xy/w^2:1)=(x:y:z:w)$ since $xy=zw$ gives $xy/w^2=z/w$. So $U\cong\mathbf A^2$, which is an open subset of $\mathbf P^2$. Hence $Q$ and $\mathbf P^2$ are birational (Cor.~4.5).

\emph{Not isomorphic.} $Q$ contains the lines $L_1=\{x=z=0\}$ and $L_2=\{y=w=0\}$ (both lie in $Q$, since $xy=0=zw$ on them), and $L_1\cap L_2=\emptyset$ because $x=y=z=w=0$ is impossible. Suppose $\theta:Q\to\mathbf P^2$ were an isomorphism. Then $\theta(L_1),\theta(L_2)$ would be closed irreducible subsets of $\mathbf P^2$ of dimension $1$ (curves, since $\theta$ is a homeomorphism and $\dim L_i=1$), with $\theta(L_1)\cap\theta(L_2)=\theta(L_1\cap L_2)=\emptyset$. But any two curves in $\mathbf P^2$ meet (Ex.~3.7(a); more generally Thm.~7.2, dimension of intersections in projective space). Contradiction. So $Q\not\cong\mathbf P^2$.\hfill$\square$

\exercise{4.6}{\emph{Plane Cremona Transformations.} A birational map of $\mathbf P^2$ into itself is called a \emph{plane Cremona transformation}. We give an example, called a \emph{quadratic transformation}. It is the rational map $\varphi:\mathbf P^2\to\mathbf P^2$ given by $(a_0,a_1,a_2)\to(a_1a_2,a_0a_2,a_0a_1)$ when no two of $a_0,a_1,a_2$ are $0$.
\begin{enumerate}[label=(\alph*),leftmargin=*]
\item Show that $\varphi$ is birational, and is its own inverse.
\item Find open sets $U,V\subseteq\mathbf P^2$ such that $\varphi:U\to V$ is an isomorphism.
\item Find the open sets where $\varphi$ and $\varphi^{-1}$ are defined, and describe the corresponding morphisms. See also (V,~4.2.3).
\end{enumerate}}
\pf
Let $e_0=(1:0:0)$, $e_1=(0:1:0)$, $e_2=(0:0:1)$, $U_1'=\mathbf P^2\setminus\{e_0,e_1,e_2\}$ and $T=\{a_0a_1a_2\neq0\}$. The three polynomials $a_1a_2,a_0a_2,a_0a_1$ vanish simultaneously exactly when at least two of the $a_i$ vanish, i.e.\ at $e_0,e_1,e_2$. So the formula defines a morphism $\varphi:U_1'\to\mathbf P^2$ (this is the meaning of ``when no two of the $a_i$ are $0$''), which is a representative of the rational map $\varphi$.

(a) For $a\in T$, put $b=\varphi(a)=(a_1a_2,a_0a_2,a_0a_1)$. All $b_i\ne0$, so $\varphi(a)\in T$, and
\[
\varphi(b)=(b_1b_2,b_0b_2,b_0b_1)=(a_0^2a_1a_2,\ a_0a_1^2a_2,\ a_0a_1a_2^2)=a_0a_1a_2\cdot(a_0,a_1,a_2),
\]
which is the point $a$. So $\varphi\circ\varphi=\mathrm{id}$ on the nonempty open set $T$ and $\varphi(T)\subseteq T$. In particular $\varphi$ is dominant (its image contains $\varphi(\varphi(T))=T$, dense), so $\varphi\circ\varphi$ is defined as a composition of rational maps and equals $\mathrm{id}$. Hence $\varphi$ is birational and $\varphi^{-1}=\varphi$.

(b) Take $U=V=T=\{a_0a_1a_2\ne0\}$. $\varphi|_T:T\to T$ is a morphism with $\varphi\circ\varphi=\mathrm{id}_T$, hence an isomorphism (it is its own inverse).

(c) \emph{Answer: $\varphi$ and $\varphi^{-1}=\varphi$ are defined exactly on $U_1'=\mathbf P^2\setminus\{e_0,e_1,e_2\}$, where they are given by the morphism $(a_0:a_1:a_2)\mapsto(a_1a_2:a_0a_2:a_0a_1)$.} It remains to show that $\varphi$ is not defined at any $e_i$. Since $\varphi$ commutes with permutations of the coordinates it suffices to treat $e_0$. Suppose $\tilde\varphi:W\to\mathbf P^2$ is a morphism on an open $W\ni e_0$ agreeing with $\varphi$ on $W\cap T$. Fix $m\in\{0,1\}$ and the morphism $\iota_m:\mathbf A^1\to\mathbf P^2$, $s\mapsto(1:s:ms)$ (a line through $e_0$, $\iota_m(0)=e_0$). $N_m=\iota_m^{-1}(W)$ is an open subset of $\mathbf A^1$ containing $0$, and for $s\in N_m\setminus\{0\}$, $\iota_m(s)\in U_1'$ (it has $a_0=1,a_1=s\ne0$), where $\tilde\varphi$ agrees with the formula for $\varphi$ (two morphisms on the open set $W\cap U_1'$ agreeing on the dense open $W\cap T$ agree by Lemma~4.1):
\[
\tilde\varphi(\iota_m(s))=(s\cdot ms:\ 1\cdot ms:\ 1\cdot s)=(ms^2:ms:s)=(ms:m:1).
\]
So $\tilde\varphi\circ\iota_m$ and $g_m:\mathbf A^1\to\mathbf P^2$, $g_m(s)=(ms:m:1)$, are morphisms $N_m\to\mathbf P^2$ agreeing on $N_m\setminus\{0\}$, which is dense open; by Lemma~4.1 they agree on $N_m$, in particular $\tilde\varphi(e_0)=g_m(0)=(0:m:1)$. For $m=0$ this is $(0:0:1)$ and for $m=1$ it is $(0:1:1)$, a contradiction. Hence $\varphi$ is not defined at $e_0$, and likewise at $e_1,e_2$.

Geometrically: $\varphi|_T$ is an automorphism of $T$; the line $\{a_i=0\}$ minus its two vertices $e_j,e_k$ is contracted by $\varphi$ to the opposite vertex $e_i$ (e.g.\ $\varphi(0:a_1:a_2)=(a_1a_2:0:0)=e_0$ for $a_1a_2\ne0$); and the three vertices are the indeterminacy points (blown up to the opposite lines).\hfill$\square$

\exercise{4.7}{Let $X$ and $Y$ be two varieties. Suppose there are points $P\in X$ and $Q\in Y$ such that the local rings $\mathcal O_{P,X}$ and $\mathcal O_{Q,Y}$ are isomorphic as $k$-algebras. Then show that there are open sets $P\in U\subseteq X$ and $Q\in V\subseteq Y$ and an isomorphism of $U$ to $V$ which sends $P$ to $Q$.}
\pf
By Prop.~4.3 we may replace $X,Y$ by affine open neighbourhoods of $P,Q$ (local rings and function fields are unchanged); so $X\subseteq\mathbf A^m$, $Y\subseteq\mathbf A^N$ are affine varieties with coordinate functions $x_1,\dots,x_m$, resp.\ $y_1,\dots,y_N$. Let $\theta:\mathcal O_P\to\mathcal O_Q$ be a $k$-algebra isomorphism. Recall $\mathcal O_P\subseteq K(X)$, $\mathcal O_Q\subseteq K(Y)$, and $\mathcal O_P$ has fraction field $K(X)$; so $\theta$ extends to an isomorphism $K(X)\to K(Y)$, still called $\theta$.

\emph{Construction of $f$.} $\theta(x_i)\in\mathcal O_Q$, so there are $s,b_1,\dots,b_m\in A(Y)$ with $s(Q)\ne0$ and $\theta(x_i)=b_i/s$. On the open neighbourhood $V_0=\{s\ne0\}$ of $Q$ in $Y$ define $f:V_0\to\mathbf A^m$, $f(y)=(b_1(y)/s(y),\dots,b_m(y)/s(y))$. This is a morphism. For $h\in I(X)$, the regular function $h\circ f=h(b_1/s,\dots,b_m/s)$ on $V_0$ equals $h(\theta(x_1),\dots,\theta(x_m))=\theta(h(x_1,\dots,x_m))=\theta(0)=0$ in $K(Y)$; a regular function on an open subset of the variety $Y$ which is $0$ in $K(Y)$ is the zero function. So $f(V_0)\subseteq X$, and $f:V_0\to X$ is a morphism. Next, $f(Q)=P$: the composite $\mathrm{ev}_Q\circ\theta:\mathcal O_P\to k$ is a surjective $k$-algebra homomorphism, so its kernel is a maximal ideal of the local ring $\mathcal O_P$, i.e.\ $\mathfrak m_P$; and $\mathrm{ev}_P$ has the same kernel and is also $k$-linear and onto $k$; two such homomorphisms coincide (both are the identity on $k$ and zero on $\mathfrak m_P$, and $\mathcal O_P=k\oplus\mathfrak m_P$). So $\mathrm{ev}_Q\circ\theta=\mathrm{ev}_P$, and $f(Q)_i=\mathrm{ev}_Q(\theta(x_i))=x_i(P)$, i.e.\ $f(Q)=P$.

Symmetrically $\theta^{-1}$ gives a morphism $g:W_0\to Y$ on an open $W_0=\{t\ne0\}\ni P$ of $X$, $g(y)$ having coordinates $d_j/t$ where $\theta^{-1}(y_j)=d_j/t$, with $g(P)=Q$.

\emph{$f,g$ are inverse near $Q,P$.} Let $V_1=f^{-1}(W_0)$, an open neighbourhood of $Q$. For $j=1,\dots,N$, the regular function $y_j\circ g\circ f$ on $V_1$ is $f^*(d_j/t)=d_j(f)/t(f)$, and $d_j(f)=d_j(\theta(x_1),\dots,\theta(x_m))=\theta(d_j)$, $t(f)=\theta(t)$ ($\theta$ is a $k$-algebra homomorphism); so $y_j\circ g\circ f=\theta(d_j/t)=\theta(\theta^{-1}(y_j))=y_j$ in $K(Y)$, hence as functions on $V_1$. Thus $g(f(y))=y$ for all $y\in V_1$. Likewise $f(g(x))=x$ for all $x\in W_1$, where $W_1=g^{-1}(V_0)$ is an open neighbourhood of $P$.

\emph{Conclusion.} Put $U=W_1\ni P$ and $V=f^{-1}(W_1)\cap V_1\ni Q$ (note $f(Q)=P\in W_1$). If $x\in U$, then $g(x)\in V_0$ and $f(g(x))=x\in W_1\subseteq W_0$, so $g(x)\in f^{-1}(W_0)=V_1$ and $g(x)\in f^{-1}(W_1)$; thus $g(U)\subseteq V$. If $y\in V$, then $f(y)\in W_1$ and $g(f(y))=y$ because $y\in V_1$. Hence $f|_V:V\to U$ and $g|_U:U\to V$ are mutually inverse morphisms, so $f$ is an isomorphism $V\to U$ with $f(Q)=P$, and its inverse is an isomorphism $U\to V$ sending $P$ to $Q$.\hfill$\square$

\exercise{4.8}{
\begin{enumerate}[label=(\alph*),leftmargin=*]
\item Show that any variety of positive dimension over $k$ has the same cardinality as $k$. [\emph{Hints:} Do $\mathbf A^n$ and $\mathbf P^n$ first. Then for any $X$, use induction on the dimension $n$. Use (4.9) to make $X$ birational to a hypersurface $H\subseteq\mathbf P^{n+1}$. Use (Ex.~3.7) to show that the projection of $H$ to $\mathbf P^n$ from a point not on $H$ is finite-to-one and surjective.]
\item Deduce that any two curves over $k$ are homeomorphic (cf.\ Ex.~3.1).
\end{enumerate}}
\pf
Write $\kappa=|k|$; $k$ is infinite, so $\kappa\cdot\kappa=\kappa$, $|k^n|=\kappa$ for $n\ge1$, and a union of at most $\kappa$ sets each of cardinality $\le\kappa$ has cardinality $\le\kappa$. A (quasi-projective) variety is a subset of some $\mathbf P^N$, so $|X|\le|\mathbf P^N|\le|k^{N+1}|=\kappa$ always. Thus it suffices to prove $|X|\ge\kappa$. (I use here a route different from the hint, replacing the induction by Noether normalization, a standard fact of commutative algebra beyond Chapter~I; this lets us handle the lower bound cleanly.)

(a) \emph{$\mathbf A^n$, $\mathbf P^n$.} $|\mathbf A^n|=|k^n|=\kappa$ for $n\ge1$, and $\mathbf A^n\subseteq\mathbf P^n$ with $|\mathbf P^n|\le|k^{n+1}|=\kappa$ gives $|\mathbf P^n|=\kappa$ too.

\emph{Affine varieties.} Let $Y$ be an affine variety of dimension $n\ge1$, $B=A(Y)$. By Noether normalization there are $z_1,\dots,z_n\in B$, algebraically independent over $k$, with $B$ a finitely generated module over $R=k[z_1,\dots,z_n]$ ($n=\operatorname{tr.deg}K(B)=\dim Y$, Thm.~1.8A). The map $\rho:Y\to\mathbf A^n$, $y\mapsto(z_1(y),\dots,z_n(y))$ is surjective with finite fibres: points of $Y$ are the maximal ideals of $B$ (Nullstellensatz); the points of $\mathbf A^n$ are the maximal ideals $\mathfrak n_c=(z_i-c_i)$ of $R$, and $\rho(\mathfrak M)=\mathfrak M\cap R$. By lying-over (integrality of $B$ over $R$) every $\mathfrak n_c$ is of the form $\mathfrak q\cap R$ for a prime $\mathfrak q$ of $B$, which is maximal because $\mathfrak q\cap R$ is; so $\rho$ is onto. The fibre over $c$ is the set of maximal ideals of $B/\mathfrak n_cB$, a finite-dimensional $k$-algebra (spanned by the images of finitely many $R$-module generators of $B$), which has finitely many maximal ideals. Hence $|\mathbf A^n|=\kappa\le|Y|\le\aleph_0\cdot\kappa=\kappa$, so $|Y|=\kappa$.

\emph{General varieties.} Let $X$ be a variety of dimension $n\ge1$. $X$ is covered by open affine subsets (Prop.~4.3), and is quasi-compact (it is noetherian, Ex.~1.7(b)), so a finite number of nonempty affine opens $U_1,\dots,U_s$ cover it. Each $U_i$ is an affine variety with $K(U_i)=K(X)$, so $\dim U_i=\operatorname{tr.deg}_kK(X)=\dim X=n\ge1$ (Thm.~3.2(d) for each $U_i$, and Ex.~1.10(b) for $\dim X=\sup\dim U_i$). By the affine case $|U_i|=\kappa$, so $\kappa=|U_1|\le|X|\le\sum|U_i|=\kappa$.

(b) A curve is a variety $C$ of dimension $1$. Let $Z\subsetneq C$ be a proper closed subset. $C$ is noetherian (Ex.~1.7), so $Z$ is a finite union of irreducible closed subsets $Z_i$ (Prop.~1.5), and $\dim Z_i<1$ by Ex.~1.10(d) (if $\dim Z_i=1=\dim C$ then $Z_i=C$). If $Z_i$ had two distinct points $p\ne q$, then $\{p\}\subsetneq Z_i$ would be a chain of irreducible closed sets of length $1$ (points are closed in a variety), contradicting $\dim Z_i=0$; so each $Z_i$ is a point and $Z$ is finite. Conversely every finite subset of $C$ is closed. So the closed subsets of $C$ are exactly $C$ and the finite subsets: the topology of a curve is the cofinite topology on its underlying set. If $C,C'$ are curves, then $|C|=|k|=|C'|$ by (a), so there is a bijection $C\to C'$; it carries finite sets to finite sets in both directions, hence is a homeomorphism.\hfill$\square$

\exercise{4.9}{Let $X$ be a projective variety of dimension $r$ in $\mathbf P^n$, with $n\ge r+2$. Show that for suitable choice of $P\notin X$, and a linear $\mathbf P^{n-1}\subseteq\mathbf P^n$, the projection from $P$ to $\mathbf P^{n-1}$ (Ex.~3.14) induces a birational morphism of $X$ onto its image $X'\subseteq\mathbf P^{n-1}$. You will need to use (4.6A), (4.7A), and (4.8A). This shows in particular that the birational map of (4.9) can be obtained by a finite number of projections.}
\pf
\emph{Notation and a lemma.} Recall (4.6A--4.8A): $k$ is perfect, so every finitely generated field extension $K/k$ is separably generated (4.8A), and every set of generators contains a separating transcendence basis (4.7A), i.e.\ a transcendence basis $B$ with $K/k(B)$ separable algebraic (hence finite).

\smallskip
\textbf{Lemma} (primitive element with coefficients in $k$). Let $F\supseteq k$ be fields, $k$ infinite, and let $L=F(\alpha,\beta)$ be a finite separable extension of $F$. Then for all but finitely many $c\in k$, $L=F(\alpha+c\beta)$.

\emph{Proof.} (This is the proof of 4.6A, noting the coefficients may be taken in the infinite subfield $k$.) Let $f,g\in F[x]$ be the minimal polynomials of $\alpha,\beta$, and in an algebraic closure $\Omega\supseteq L$ let $\alpha=\alpha_1,\dots,\alpha_p$ and $\beta=\beta_1,\dots,\beta_q$ be their roots; they are distinct by separability. Let $c\in k$ avoid the finitely many values $(\alpha_i-\alpha)/(\beta-\beta_j)$, $1\le i\le p$, $2\le j\le q$, and put $\gamma=\alpha+c\beta$, $h(x)=f(\gamma-cx)\in F(\gamma)[x]$. Then $g(\beta)=h(\beta)=0$. If $\beta_j$ is a common root of $g$ and $h$ then $\gamma-c\beta_j=\alpha_i$ for some $i$, i.e.\ $c(\beta-\beta_j)=\alpha_i-\alpha$, forcing $j=1$ by the choice of $c$. As $\beta$ is a simple root of $g$, $\gcd(g,h)=x-\beta$ in $\Omega[x]$; the gcd of two polynomials in $F(\gamma)[x]$ has coefficients in $F(\gamma)$ (Euclid's algorithm), so $\beta\in F(\gamma)$ and $\alpha=\gamma-c\beta\in F(\gamma)$. Thus $L=F(\gamma)$. $\square$

\smallskip
\emph{Setup.} If $r=0$, $X$ is a point and any $P\notin X$ works; assume $r\ge1$ (the argument below also works for $r=0$). After a projective linear change of coordinates assume $X\not\subseteq H_\infty=\{x_0=0\}$. Then $X_0=X\cap\{x_0\ne0\}$ is a nonempty affine open subset of $X$, $K=K(X)=K(X_0)$ is generated over $k$ by $t_i=x_i/x_0$ ($1\le i\le n$), and $\operatorname{tr.deg}_kK=r$ (Thm.~3.2(d) for each nonempty affine open $X\cap\{x_i\ne0\}$, all with function field $K$, together with Ex.~1.10(b)). Let $\Sigma=X\cap H_\infty$. Since $\dim H_\infty=n-1>r$, $H_\infty\not\subseteq X$, so $\Sigma\subsetneq H_\infty\cong\mathbf P^{n-1}$ is a proper closed subset.

\emph{Step 1: choice of coordinates.} I choose an invertible matrix $b\in GL_n(k)$ with columns $b_1,\dots,b_n$ such that all points $P_j=(0:b_j)\in H_\infty$ ($j=1,\dots,n$) are outside $X$. This is possible inductively: given $b_1,\dots,b_{j-1}$, the set $\Sigma\cup\mathbf P(\langle b_1,\dots,b_{j-1}\rangle)$ is a proper closed subset of $H_\infty\cong\mathbf P^{n-1}$ (a union of two proper closed subsets of an irreducible space is proper; the span is a proper linear subspace as $j\le n$), so we may choose $b_j$ with $(0:b_j)$ outside it. Put $a=b^{-1}$ and $s_i=\sum_ja_{ij}t_j\in K$; the $s_i$ again generate $K$ over $k$.

\emph{Step 2: separating basis.} By (4.8A) and (4.7A) some $r$ of $s_1,\dots,s_n$ form a separating transcendence basis. Permute the rows of $a$ (equivalently the columns of $b$; the set of points $P_j$ is unchanged, so they all remain outside $X$) so that $s_1,\dots,s_r$ is this basis. Since $r\le n-2$, the basis lies in $s_1,\dots,s_{n-2}$. Let $F=k(s_1,\dots,s_{n-2})\supseteq k(s_1,\dots,s_r)$. Then $K=F(s_{n-1},s_n)$ and $K/F$ is finite and separable ($K/k(s_1,\dots,s_r)$ is separable algebraic, so is its subextension $K/F$).

\emph{Step 3: primitive element and the point $P$.} By the Lemma there is a finite set $\Sigma_1\subseteq k$ with $K=F(s_{n-1}+cs_n)$ for $c\in k\setminus\Sigma_1$. For $c\in k$ let $S_c\in GL_n(k)$ be the identity with an extra entry $c$ at position $(n-1,n)$, $a_c=S_ca$, $b_c=a_c^{-1}=bS_c^{-1}$. The new coordinates $s^{(c)}=a_ct$ are $s_i$ for $i\ne n-1$ and $s_{n-1}+cs_n$ for $i=n-1$. The last column of $b_c$ is $b_n-cb_{n-1}$ (as $S_c^{-1}e_n=e_n-ce_{n-1}$), so $P(c)=(0:b_n-cb_{n-1})$ lies on the line $\ell\subseteq H_\infty$ through $P_{n-1}$ and $P_n$, and distinct $c$ give distinct points of $\ell$. Since $P_n\in\ell\setminus X$, $\ell\cap X$ is a proper closed subset of $\ell\cong\mathbf P^1$, hence finite. As $k$ is infinite we can choose $c\in k$ with $c\notin\Sigma_1$ and $P:=P(c)\notin X$. Write $s'_i=s^{(c)}_i$ and $a'=a_c$.

\emph{Step 4: the projection.} Let $T:\mathbf P^n\to\mathbf P^n$ be the projective linear automorphism $(x_0:x)\mapsto(x_0:a'x)=:(x_0:y_1:\dots:y_n)$. Then $T(P)=(0:a'b_c e_n)=(0:0:\dots:0:1)=:\tilde P\notin\tilde X:=T(X)$, and $T:X\to\tilde X$ is an isomorphism of varieties. Let $\pi:\mathbf P^n\setminus\{\tilde P\}\to\mathbf P^{n-1}$, $\pi(x_0:y_1:\dots:y_n)=(x_0:y_1:\dots:y_{n-1})$ be the projection from $\tilde P$ onto the hyperplane $\{y_n=0\}\not\ni\tilde P$; it is a morphism (Ex.~3.14(a)), so its restriction $\pi|_{\tilde X}$ is a morphism from $\tilde X\cong X$. (Projection from $P$ to another hyperplane differs from this by an isomorphism of $\mathbf P^{n-1}$.) Let $X'=\overline{\pi(\tilde X)}\subseteq\mathbf P^{n-1}$, an irreducible closed subset (Ex.~1.6), i.e.\ a projective variety; $\pi:\tilde X\to X'$ is a dominant morphism. (In fact $\pi(\tilde X)=X'$ is closed, since $\pi|_{\tilde X}$ is a finite morphism of projective varieties, cf.\ Thm.~II.4.9; this closedness uses later results and is not needed for the birationality statement.)

\emph{Step 5: birationality.} $X'\cap\{x_0\ne0\}\supseteq\pi(\tilde X\cap\{x_0\ne0\})\ne\emptyset$, so $K(X')$ is the field generated by the restrictions of the affine coordinates $y_i/x_0$ ($1\le i\le n-1$) of $\mathbf P^{n-1}\cap\{x_0\ne0\}=\mathbf A^{n-1}$ to $X'$. The pullback $\pi^*:K(X')\to K(\tilde X)=K(X)=K$ sends $y_i/x_0$ to $y_i/x_0|_{\tilde X}$, which under $T$ corresponds to $s'_i$. Hence $\pi^*(K(X'))=k(s'_1,\dots,s'_{n-1})$. This contains $F=k(s_1,\dots,s_{n-2})$ and $s'_{n-1}=s_{n-1}+cs_n$, hence $F(s'_{n-1})=K$ by the choice of $c$. So $\pi^*:K(X')\to K(X)$ is an isomorphism, and by Thm.~4.4/Cor.~4.5 the dominant morphism $\pi|_X:X\to X'$ is a birational map, i.e.\ a birational morphism of $X$ onto $X'$.

\emph{Remark.} $\dim X'=r$ (same function field transcendence degree). If $n-1\ge r+2$ one can repeat the construction with $X'\subseteq\mathbf P^{n-1}$; after $n-r-1$ projections one obtains a variety in $\mathbf P^{r+1}$ birational to $X$, i.e.\ a hypersurface. This shows that the birational map of Prop.~4.9 can be obtained by finitely many projections.\hfill$\square$

\exercise{4.10}{Let $Y$ be the cuspidal cubic $y^2=x^3$ in $\mathbf A^2$. Blow up the point $O=(0,0)$, let $E$ be the exceptional curve, and let $\tilde Y$ be the strict transform of $Y$. Show that $E$ meets $\tilde Y$ in one point, and that $\tilde Y\cong\mathbf A^1$. In this case the morphism $\varphi:\tilde Y\to Y$ is bijective and bicontinuous, but it is not an isomorphism.}
\pf
Recall: with homogeneous coordinates $(t:u)$ on $\mathbf P^1$, the blow-up of $\mathbf A^2$ at $O$ is $X=\{((x,y),(t:u)):xu=yt\}\subseteq\mathbf A^2\times\mathbf P^1$ with projection $\varphi:X\to\mathbf A^2$; the exceptional curve is $E=\varphi^{-1}(O)=\{O\}\times\mathbf P^1$, and the strict transform is $\tilde Y=\overline{\varphi^{-1}(Y\setminus\{O\})}$ (closure in $X$).

\emph{Computation of $\tilde Y$.} $X$ is covered by the open sets $X_t=\{t\ne0\}$ and $X_u=\{u\ne0\}$. On $X_t$ put $t=1$: the equation is $y=xu$, so $X_t\cong\mathbf A^2$ with coordinates $(x,u)$, $y=xu$. On $X_u$ put $u=1$: $x=yt$, $X_u\cong\mathbf A^2$ with coordinates $(y,t)$.

A point of $Y\setminus\{O\}$ has $x\ne0$ and $y\neq0$ (since $y^2=x^3$; $x=0$ forces $y=0$ and conversely). If $((x,y),(t:u))\in X$ lies over such a point, then $xu=yt$ with $x,y\ne0$ shows $t\ne0$ and $u\neq0$. Hence $S:=\varphi^{-1}(Y\setminus\{O\})\subseteq X_t\cap X_u$. In the chart $X_t$ with coordinates $(x,u)$: $S=\{(x,u):x\neq0,\ y^2=x^3\}=\{x\ne0,\ x^2u^2=x^3\}=\{x\ne0,\ x=u^2\}=\{(u^2,u):u\ne0\}$. The set $C=\{(x,u):x=u^2\}\subseteq X_t\cong\mathbf A^2$ is closed, irreducible (isomorphic to $\mathbf A^1$ via $u$), and $S=C\setminus\{(0,0)\}$ is dense in $C$; so the closure of $S$ in $X_t$ is $C$.

In the chart $X_u$ with coordinates $(y,t)$: $S=\{(y,t):y\ne0,\,t\neq0,\ y^2=(yt)^3\}=\{yt^3=1\}$, which is already closed in $X_u$, so the closure of $S$ in $X_u$ is $S$ itself. Since $X=X_t\cup X_u$, the closure of $S$ in $X$ is $\tilde Y=\overline S=(\overline S\cap X_t)\cup(\overline S\cap X_u)=C\cup S=C$. So
\[
\tilde Y=C\subseteq X_t,\qquad C=\{(x,u)=(u^2,u)\},
\]
and $u:\tilde Y\to\mathbf A^1$ is an isomorphism (inverse $u\mapsto(u^2,u)$, with $y=xu=u^3$). Hence $\tilde Y\cong\mathbf A^1$.

\emph{$E\cap\tilde Y$.} Points of $\tilde Y=C$ over $O$ have $x=0=y$, so $u^2=0$, $u=0$: the single point $((0,0),(1:0))$. (And $E\cap X_u$ contains no point of $\tilde Y$, consistent with $\tilde Y\subseteq X_t$.) So $E$ meets $\tilde Y$ in exactly one point.

\emph{The morphism $\varphi|_{\tilde Y}$.} In terms of the parameter $u$, $\varphi:\mathbf A^1\to Y$ is $u\mapsto(u^2,u^3)$.
\begin{itemize}
\item \emph{Bijective.} Injective: if $u^2=u'^2$ and $u^3=u'^3$ and $u\neq0$, then $u=u^3/u^2=u'^3/u'^2=u'$; if $u=0$ then $u'^2=0$, so $u'=0$. Surjective: $(0,0)=\varphi(0)$, and for $(a,b)\in Y$, $a\ne0$, put $u=b/a$: then $u^2=b^2/a^2=a^3/a^2=a$ and $u^3=u\cdot u^2=(b/a)a=b$.
\item \emph{Bicontinuous.} $Y$ is an irreducible curve ($y^2-x^3$ is irreducible by Ex.~4.4(b), so $Y$ is a hypersurface and $\dim Y=1$ by Prop.~1.13), so as in Ex.~4.8(b) its closed sets are $Y$ and the finite subsets; likewise for $\mathbf A^1$. A bijection between spaces with the cofinite topology is a homeomorphism.
\item \emph{Not an isomorphism.} If $\varphi:\mathbf A^1\to Y$ were an isomorphism of affine varieties, the induced map $\varphi^*:A(Y)\to A(\mathbf A^1)=k[u]$, $x\mapsto u^2$, $y\mapsto u^3$, would be an isomorphism of $k$-algebras (an isomorphism of varieties induces an isomorphism of rings of regular functions, and $\mathcal O(Y)=A(Y)$, $\mathcal O(\mathbf A^1)=k[u]$ by Thm.~3.2(a)). Its image is $k[u^2,u^3]$, which does not contain $u$ (no polynomial in $u^2,u^3$ has a nonzero term of degree $1$). So $\varphi^*$ is not surjective, and $\varphi$ is not an isomorphism. (This is Ex.~3.2(a) again.)
\end{itemize}
\hfill$\square$

\section*{5. Nonsingular varieties}

\textbf{Conventions and three lemmas used repeatedly.} Throughout, $k$ is algebraically closed. If $g$ is a polynomial, $\nabla g$ denotes the vector of its partial derivatives, and $\operatorname{ord}$ denotes the order of a polynomial or power series at the origin (the degree of its lowest non-zero homogeneous part). Where a characteristic restriction is needed it is stated.

\textbf{Lemma A (hypersurfaces).} Let $g\in k[x_1,\dots,x_n]$ be non-constant.
\begin{enumerate}[label=(\roman*),leftmargin=*]
\item If $g$ is irreducible, then $I(Z(g))=(g)$, $\dim Z(g)=n-1$, and a point $P\in Z(g)$ is nonsingular if and only if $\nabla g(P)\neq 0$.
\item If $P\in Z(g)$ and $\nabla g(P)\neq0$, then $P$ lies on exactly one irreducible component $W$ of $Z(g)$; moreover $g=hq$ with $h$ irreducible, $W=Z(h)$, $q(P)\neq 0$, and $P$ is a nonsingular point of $W$.
\end{enumerate}
\emph{Proof.} (i) $(g)$ is a prime ideal, hence radical, so $I(Z(g))=(g)$ by the Nullstellensatz, and $\dim Z(g)=n-1$ (Prop.~1.13). The Jacobian matrix of the single generator $g$ is the row $\nabla g(P)$, and the Jacobian criterion requires rank $n-r=1$, i.e.\ $\nabla g(P)\ne 0$. (ii) Write $g=\prod h_j^{e_j}$ with distinct irreducible $h_j$; the components of $Z(g)$ are the distinct hypersurfaces $Z(h_j)$. If $P\in Z(h_i)$ and either $e_i\ge2$ or $P\in Z(h_j)$ for some $j\ne i$, then $g=h_iq'$ with $q'(P)=0$, so $\nabla g(P)=q'(P)\nabla h_i(P)+h_i(P)\nabla q'(P)=0$, a contradiction. Hence $P$ lies on exactly one $Z(h)$ ($h=h_1$), with $e_1=1$, $g=hq$, $q(P)\neq0$, and $\nabla g(P)=q(P)\nabla h(P)$ gives $\nabla h(P)\neq0$. Apply (i) to $h$. \hfill$\square$

\textbf{Lemma B (blowing up a plane curve at the origin).} Let $f\in k[x,y]$ be squarefree, $f(0,0)=0$, $r=\operatorname{ord} f\ge1$, $f=f_r+f_{r+1}+\cdots$, $Y=Z(f)\subseteq\mathbf A^2$. Let $X=\{xu=yt\}\subseteq \mathbf A^2\times\mathbf P^1$ be the blow-up of $\mathbf A^2$ at $O=(0,0)$, with homogeneous coordinates $(t:u)$ on $\mathbf P^1$, and $\widetilde Y$ the closure in $X$ of $\varphi^{-1}(Y-O)$ (for reducible $Y$ apply this to each component; everything below is local and the argument is verbatim the same). $X$ is covered by $X_1=\{t\neq0\}\cong\mathbf A^2$ with coordinates $(x,u)$, $y=xu$, and $X_2=\{u\ne0\}\cong\mathbf A^2$ with coordinates $(y,t)$, $x=yt$. Then
\[
\widetilde Y\cap X_1=Z(g_1),\quad g_1(x,u)=x^{-r}f(x,xu)=f_r(1,u)+xf_{r+1}(1,u)+\cdots,
\]
\[
\widetilde Y\cap X_2=Z(g_2),\quad g_2(y,t)=y^{-r}f(yt,y)=f_r(t,1)+yf_{r+1}(t,1)+\cdots .
\]
In particular the points of $\widetilde Y$ over $O$ are the points $(x,u)=(0,u_0)$ with $f_r(1,u_0)=0$ in $X_1$ and $(y,t)=(0,t_0)$ with $f_r(t_0,1)=0$ in $X_2$, i.e.\ they correspond to the tangent directions at $O$.

\emph{Proof.} In $X_1$, $\varphi^{-1}(Y)=Z(f(x,xu))=Z(x^rg_1)=Z(x)\cup Z(g_1)$, and $\varphi^{-1}(Y-O)\cap X_1=Z(g_1)-Z(x)$ (the line $x=0$ is $\varphi^{-1}(O)$). Since $g_1(0,u)=f_r(1,u)$ is a non-zero polynomial, $x\nmid g_1$, so no irreducible component of $Z(g_1)$ lies in $Z(x)$; so $Z(g_1)-Z(x)$ is dense in $Z(g_1)$ and its closure is $Z(g_1)$. The same for $X_2$. \hfill$\square$

A linear change of coordinates $A\in GL_2$ acts on $\mathbf A^2\times\mathbf P^1$ by $(v,[w])\mapsto(Av,[Aw])$; since $xu-yt=\det\begin{pmatrix}x&y\\t&u\end{pmatrix}$ is multiplied by $\det A$, this preserves $X$ and is compatible with $\varphi$. Hence the blow-up of a plane curve at $O$ may be computed in any linear coordinates, and (after a translation) at any point.

\exercise{5.1}{Locate the singular points and sketch the following curves in $\mathbf A^2$ (assume $\operatorname{char}k\neq 2$). Which is which in Figure~4 (node, triple point, cusp, tacnode)?
\begin{enumerate}[label=(\alph*),leftmargin=*]
\item $x^2=x^4+y^4$;
\item $xy=x^6+y^6$;
\item $x^3=y^2+x^4+y^4$;
\item $x^2y+xy^2=x^4+y^4$.
\end{enumerate}}
\pf
Write $f$ for the polynomial obtained by moving everything to one side. In each case we find all common zeros of $f,f_x,f_y$; these form a finite set, so $f$ is squarefree (if $f=h^2q$ then $f,f_x,f_y$ would vanish on the infinite set $Z(h)$). By Lemma A, a point of $Z(f)$ with $\nabla f\neq0$ is nonsingular, and a point with $\nabla f=0$ is singular (if it lies on two components it is singular for trivial reasons; if on one component $Z(f_1)$, $f=f_1q$, then $\nabla f(P)=q(P)\nabla f_1(P)$ with $q(P)\ne0$, so $\nabla f_1(P)=0$).

\textbf{(a)} $f=x^2-x^4-y^4$, $f_x=2x-4x^3$, $f_y=-4y^3$. Since $\operatorname{char}k\ne2$, $f_y=0$ forces $y=0$, then $f=x^2(1-x^2)=0$ gives $x\in\{0,1,-1\}$; at $x=\pm1$, $f_x=\pm2(1-2)\neq0$. So the only singular point is $O=(0,0)$. The lowest form of $f$ is $x^2$: multiplicity $2$ with a single (double) tangent direction $x=0$. In fact, in $k[[x,y]]$ put $u=x\sqrt{1-x^2}$ (the square root exists as $\operatorname{char}k\neq2$); $(u,y)$ is a coordinate system there and $f=u^2-y^4=(u-y^2)(u+y^2)$: two smooth branches tangent to each other to order $2$. This is a \emph{tacnode}. (Real picture: $y^4=x^2(1-x^2)$ is the figure-eight lying on its side with the crossing point at the origin: two lobes, $-1\le x\le 1$, touching at $O$ with common vertical tangent.)

\textbf{(b)} $f=xy-x^6-y^6$, $f_x=y-6x^5$, $f_y=x-6y^5$. At a singular point we have $f=0$, so $x^6+y^6=xy$ and
\[
0=xf_x+yf_y=2xy-6(x^6+y^6)=2xy-6xy=-4xy,
\]
so $xy=0$ ($\operatorname{char}k\ne2$), and then $x^6+y^6=xy=0$ forces $x=y=0$. Thus $O$ is the only singular point, with lowest form $xy$: multiplicity $2$ and two distinct tangent directions $x=0,\ y=0$. This is a \emph{node}. (Real picture: two branches crossing at the origin like the coordinate axes, closing up into loops in the first and third quadrants.)

\textbf{(c)} $f=x^3-y^2-x^4-y^4$, $f_x=3x^2-4x^3=x^2(3-4x)$, $f_y=-2y(1+2y^2)$. If $x=0$: $f=-y^2(1+y^2)=0$ gives $y=0$ or $y^2=-1$, while $f_y=0$ gives $y=0$ or $y^2=-\tfrac12$; if $y\ne0$ we would need $-1=-\tfrac12$, i.e.\ $1=0$; so $y=0$. If $x\ne0$ then $f_x=0$ forces $x=\tfrac34$ (and $\operatorname{char}k\neq3$); $f_y=0$ gives $y=0$ or $y^2=-\tfrac12$; for $y=0$, $f=x^3(1-x)=\tfrac{27}{256}\ne0$; for $y^2=-\tfrac12$, $f=x^3(1-x)-y^2(1+y^2)=\tfrac{27}{256}+\tfrac14=\tfrac{91}{256}$, which vanishes only if $\operatorname{char}k\in\{7,13\}$. So for $\operatorname{char}k\ne2,7,13$ the origin is the only singular point; in characteristic $7$ or $13$ there are in addition the two points $(\tfrac34,\pm\sqrt{-1/2})$, which are nodes (there $f_{xy}=0$, $f_{xx}=6x(1-2x)=-\tfrac94\neq0$, $f_{yy}=-2-12y^2=4\neq0$, so $f_2=-\tfrac98X^2+2Y^2$ has two distinct factors). At $O$ the lowest form is $-y^2$: multiplicity $2$, one double tangent direction $y=0$. Moreover $f=(1-x)\bigl(x^3-Y^2\bigr)$ with $Y=y\sqrt{1+y^2}/\sqrt{1-x}\in k[[x,y]]$, and $(x,Y)$ are coordinates; so $O$ is analytically the singularity $Y^2=x^3$: a \emph{cusp}. (Real picture: a teardrop-shaped loop over $0\le x\le1$ with a cuspidal point at the origin pointing into $x<0$.)

\textbf{(d)} $f=x^2y+xy^2-x^4-y^4$, $f_x=2xy+y^2-4x^3$, $f_y=x^2+2xy-4y^3$. Now
\[
xf_x+yf_y=3(x^2y+xy^2)-4(x^4+y^4).
\]
At a singular point $f=0$, so $x^2y+xy^2=x^4+y^4=:t$ and $0=xf_x+yf_y=3t-4t=-t$. Hence $x^4+y^4=0$ and $xy(x+y)=0$: if $x=0$ then $y=0$; if $y=0$ then $x=0$; if $y=-x$ then $2x^4=0$ so $x=0$ ($\operatorname{char}k\neq2$). So $O$ is the only singular point. The lowest form is $f_3=xy(x+y)$: multiplicity $3$ with three distinct tangent directions $x=0$, $y=0$, $x+y=0$, an ordinary \emph{triple point}. (Real picture: three branches through the origin, forming loops.)

\textbf{Identification with Figure~4:} (a) tacnode, (b) node, (c) cusp, (d) triple point. \hfill$\square$

\exercise{5.2}{Locate the singular points and describe the singularities of the following surfaces in $\mathbf A^3$ (assume $\operatorname{char}k\neq2$). Which is which in Figure~5 (conical double point, double line, pinch point)?
\begin{enumerate}[label=(\alph*),leftmargin=*]
\item $xy^2=z^2$;
\item $x^2+y^2=z^2$;
\item $xy+x^3+y^3=0$.
\end{enumerate}}
\pf
Each polynomial $f$ below is irreducible, so $I(Z(f))=(f)$ and by Lemma A(i) a point of the surface $Z(f)$ is singular iff $\nabla f=0$ there.

\textbf{(a)} $f=xy^2-z^2$ is irreducible: as a polynomial in $z$ it is irreducible over $k(x,y)$ because $xy^2=y^2\cdot x$ is not a square in $k(x,y)$. $\nabla f=(y^2,\,2xy,\,-2z)$ vanishes iff $y=z=0$ (using $\operatorname{char}k\ne2$), and such points lie on the surface. So the singular locus is the whole $x$-axis $\{y=z=0\}$. At a point $(a,0,0)$ with $a\neq0$ the surface is locally $z^2=xy^2$ with $x$ a unit near $a$, i.e.\ two smooth sheets $z=\pm\sqrt{x}\,y$ crossing transversally along the axis. At the origin $x$ is not a unit, the two sheets are interchanged when $\sqrt x$ changes sign and cannot be separated: a \emph{pinch point} (Whitney umbrella) on the double line.

\textbf{(b)} $f=x^2+y^2-z^2$ is irreducible (as a polynomial in $z$, $x^2+y^2=(x+iy)(x-iy)$ is not a square). $\nabla f=(2x,2y,-2z)=0$ iff $x=y=z=0$. The only singular point is the origin, the vertex of the cone: a \emph{conical double point}.

\textbf{(c)} $f=xy+x^3+y^3$ does not involve $z$; the surface is the cylinder $C\times\mathbf A^1$ over the plane cubic $C:\ xy+x^3+y^3=0$. $f$ is irreducible: as a monic cubic in $y$ over $k[x]$, a factorization would give a root $y_0\in k[x]$ of $y^3+xy+x^3$ dividing the constant term $x^3$, so $y_0=cx^j$ ($0\le j\le3$), and the lowest and highest terms of $c^3x^{3j}+cx^{j+1}+x^3$ never cancel to zero (check $j=0,1,2,3$). Now $\nabla f=(y+3x^2,\;x+3y^2,\;0)$ and
\[
xf_x+yf_y=2xy+3(x^3+y^3)=2xy-3xy=-xy
\]
on the surface (where $x^3+y^3=-xy$). So at a singular point $xy=0$, hence $x^3+y^3=0$ and $x=y=0$. Conversely every point $(0,0,z)$ is on the surface with $\nabla f=0$. The singular locus is the $z$-axis. Along it the surface is locally the product of the node $xy+\dots=0$ of $C$ (two smooth branches with tangents $x=0$, $y=0$) with a line: a \emph{double line} (two sheets crossing transversally along a line, all along it).

\textbf{Identification with Figure~5:} (b) is the conical double point, (c) the double line, (a) the pinch point. \hfill$\square$

\exercise{5.3}{\emph{Multiplicities.} Let $Y\subseteq\mathbf A^2$ be a curve defined by the equation $f(x,y)=0$. Let $P=(a,b)$ be a point of $\mathbf A^2$. Make a linear change of coordinates so that $P$ becomes the point $(0,0)$. Then write $f$ as a sum $f=f_0+f_1+\dots+f_d$, where $f_i$ is a homogeneous polynomial of degree $i$ in $x$ and $y$. Then we define the \emph{multiplicity} of $P$ on $Y$, denoted $\mu_P(Y)$, to be the least $r$ such that $f_r\neq0$. (Note that $P\in Y\Leftrightarrow\mu_P(Y)>0$.) The linear factors of $f_r$ are called the \emph{tangent directions} at $P$.
\begin{enumerate}[label=(\alph*),leftmargin=*]
\item Show that $\mu_P(Y)=1\Leftrightarrow P$ is a nonsingular point of $Y$.
\item Find the multiplicity of each of the singular points in (Ex.~5.1) above.
\end{enumerate}}
\pf
Here $f$ is irreducible (or at least squarefree), so that $I(Y)=(f)$. First, $\mu_P(Y)$ does not depend on the choice of coordinates: if $\fm_P=(x-a,y-b)\subseteq k[x,y]$, then $\mu_P(Y)=\max\{r: f\in\fm_P^r\}$, since $\fm_P^r$ consists of the polynomials whose Taylor expansion at $P$ starts in degree $\ge r$, and this is coordinate-free.

\textbf{(a)} Take coordinates with $P=(0,0)$. Then $f_0=f(P)=0$ and $f_1=\frac{\partial f}{\partial x}(P)x+\frac{\partial f}{\partial y}(P)y$. So $\mu_P(Y)=1\iff f_1\ne0\iff\nabla f(P)\neq 0\iff P$ is nonsingular (Lemma A(i), Theorem~5.1; nonsingularity does not depend on the coordinates, because it is a property of the local ring). If $f$ is only known to be squarefree, use Lemma A(ii) and the remark in Ex.~5.1 in the same way.

\textbf{(b)} In each case the singular point is $O$ (cf.\ Ex.~5.1), and the multiplicity is read off from the lowest-degree form:
\begin{enumerate}[label=(\alph*),leftmargin=*]
\item $x^2=x^4+y^4$: $f_2=x^2$, so $\mu_O=2$ (one double tangent direction $x=0$);
\item $xy=x^6+y^6$: $f_2=xy$, $\mu_O=2$ (tangent directions $x=0,\ y=0$);
\item $x^3=y^2+x^4+y^4$: $f_2=-y^2$, $\mu_O=2$ (one double tangent direction $y=0$);
\item $x^2y+xy^2=x^4+y^4$: $f_3=xy(x+y)$, $\mu_O=3$ (tangent directions $x=0$, $y=0$, $x+y=0$).
\end{enumerate}
(In characteristic $7$ or $13$ the two extra singular points of (c) have $\mu=2$ with two distinct tangent directions, see Ex.~5.1(c).) \hfill$\square$

\exercise{5.4}{\emph{Intersection multiplicity.} If $Y,Z\subseteq\mathbf A^2$ are two distinct curves, given by equations $f=0$, $g=0$, and if $P\in Y\cap Z$, we define the \emph{intersection multiplicity} $(Y\cdot Z)_P$ of $Y$ and $Z$ at $P$ to be the length of the $\mathcal O_P$-module $\mathcal O_P/(f,g)$.
\begin{enumerate}[label=(\alph*),leftmargin=*]
\item Show that $(Y\cdot Z)_P$ is finite, and $(Y\cdot Z)_P\ge\mu_P(Y)\cdot\mu_P(Z)$.
\item If $P\in Y$, show that for almost all lines $L$ through $P$ (i.e.\ all but a finite number), $(L\cdot Y)_P=\mu_P(Y)$.
\item If $Y$ is a curve of degree $d$ in $\mathbf P^2$, and if $L$ is a line, $L\ne Y$, show that $(L\cdot Y)=d$. Here we define $(L\cdot Y)=\sum(L\cdot Y)_P$ taken over all points $P\in L\cap Y$, where $(L\cdot Y)_P$ is defined using a suitable affine cover of $\mathbf P^2$.
\end{enumerate}}
\pf
Here $\mathcal O_P=k[x,y]_{\fm_P}$ is the local ring of $\mathbf A^2$ at $P$, and the curves are $Y=Z(f)$, $Z=Z(g)$ with $f,g$ irreducible and not associate (distinct curves), hence coprime. Translate so that $P=O$, and write $\fm=\fm_O$ also for the maximal ideal of $\mathcal O_O$. Since the residue field is $k$, length equals $\dim_k$.

\textbf{(a) Finiteness.} $Y\not\subseteq Z$ (else $g\in I(Y)=(f)$, so $f\mid g$ and $f,g$ would be associate), so $Y\cap Z$ is a proper closed subset of the irreducible curve $Y$, hence of dimension $0$, hence finite. Let $A=k[x,y]$ and $I=(f,g)$. Every prime $\fp\supseteq I$ satisfies $Z(\fp)\subseteq Y\cap Z$, a finite set, and $Z(\fp)$ is irreducible, so it is one point and $\fp=I(Z(\fp))$ is maximal. So $A/I$ is a noetherian ring of dimension $0$, hence Artinian, and being a finitely generated $k$-algebra, finite dimensional over $k$. Then $\mathcal O_P/(f,g)=(A/I)_{\fm_P}$ is a localization of a finite dimensional algebra, hence a direct factor, and finite dimensional.

\textbf{(a) The inequality.} Let $r=\mu_O(Y)$ and $s=\mu_O(Z)\ge1$, so $f\in\fm^r$, $g\in\fm^s$. Put $N=r+s-1$ and $R_j=\mathcal O/\fm^j\cong k[x,y]/(x,y)^j$, of dimension $j(j+1)/2$ for $j\ge0$. The quotient $\mathcal O/(f,g,\fm^N)$ of $\mathcal O/(f,g)$ equals $R_N/(fR_N+gR_N)$. Multiplication by $f$ maps $\fm^{N-r}$ into $\fm^N$, so $fR_N$ is the image of $\mathcal O/\fm^{N-r}=R_{s-1}$, and $\dim fR_N\le s(s-1)/2$; likewise $\dim gR_N\le r(r-1)/2$. Therefore
\begin{align*}
(Y\cdot Z)_O&\ge\dim_k R_N/(fR_N+gR_N)\ge \frac{(r+s-1)(r+s)}2-\frac{s(s-1)}2-\frac{r(r-1)}2\\
&=\frac{(r+s)^2-(r+s)-s^2+s-r^2+r}{2}=rs.
\end{align*}
(If $P\notin Y$ or $P\notin Z$ the inequality is trivial.)

\textbf{(b)} Let $r=\mu_O(Y)$, $f=f_r+f_{r+1}+\cdots$. A line $L$ through $O$ is $\{(ut,vt)\}$ for a direction $(u:v)\in\mathbf P^1$, with equation $\ell=vx-uy$. If $u\neq0$ then $k[x,y]/(\ell)\cong k[t]$ via $x\mapsto ut$, $y\mapsto vt$ (for $u=0$ use $y\mapsto t$, $x\mapsto 0$), hence $\mathcal O/(\ell)\cong k[t]_{(t)}$, a discrete valuation ring with uniformizer $t$, and the image of $f$ is $f(ut,vt)=t^rf_r(u,v)+t^{r+1}f_{r+1}(u,v)+\cdots$. Thus
\[
(L\cdot Y)_O=\dim_k k[t]_{(t)}/(f(ut,vt))=\operatorname{ord}_tf(ut,vt)=r\quad\text{whenever } f_r(u,v)\neq0 .
\]
The non-zero form $f_r$ has at most $r$ zeros in $\mathbf P^1$ (its tangent directions), so the equality holds for all lines through $O$ except at most $r$ of them. (If $Y$ itself is a line through $P$ then $L=Y$ is among the exceptional lines.)

\textbf{(c)} Let $Y=Z(F)$ with $F$ homogeneous irreducible of degree $d$ and $L=Z(\lambda)$, $\lambda$ a linear form not proportional to $F$, so $\lambda\nmid F$. The local intersection multiplicity $(L\cdot Y)_P=\dim_k\mathcal O_{P,\mathbf P^2}/(\lambda/\mu,F/\mu^d)$, with $\mu$ any linear form with $\mu(P)\ne0$ (i.e.\ the dehomogenized equations in an affine chart containing $P$), does not depend on the chart, because dehomogenizations with respect to different $\mu$ differ by units of the local ring; it is also invariant under projective changes of coordinates. So choose coordinates $(x:y:z)$ with $L=\{z=0\}$. Put $\varphi(s,t)=F(s,t,0)$, a binary form of degree $d$; it is non-zero because $z\nmid F$. As $k$ is algebraically closed, $\varphi=c\prod_{i}(\beta_is-\alpha_it)^{m_i}$ with $(\alpha_i:\beta_i)$ distinct points of $\mathbf P^1$ and $\sum m_i=d$, and $L\cap Y=\{(\alpha_i:\beta_i:0)\}$.

Take $P_i=(\alpha_i:\beta_i:0)$ with, say, $\alpha_i\ne0$ (otherwise use the chart $y\neq0$). Use the chart $x\ne0$ with affine coordinates $Y'=y/x$, $Z'=z/x$; $L=\{Z'=0\}$, $P_i=(c,0)$ with $c=\beta_i/\alpha_i$. Then $\mathcal O_P/(Z',F(1,Y',Z'))\cong k[Y']_{(Y'-c)}/(\varphi(1,Y'))$, whose length is the order of vanishing of $\varphi(1,Y')$ at $Y'=c$, namely the multiplicity $m_i$ of the factor $\beta_is-\alpha_it$ in $\varphi$. Hence $(L\cdot Y)=\sum_im_i=d$. \hfill$\square$

\exercise{5.5}{For every degree $d>0$, and every $p=0$ or a prime number, give the equation of a nonsingular curve of degree $d$ in $\mathbf P^2$ over a field $k$ of characteristic $p$.}
\pf
We use Ex.~5.9 (proved below, independently of the present exercise): if $F$ is a form in $x,y,z$ such that at every point of $Z(F)\subseteq\mathbf P^2$ some partial derivative of $F$ is non-zero, then $F$ is irreducible and $Z(F)$ is a nonsingular curve (of degree $\deg F$).

\textbf{Case $p=0$ or $p\nmid d$.} Take the Fermat curve $F=x^d+y^d+z^d$. Its partials are $dx^{d-1},dy^{d-1},dz^{d-1}$ with $d\ne0$ in $k$; they vanish simultaneously only at $x=y=z=0$, which is not a point of $\mathbf P^2$.

\textbf{Case $p\mid d$.} Then $d\ge p\ge2$, and $d=0$ in $k$. Take
\[
F=x^{d-1}y+y^{d-1}z+z^d .
\]
Then $F_x=(d-1)x^{d-2}y$, $F_y=x^{d-1}+(d-1)y^{d-2}z$, $F_z=y^{d-1}+dz^{d-1}=y^{d-1}$. Let $P=(a:b:c)\in Z(F)$ with $F_z(P)=0$. Then $b=0$, so $0=F(P)=c^d$ gives $c=0$ and $P=(1:0:0)$, where $F_y(P)=1+(d-1)\cdot0=1\neq0$. Hence at every point of $Z(F)$ some partial is non-zero, and $Z(F)$ is a nonsingular curve of degree $d$. (E.g.\ $d=p=2$: $xy+yz+z^2$.) \hfill$\square$

\exercise{5.6}{\emph{Blowing Up Curve Singularities.}
\begin{enumerate}[label=(\alph*),leftmargin=*]
\item Let $Y$ be the cusp or node of (Ex.~5.1). Show that the curve $\widetilde Y$ obtained by blowing up $Y$ at $O=(0,0)$ is nonsingular (cf.\ (4.9.1) and (Ex.~4.10)).
\item We define a \emph{node} (also called \emph{ordinary double point}) to be a double point (i.e.\ a point of multiplicity $2$) of a plane curve with distinct tangent directions (Ex.~5.3). If $P$ is a node on a plane curve $Y$, show that $\varphi^{-1}(P)$ consists of two distinct nonsingular points on the blown-up curve $\widetilde Y$. We say that ``blowing up $P$ resolves the singularity at $P$''.
\item Let $P\in Y$ be the tacnode of (Ex.~5.1). If $\varphi:\widetilde Y\to Y$ is the blowing-up at $P$, show that $\varphi^{-1}(P)$ is a node. Using (b) we see that the tacnode can be resolved by two successive blowings-up.
\item Let $Y$ be the plane curve $y^3=x^5$, which has a ``higher order cusp'' at $O$. Show that $O$ is a triple point; that blowing up $O$ gives rise to a double point (what kind?) and that one further blowing up resolves the singularity.
\end{enumerate}
\emph{Note:} We will see later (V, 3.8) that any singular point of a plane curve can be resolved by a finite sequence of successive blowings-up.}
\pf
We use Lemma B and Lemma A(ii). Note that away from $\varphi^{-1}(O)$ the map $\varphi:\widetilde Y\to Y$ is an isomorphism $\widetilde Y-\varphi^{-1}(O)\cong Y-O$ (Section~4), so singular points of $\widetilde Y$ outside $\varphi^{-1}(O)$ are exactly the singular points of $Y$ other than $O$.

\textbf{(b)} Let $P$ be a node of $Y=Z(f)$; translating $P$ to $O$ we may, by the remark after Lemma B, use any linear coordinates. The lowest form is $f_2=\ell_1\ell_2$ with $\ell_1,\ell_2$ non-proportional linear forms; take them as new coordinates and scale $f$, so $f_2=xy$. Then $f=xy+f_3+f_4+\cdots$.

In $X_1$: $g_1(x,u)=u+xf_3(1,u)+x^2f_4(1,u)+\cdots$, so $g_1(0,u)=u$ and the only point of $\widetilde Y$ over $O$ in $X_1$ is $(x,u)=(0,0)$, i.e.\ the point $(O,(1:0))$ (the tangent direction $y=0$). There $\partial g_1/\partial u=1\neq0$, so by Lemma A(ii) $\widetilde Y=Z(g_1)$ is nonsingular at this point.

In $X_2$: $g_2(y,t)=t+yf_3(t,1)+\cdots$, with the only point over $O$ being $(y,t)=(0,0)$, i.e.\ $(O,(0:1))$ (direction $x=0$), and $\partial g_2/\partial t=1\ne0$ there: nonsingular. These two points are distinct, and these are all the points of $\varphi^{-1}(O)\cap\widetilde Y$.

\textbf{(a)} \emph{The node} (Ex.~5.1(b)): it is a node by (b), and $Y-O$ is nonsingular by Ex.~5.1(b); so $\widetilde Y$ is nonsingular.

\emph{The cusp} (Ex.~5.1(c)), $f=x^3-y^2-x^4-y^4$, $r=2$. In $X_1$ ($y=xu$):
\[
f(x,xu)=x^3-x^2u^2-x^4-x^4u^4=x^2\bigl(x-u^2-x^2-x^2u^4\bigr),\quad g_1=x-u^2-x^2(1+u^4).
\]
$g_1(0,u)=-u^2$, so the only point over $O$ is $(x,u)=(0,0)$, and $\partial g_1/\partial x=1-2x(1+u^4)=1\ne0$ there: nonsingular. In $X_2$ ($x=yt$): $g_2=y^{-2}(y^3t^3-y^2-y^4t^4-y^4)=yt^3-1-y^2(t^4+1)$, which is $-1\neq 0$ on $y=0$; so no points over $O$. Hence $\widetilde Y$ is nonsingular along $\varphi^{-1}(O)$, and elsewhere it is isomorphic to $Y-O$, nonsingular if $\operatorname{char}k\neq7,13$ (Ex.~5.1(c)); in characteristic $7,13$ the additional singular points of $Y$ remain singular on $\widetilde Y$, as they must, and have nothing to do with the blow-up at $O$. (Similarly for the standard cusp $y^2=x^3$: $g_1=x-u^2$ and $\partial g_1/\partial x=1$.)

\textbf{(c)} The tacnode (Ex.~5.1(a)): $f=x^2-x^4-y^4$, $r=2$, $f_2=x^2$.

In $X_1$: $g_1=1-x^2-x^2u^4$, and $g_1(0,u)=1\ne0$: no points over $O$ (the tangent direction is $x=0$, not in this chart).

In $X_2$: $f(yt,y)=y^2t^2-y^4t^4-y^4=y^2\bigl(t^2-y^2t^4-y^2\bigr)$, so $g_2=t^2-y^2(1+t^4)$. $g_2(0,t)=t^2$: the only point over $O$ is $Q=(y,t)=(0,0)$. The lowest form of $g_2$ at $Q$ is $t^2-y^2=(t-y)(t+y)$: two distinct tangent directions (char $\ne2$). So $Q$ is a point of multiplicity $2$ with distinct tangents: a \emph{node}, and $\varphi^{-1}(O)$ is the single point $Q$, a node of $\widetilde Y$. Now $\widetilde Y\cap X_2=Z(g_2)\subseteq\mathbf A^2$ is an affine neighbourhood of $Q$ in $\widetilde Y$; blowing up $\widetilde Y$ at $Q$ (an isomorphism outside $Q$) is computed in this plane curve, and by (b) the preimage of $Q$ in the second blow-up $\widetilde{\widetilde Y}$ consists of two distinct nonsingular points. As $\widetilde Y-Q$ is nonsingular, $\widetilde{\widetilde Y}$ is nonsingular: the tacnode is resolved by two successive blow-ups.

\textbf{(d)} $f=y^3-x^5$. The lowest form is $y^3$, so $\mu_O=3$: a triple point (one triple tangent direction $y=0$). Also $Y-O$ is nonsingular: $f_x=-5x^4$, $f_y=3y^2$ vanish together with $f$ only at $O$ (if $\operatorname{char}k=5$: $f_y=0\Rightarrow y=0\Rightarrow x^5=0$; if $\operatorname{char}k=3$: $f_x=0\Rightarrow x=0\Rightarrow y^3=0$; otherwise both $x=y=0$).

Blow up $O$. In $X_1$ ($y=xu$): $f(x,xu)=x^3u^3-x^5=x^3(u^3-x^2)$, so $g_1=u^3-x^2$ and $g_1(0,u)=u^3$: the only point over $O$ is $Q=(0,0)$. The lowest form of $g_1$ at $Q$ is $-x^2$: multiplicity $2$ with a single double tangent direction $x=0$; indeed $g_1=0$ is exactly the cusp $x^2=u^3$. In $X_2$ ($x=yt$): $f(yt,y)=y^3(1-y^2t^5)$, $g_2=1-y^2t^5$ is $1$ on $y=0$: no points over $O$. So after one blow-up $\widetilde Y$ has a double point of \emph{cusp} type at $Q$, and is nonsingular elsewhere.

Blow up $Q$ in the $(x,u)$-plane, $\widetilde Y\cap X_1=Z(u^3-x^2)$. We use the charts $u=xv$ (coordinates $(x,v)$) and $x=uw$ (coordinates $(u,w)$). First chart: $g(x,xv)=x^3v^3-x^2=x^2(xv^3-1)$, new equation $xv^3-1$, which is $-1$ on $x=0$: no points over $Q$. Second chart: $g(uw,u)=u^3-u^2w^2=u^2(u-w^2)$, new equation $h=u-w^2$; points over $Q$ ($u=0$): $w=0$; $\partial h/\partial u=1\ne0$: nonsingular. (Here the exceptional curve is $u=0$ and $\widetilde{\widetilde Y}$ is tangent to it.) Hence the second blow-up $\widetilde{\widetilde Y}$ is nonsingular: one further blowing up resolves the singularity. \hfill$\square$

\exercise{5.7}{Let $Y\subseteq\mathbf P^2$ be a nonsingular plane curve of degree $>1$, defined by the equation $f(x,y,z)=0$. Let $X\subseteq\mathbf A^3$ be the affine variety defined by $f$ (this is the cone over $Y$; see Ex.~2.10). Let $P$ be the point $(0,0,0)$, which is the \emph{vertex} of the cone. Let $\varphi:\widetilde X\to X$ be the blowing-up of $X$ at $P$.
\begin{enumerate}[label=(\alph*),leftmargin=*]
\item Show that $X$ has just one singular point, namely $P$.
\item Show that $\widetilde X$ is nonsingular (cover it with open affines).
\item Show that $\varphi^{-1}(P)$ is isomorphic to $Y$.
\end{enumerate}}
\pf
Let $d=\deg f\ge2$. As usual ``$Y$ is defined by $f$'' (and has degree $\deg f$) means that $f$ generates $I(Y)$; since $Y$ is irreducible, $f$ is then irreducible. So $I(X)=(f)\subset k[x,y,z]$, $X$ is irreducible of dimension $2$. We use Ex.~5.8 (below): a point $Q=(a:b:c)\in Y$ is nonsingular iff $\nabla f(a,b,c)\neq0$ (here $n=2$, $r=1$, so the rank must be $1$). So $\nabla f\neq0$ at every non-zero $(a,b,c)$ with $f(a,b,c)=0$.

\textbf{(a)} By Lemma A(i), a point $Q$ of $X$ is singular iff $\nabla f(Q)=0$. At $Q=0$, $\nabla f(0)=0$ since the partials have degree $d-1\ge1$; so $P$ is singular. A point $Q\ne0$ of $X$ defines a point $[Q]$ of $Y$, where $\nabla f(Q)\ne0$ as just noted; so $Q$ is nonsingular.

\textbf{(b)} Write $(x_1,x_2,x_3)$ for $(x,y,z)$. The blow-up of $\mathbf A^3$ at $P$ is $\{x_iu_j=x_ju_i\}\subseteq\mathbf A^3\times\mathbf P^2$, with homogeneous coordinates $(u_1:u_2:u_3)$ on $\mathbf P^2$, and $\widetilde X$ is the closure of $\varphi^{-1}(X-P)$. It is covered by the three charts $\{u_i\ne0\}\cong\mathbf A^3$. In $\{u_1\ne0\}$: coordinates $(x_1,u_2,u_3)$, $x_2=x_1u_2$, $x_3=x_1u_3$, and
\[
f(x_1,x_1u_2,x_1u_3)=x_1^d\,g(u_2,u_3),\qquad g(u_2,u_3)=f(1,u_2,u_3).
\]
Thus $\varphi^{-1}(X)\cap\{u_1\ne0\}=Z(x_1)\cup Z(g)$, and $Z(g)=\mathbf A^1\times\{g=0\}$ has no component inside $\{x_1=0\}$; so as in Lemma B, $\widetilde X\cap\{u_1\ne0\}=Z(g)\subseteq\mathbf A^3$.

We claim $\nabla g\neq0$ at every point of $Z(g)$, as a polynomial in $(x_1,u_2,u_3)$. Here $\nabla g=(0,g_{u_2},g_{u_3})$ with $g_{u_i}=f_{x_i}(1,u_2,u_3)$. Let $(x_1,u_2,u_3)\in Z(g)$, so $f(1,u_2,u_3)=0$. If $f_{x_2}(1,u_2,u_3)=f_{x_3}(1,u_2,u_3)=0$, then by Euler's identity $x_1f_{x_1}+x_2f_{x_2}+x_3f_{x_3}=d\cdot f$ (Ex.~5.8) at $(1,u_2,u_3)$ we get $f_{x_1}=0$ too, contradicting $\nabla f\neq0$ at the non-zero zero $(1,u_2,u_3)$ of $f$. Hence the claim. Since $\widetilde X$ is irreducible, $\widetilde X\cap\{u_1\ne0\}=Z(g)$ is irreducible (when non-empty) and by Lemma A(ii) every point of it is nonsingular. The same holds in the charts $u_2\ne0$, $u_3\ne0$. So $\widetilde X$ is nonsingular.

\textbf{(c)} $\varphi^{-1}(P)=\widetilde X\cap(\{P\}\times\mathbf P^2)$. In the chart $u_1\ne0$ it is $Z(g)\cap\{x_1=0\}=\{0\}\times\{(u_2,u_3):f(1,u_2,u_3)=0\}$, and similarly in the other charts. Hence $\varphi^{-1}(P)=\{P\}\times Y\subseteq\mathbf A^3\times\mathbf P^2$, a closed subvariety, which maps isomorphically to $Y$ by the second projection. \hfill$\square$

\exercise{5.8}{Let $Y\subseteq\mathbf P^n$ be a projective variety of dimension $r$. Let $f_1,\dots,f_t\in S=k[x_0,\dots,x_n]$ be homogeneous polynomials which generate the ideal of $Y$. Let $P\in Y$ be a point, with homogeneous coordinates $P=(a_0,\dots,a_n)$. Show that $P$ is nonsingular on $Y$ if and only if the rank of the matrix $\|(\partial f_i/\partial x_j)(a_0,\dots,a_n)\|$ is $n-r$. [\emph{Hint:} (a) Show that this rank is independent of the homogeneous coordinates chosen for $P$; (b) pass to an open affine $U_i\subseteq\mathbf P^n$ containing $P$ and use the affine Jacobian matrix; (c) you will need Euler's lemma, which says that if $f$ is a homogeneous polynomial of degree $d$, then $\sum x_i(\partial f/\partial x_i)=d\cdot f$.]}
\pf
\textbf{Euler's lemma.} For a monomial $x^\alpha$ of degree $d$, $\sum_ix_i\,\partial x^\alpha/\partial x_i=\sum_i\alpha_ix^\alpha=d\,x^\alpha$ (as an identity in $k[x_0,\ldots,x_n]$, with $d$ read in $k$); by linearity it holds for every form of degree $d$.

\textbf{(a)} Let $d_i=\deg f_i$. The partial $\partial f_i/\partial x_j$ is homogeneous of degree $d_i-1$ (or zero), so replacing $(a_j)$ by $(\lambda a_j)$, $\lambda\ne0$, multiplies the $i$-th row of the matrix by $\lambda^{d_i-1}\neq0$. This does not change the rank. (If $d_i=0$ the $f_i$ would be constants, impossible as $P\in Y$.)

\textbf{(b)} Choose $i$ with $a_i\ne0$; by (a) we may take $a_i=1$. Say $i=0$ (the other cases are the same). On $U_0=\{x_0\neq0\}\cong\mathbf A^n$, with coordinates $y_j=x_j/x_0$, the open subvariety $Y_0=Y\cap U_0$ is an affine variety of dimension $r$ (non-empty, as it contains $P$), and $O_{P,Y}=O_{P,Y_0}$. Let $g_l(y)=f_l(1,y_1,\dots,y_n)$. We claim $I(Y_0)=(g_1,\dots,g_t)$. Clearly $g_l\in I(Y_0)$. Conversely if $g\in I(Y_0)$ has degree $N$, put $\beta(g)=x_0^Ng(x_1/x_0,\dots,x_n/x_0)$, a form vanishing on $Y_0$, which is dense in $Y$ ($Y$ irreducible), hence $\beta(g)\in I(Y)=(f_1,\dots,f_t)$, $\beta(g)=\sum h_lf_l$, and putting $x_0=1$ gives $g=\sum h_l(1,y)g_l$.

By the definition of nonsingularity (Jacobian rank $n-r$; this does not depend on the generators, Thm~5.1), $P$ is nonsingular on $Y$ iff $\operatorname{rank}\|\partial g_l/\partial y_j(P)\|_{j=1..n}=n-r$. Here $\dfrac{\partial g_l}{\partial y_j}(P)=\dfrac{\partial f_l}{\partial x_j}(1,a_1,\dots,a_n)$ for $j\ge1$.

\textbf{(c)} Consider the full matrix $M=\|\partial f_l/\partial x_j(1,a_1,\ldots,a_n)\|_{j=0..n}$. Euler's lemma at the point $(1,a_1,\dots,a_n)$, where $f_l=0$, gives $\partial f_l/\partial x_0+\sum_{j\ge1}a_j\,\partial f_l/\partial x_j=d_lf_l=0$, i.e.\ column $0$ of $M$ is $-\sum_{j\ge1}a_j\cdot(\text{column }j)$. So $\operatorname{rank}M=\operatorname{rank}$ of the matrix without column $0$, which is the affine Jacobian. Combining with (a) and (b), $P$ is nonsingular iff $\operatorname{rank}\|(\partial f_l/\partial x_j)(a_0,\dots,a_n)\|=n-r$. \hfill$\square$

\exercise{5.9}{Let $f\in k[x,y,z]$ be a homogeneous polynomial, let $Y=Z(f)\subseteq\mathbf P^2$ be the algebraic set defined by $f$, and suppose that for every $P\in Y$, at least one of $(\partial f/\partial x)(P)$, $(\partial f/\partial y)(P)$, $(\partial f/\partial z)(P)$ is non-zero. Show that $f$ is irreducible (and hence that $Y$ is a nonsingular variety). [\emph{Hint:} Use (Ex.~3.7).]}
\pf
Assume $f$ is non-constant (so $Y\neq\emptyset$) and suppose $f=gh$ with $g,h$ non-constant; factors of a homogeneous polynomial are homogeneous. Let $g'$, $h'$ be irreducible factors of $g$, $h$. Then $Z(g')$, $Z(h')$ are irreducible curves in $\mathbf P^2$ (dimension $1$), and $1+1\ge2$, so by Ex.~3.7(b) $Z(g')\cap Z(h')\ne\emptyset$. Let $Q$ be a common point; then $g(Q)=h(Q)=0$, so $f(Q)=0$, i.e.\ $Q\in Y$, and
\[
\nabla f(Q)=g(Q)\nabla h(Q)+h(Q)\nabla g(Q)=0,
\]
contradicting the hypothesis. (This covers $g=h$ and repeated factors.) Hence $f$ is irreducible. Then $Y$ is an irreducible variety of dimension $1$, $I(Y)=(f)$, and by Ex.~5.8 (with $n=2$, $r=1$) $Y$ is nonsingular at every point since the $1\times3$ matrix $\nabla f(P)$ has rank $1=n-r$. \hfill$\square$

\exercise{5.10}{For a point $P$ on a variety $X$, let $\fm$ be the maximal ideal of the local ring $\mathcal O_P$. We define the \emph{Zariski tangent space} $T_P(X)$ of $X$ at $P$ to be the dual $k$-vector space of $\fm/\fm^2$.
\begin{enumerate}[label=(\alph*),leftmargin=*]
\item For any point $P\in X$, $\dim T_P(X)\ge\dim X$, with equality if and only if $P$ is nonsingular.
\item For any morphism $\varphi:X\to Y$, there is a natural $k$-linear map $T_P(\varphi):T_P(X)\to T_{\varphi(P)}(Y)$.
\item If $\varphi$ is the vertical projection of the parabola $x=y^2$ onto the $x$-axis, show that the induced map $T_0(\varphi)$ of tangent spaces at the origin is the zero map.
\end{enumerate}}
\pf
\textbf{(a)} $\mathcal O_P$ is a noetherian local ring with residue field $k=\mathcal O_P/\fm$, so by Prop.~5.2A $\dim_k\fm/\fm^2\ge\dim\mathcal O_P$, and $\dim\mathcal O_P=\dim X$ (Ex.~3.12). As $\dim T_P(X)=\dim_k(\fm/\fm^2)^*=\dim_k\fm/\fm^2$ we get $\dim T_P(X)\ge\dim X$, with equality iff $\mathcal O_P$ is a regular local ring, i.e.\ iff $P$ is nonsingular (definition).

\textbf{(b)} Let $Q=\varphi(P)$. The morphism $\varphi$ induces a $k$-algebra homomorphism $\varphi^*:\mathcal O_{Q,Y}\to\mathcal O_{P,X}$, $h\mapsto h\circ\varphi$ (if $h$ is regular on a neighbourhood $V$ of $Q$, then $h\circ\varphi$ is regular on $\varphi^{-1}(V)\ni P$; Ex.~3.3). If $h(Q)=0$ then $(h\circ\varphi)(P)=h(Q)=0$; so $\varphi^*(\fm_Q)\subseteq\fm_P$ and consequently $\varphi^*(\fm_Q^2)\subseteq\fm_P^2$. Hence $\varphi^*$ induces a $k$-linear map $\bar\varphi^*:\fm_Q/\fm_Q^2\to\fm_P/\fm_P^2$, and its dual $T_P(\varphi)=(\bar\varphi^*)^*:T_P(X)\to T_Q(Y)$ is the required $k$-linear map. (It is functorial: $T_P(\psi\circ\varphi)=T_{\varphi(P)}(\psi)\circ T_P(\varphi)$.)

\textbf{(c)} Here $X=Z(x-y^2)\subseteq\mathbf A^2$, $Y=\mathbf A^1$ (the $x$-axis), $\varphi(x,y)=x$, $P=O$. Then $\mathcal O_{O,Y}=k[x]_{(x)}$ with $\fm_Y=(x)$, so $\fm_Y/\fm_Y^2$ is one-dimensional, spanned by the class of $x$. The coordinate ring of $X$ is $k[x,y]/(x-y^2)\cong k[y]$ (with $x=y^2$), so $\mathcal O_{O,X}=k[y]_{(y)}$, $\fm_X=(y)$, $\fm_X^2=(y^2)$. Now $\varphi^*(x)=x|_X=y^2\in\fm_X^2$, so $\bar\varphi^*:\fm_Y/\fm_Y^2\to\fm_X/\fm_X^2$ sends the generator to $0$; it is the zero map, and so is its dual $T_O(\varphi)$. (Geometrically: the tangent line at $O$ to the parabola is vertical and is killed by the projection.) \hfill$\square$

\exercise{5.11}{\emph{The Elliptic Quartic Curve in $\mathbf P^3$.} Let $Y$ be the algebraic set in $\mathbf P^3$ defined by the equations $x^2-xz-yw=0$ and $yz-xw-zw=0$. Let $P$ be the point $(x,y,z,w)=(0,0,0,1)$, and let $\varphi$ denote the projection from $P$ to the plane $w=0$. Show that $\varphi$ induces an isomorphism of $Y-P$ with the plane cubic curve $y^2z-x^3+xz^2=0$ minus the point $(1,0,-1)$. Then show that $Y$ is an irreducible nonsingular curve. It is called the \emph{elliptic quartic curve} in $\mathbf P^3$. Since it is defined by two equations it is another example of a complete intersection (Ex.~2.17).}
\pf
We assume $\operatorname{char}k\ne2$; this is necessary: in characteristic $2$ the equations are satisfied by all points $(1:0:1:w)$, so $Y$ contains the line through $P$ and $(1:0:1:0)$, and also the point $(0:1:0:0)$ off this line; thus $Y$ is reducible and the statement fails.

Put (1) $x^2-xz-yw=0$, (2) $yz-xw-zw=0$, i.e.\ $yw=x(x-z)$ and $(x+z)w=yz$. $P\in Y$. Let $C=Z(F)\subseteq\mathbf P^2$ (coordinates $(x:y:z)$), $F=y^2z-x^3+xz^2$, and $Q=(1:0:-1)$; $F(Q)=0-1+1=0$ so $Q\in C$. The projection $\varphi:\mathbf P^3-P\to\mathbf P^2$, $(x:y:z:w)\mapsto(x:y:z)$, is a morphism.

\textbf{$\varphi(Y-P)\subseteq C-Q$.} Let $(x:y:z:w)\in Y-P$; then $(x,y,z)\ne0$. Multiplying (2) by $y$ and substituting (1): $y^2z=yw(x+z)=x(x-z)(x+z)$, i.e.\ $F(x,y,z)=0$. If $(x:y:z)=Q=(1:0:-1)$ then (1) gives $0=x(x-z)=2\ne0$ (with the representative $(1,0,-1,w)$): impossible. So $\varphi(Y-P)\subseteq C-Q$.

\textbf{Inverse morphism.} Points of $C$ with $y=0$ and $x+z=0$: then $F=0$ is automatic, and they are $(x:0:-x)=Q$. Thus $C-Q=V_1\cup V_2$ with $V_1=\{y\ne0\}$, $V_2=\{x+z\ne0\}$ (open in $C-Q$). Define
\[
\psi_1:V_1\to\mathbf P^3,\ (x:y:z)\mapsto(xy:y^2:yz:x(x-z)),\qquad
\psi_2:V_2\to\mathbf P^3,\ (x:y:z)\mapsto(x(x+z):y(x+z):z(x+z):yz).
\]
Each is given by forms of the same degree without common zero on its domain (in $\psi_1$, $y^2\ne0$; in $\psi_2$, $x+z\ne0$ gives some coordinate $\ne0$), so these are morphisms. On $V_1\cap V_2$ they agree: $\psi_2$'s coordinates equal $(x+z)/y$ times those of $\psi_1$, because $\frac{x+z}{y}\,x(x-z)=\frac{x(x^2-z^2)}{y}=\frac{y^2z}y=yz$ on $C$. So we get a morphism $\psi:C-Q\to\mathbf P^3$. Its image lies in $Y-P$: for $\psi_1$, with $w=x(x-z)/y$ one has $yw=x(x-z)$ and $(x+z)w=x(x^2-z^2)/y=y^2z/y=yz$, so (1),(2) hold; and the point is not $P$ since $y^2\ne0$. For $\psi_2$ put $(X,Y,Z,W)=(x(x+z),y(x+z),z(x+z),yz)$; on $C$ one computes directly $X^2-XZ-YW=(x+z)\bigl(x^3-xz^2-y^2z\bigr)=-(x+z)F=0$ and $YZ-XW-ZW=yz(x+z)^2-(x+z)^2yz=0$, and the point is not $P$ because $x+z\neq0$ makes $(X,Y,Z)\ne0$. Clearly $\varphi\circ\psi=\mathrm{id}_{C-Q}$ (drop the last coordinate and cancel the common factor).

Conversely let $(x:y:z:w)\in Y-P$. If $y\ne0$, (1) gives $w=x(x-z)/y$ so $\psi_1(\varphi(x:y:z:w))=(xy:y^2:yz:x(x-z))=y\cdot(x:y:z:w)$, the same point. If $x+z\ne0$, (2) gives $w=yz/(x+z)$ and similarly $\psi_2\varphi=\mathrm{id}$. Every point of $Y-P$ has $y\ne0$ or $x+z\neq0$, since $y=0=x+z$ implies, with (1), $x(x-z)=2x^2=0$ so $x=z=0$ and the point is $P$. Hence $\psi\circ\varphi=\mathrm{id}_{Y-P}$ and $\varphi:Y-P\to C-Q$ is an isomorphism.

\textbf{$C$ is nonsingular and irreducible.} $F_x=-3x^2+z^2$, $F_y=2yz$, $F_z=y^2+2xz$. Suppose $F_x=F_y=F_z=0$ at a point of $C$ (off $C$ this can happen, e.g.\ at $(1:0:0)$ in characteristic $3$). $F_y=0$ gives $y=0$ or $z=0$. If $z=0$: $F=-x^3$ forces $x=0$ and then $F_z=y^2$ forces $y=0$, not a point. If $y=0$ (take $z=1$): $F_z=2x=0$ gives $x=0$, and then $F_x=1\ne0$. So $\nabla F$ never vanishes on $C$, and by Ex.~5.9 $F$ is irreducible and $C$ is a nonsingular curve.

\textbf{$Y$ is irreducible and nonsingular.} $Y-P\cong C-Q$ is an irreducible nonsingular curve. It remains to see that $P\in\overline{Y-P}$ and that $Y$ is nonsingular at $P$. In the chart $w=1$ (coordinates $x,y,z$), (1) gives $y=x^2-xz$, and substituting into (2): $h(x,z)=x^2z-xz^2-x-z=0$. So $Y\cap\{w\ne0\}$ is isomorphic (via $(x,y,z)\mapsto(x,z)$, with inverse $(x,z)\mapsto(x,x^2-xz,z)$) to $Z(h)\subseteq\mathbf A^2$, with $P\leftrightarrow(0,0)$. Since $\nabla h(0,0)=(-1,-1)\neq0$, by Lemma A(ii) $(0,0)$ lies on exactly one irreducible component $W$ of $Z(h)$, a curve ($\dim W=1$), and is a nonsingular point of $W$. $W$ has infinitely many points, so $P$ is in the closure of $\{\text{points of }Y-P\text{ on }W\}\subseteq Y-P$. Hence $Y=\overline{Y-P}$ is irreducible, $Y\cap\{w\neq0\}$ is irreducible, so $Z(h)=W$, and $Y$ is nonsingular at $P$ (Lemma A(ii)). As nonsingularity is a property of the local ring, and $Y-P\cong C-Q$ is nonsingular, $Y$ is an irreducible nonsingular curve. \hfill$\square$

\exercise{5.12}{\emph{Quadric Hypersurfaces.} Assume $\operatorname{char}k\ne2$, and let $f$ be a homogeneous polynomial of degree $2$ in $x_0,\dots,x_n$.
\begin{enumerate}[label=(\alph*),leftmargin=*]
\item Show that after a suitable linear change of variables, $f$ can be brought into the form $f=x_0^2+\dots+x_r^2$ for some $0\le r\le n$.
\item Show that $f$ is irreducible if and only if $r\ge2$.
\item Assume $r\ge2$, and let $Q$ be the quadric hypersurface in $\mathbf P^n$ defined by $f$. Show that the singular locus $Z=\operatorname{Sing}Q$ of $Q$ is a \emph{linear variety} (Ex.~2.11) of dimension $n-r-1$. In particular, $Q$ is nonsingular if and only if $r=n$.
\item In case $r<n$, show that $Q$ is a cone with axis $Z$ over a nonsingular quadric $Q'\subseteq\mathbf P^r$. (This notion of cone generalizes the one defined in Ex.~2.10.) If $Y$ is a closed subset of $\mathbf P^r$, and if $Z$ is a linear subspace of dimension $n-r-1$ in $\mathbf P^n$, we embed $\mathbf P^r$ in $\mathbf P^n$ so that $\mathbf P^r\cap Z=\emptyset$, and define the \emph{cone over $Y$ with axis $Z$} to be the union of all lines joining a point of $Y$ to a point of $Z$.
\end{enumerate}}
\pf
Throughout $f\ne0$.

\textbf{(a)} Write $f=\sum_{i\le j}c_{ij}x_ix_j$, i.e.\ $f=\sum_{i,j}a_{ij}x_ix_j$ with $a_{ij}=a_{ji}$ ($a_{ii}=c_{ii}$, $a_{ij}=c_{ij}/2$ for $i<j$; here $\operatorname{char}k\ne2$ is used). We induct on the number of variables occurring. If $f$ is the zero form there is nothing to do; otherwise: if some $a_{ii}\ne0$, permute variables so that $a_{00}\ne0$ and complete the square: $f=a_{00}\bigl(x_0+\sum_{j>0}\tfrac{a_{0j}}{a_{00}}x_j\bigr)^2+f'(x_1,\dots,x_n)$. If all $a_{ii}=0$ but some $a_{ij}\ne0$ ($i\ne j$), substitute $x_i\mapsto x_i+x_j$; then the coefficient of $x_j^2$ becomes $2a_{ij}\ne0$ and we are in the previous case. Putting $x_0'=x_0+\sum_{j>0}\frac{a_{0j}}{a_{00}}x_j$ and applying induction to $f'$ in $x_1,\dots,x_n$, we reach $f=\sum_{i=0}^rc_ix_i^2$ with all $c_i\ne0$ (after renaming variables). Replacing $x_i$ by $c_i^{-1/2}x_i$ ($k$ is algebraically closed) gives $f=x_0^2+\dots+x_r^2$.

\textbf{(b)} Define the \emph{rank} of $f$ as the rank of its symmetric matrix $A=(a_{ij})$ ($f=x^TAx$). Under a linear change of variables $x=My$ ($M$ invertible), $A$ becomes $M^TAM$, of the same rank. So the rank of $f$ is $r+1$ in the form of (a). If $r=0$, $f=x_0^2$ is a square; if $r=1$, $f=(x_0+ix_1)(x_0-ix_1)$ with $i^2=-1$. So $f$ is reducible for $r\le1$. If $r\ge 2$ and $f=\ell_1\ell_2$ with linear forms $\ell_1,\ell_2$ (a reducible quadratic form is a product of two linear forms), then $f$ has matrix $\tfrac12(\ell_1\ell_2^T+\ell_2\ell_1^T)$ of rank $\le2<r+1$, contradiction. So $f$ is irreducible iff $r\ge2$.

\textbf{(c)} Since $r\ge2$ and $f=x_0^2+\dots+x_r^2$, $f$ is irreducible (b), $I(Q)=(f)$ and $\dim Q=n-1$. By Ex.~5.8 (with $r_{Q}=n-1$), a point $P=(a_0:\dots:a_n)\in Q$ is singular iff the $1\times(n+1)$ matrix $\nabla f(a)=(2a_0,\dots,2a_r,0,\dots,0)$ has rank $\ne1$, i.e.\ iff it is zero, iff $a_0=\dots=a_r=0$ (as $\operatorname{char}k\ne 2$). Such points automatically satisfy $f=0$. Hence
\[
Z=\operatorname{Sing}Q=\{x_0=\dots=x_r=0\}\subseteq\mathbf P^n,
\]
a linear variety defined by $r+1$ independent linear forms, of dimension $n-(r+1)=n-r-1$ (empty iff $r=n$). So $Q$ is nonsingular iff $r=n$. (A linear change of coordinates preserves the notions involved.)

\textbf{(d)} Let $r<n$, embed $\mathbf P^r=\{x_{r+1}=\dots=x_n=0\}$ in $\mathbf P^n$; $Z=\{x_0=\dots=x_r=0\}$ has dimension $n-r-1$ and $\mathbf P^r\cap Z=\emptyset$. Let $Q'=Z(x_0^2+\dots+x_r^2)\subseteq\mathbf P^r$; by (c) applied in $\mathbf P^r$ (the rank $r+1$ equals the number of variables), $Q'$ is a nonsingular quadric, non-empty (it contains $(1:i:0:\dots:0)$). Claim: $Q$ is the union of the lines joining a point of $Q'$ to a point of $Z$.

($\supseteq$) If $[b]\in Q'$ and $[c]\in Z$, then a point of the line is $[\lambda b+\mu c]$, with $b=(b_0,..,b_r,0,..,0)$, $c=(0,..,0,c_{r+1},..,c_n)$, and $f(\lambda b+\mu c)=\sum_{i\le r}\lambda^2b_i^2=\lambda^2f(b)=0$.

($\subseteq$) Let $[a]\in Q$, $a'=(a_0,\dots,a_r,0,\dots,0)$, $a''=(0,\dots,0,a_{r+1},\dots,a_n)$, so $f(a')=f(a)=0$. If $a'\neq0$ and $a''\ne0$ then $[a']\in Q'$, $[a'']\in Z$ and $[a]=[a'+a'']$ lies on the line joining them. If $a''=0$, $[a]=[a']\in Q'$ lies on the line joining itself to any point of $Z\neq\emptyset$. If $a'=0$, $[a]\in Z$ lies on the line joining it to any point of $Q'\ne\emptyset$.

Hence $Q$ is the cone over the nonsingular quadric $Q'\subseteq\mathbf P^r$ with axis $Z=\operatorname{Sing}Q$. \hfill$\square$

\exercise{5.13}{It is a fact that any regular local ring is an integrally closed domain (Matsumura [2, Th.~36, p.~121]). Thus we see from (5.3) that any variety has a nonempty open subset of normal points (Ex.~3.17). In this exercise, show directly (without using (5.3)) that the set of non-normal points of a variety is a proper closed subset. (You will need the finiteness of integral closure: see (3.9A).)}
\pf
Let $N=\{P\in X:\ \mathcal O_P\text{ is not integrally closed}\}$, $X$ a variety. Since $\mathcal O_{P,X}=\mathcal O_{P,U}$ for every open $U\ni P$, $N\cap U$ is the set of non-normal points of $U$. Closedness of a subset is local, and $X$ is irreducible, so a non-empty open affine $U$ shows: if $N\cap U$ is closed and proper in $U$, then $N$ is proper; and if $X=\bigcup U_\alpha$ with $N\cap U_\alpha$ closed in $U_\alpha$ then $N$ is closed. Hence it suffices to treat $X$ affine, and to show $N$ is closed and $\neq X$.

So let $A=A(X)$, a finitely generated $k$-domain with quotient field $K$, and $B$ the integral closure of $A$ in $K$. By Prop.~3.9A, $B$ is a finitely generated $A$-module: $B=Ab_1+\dots+Ab_s$, $b_i\in K$. Let $\mathfrak c=\{a\in A:\ aB\subseteq A\}$, the conductor, an ideal of $A$. It is non-zero: writing $b_i=p_i/q_i$ with $p_i,q_i\in A$, $q_i\ne0$, the product $q=\prod q_i\neq0$ satisfies $qb_i\in A$ for all $i$, so $qB\subseteq A$, $q\in\mathfrak c$.

\textbf{Claim:} $N=Z(\mathfrak c)$. This shows $N$ is closed, and $N\ne X$ since $\mathfrak c\ne0$: choose $0\ne q\in \mathfrak c$; the function $q$ is not identically zero on the irreducible $X$, so $Z(\mathfrak c)\subseteq Z(q)\subsetneq X$.

\emph{$N\subseteq Z(\mathfrak c)$.} Let $P\notin Z(\mathfrak c)$: there is $a\in\mathfrak c$ with $a(P)\ne0$. Then $B\subseteq a^{-1}A\subseteq A_a$, so $B_a=A_a$. Since the integral closure of $A_a$ in $K$ is $B_a$ (integral closure commutes with localization, Atiyah--Macdonald 5.12), $A_a$ is integrally closed. And $\mathcal O_P$ is a localization of $A_a$ (as $a\notin\fm_P$), hence integrally closed (a localization of an integrally closed domain is integrally closed). So $P\notin N$.

\emph{$Z(\mathfrak c)\subseteq N$.} Let $P\notin N$, i.e.\ $\mathcal O_P$ is integrally closed. Since $B$ is integral over $A\subseteq\mathcal O_P$, $B\subseteq\mathcal O_P$; so $b_i=u_i/v_i$ with $u_i,v_i\in A$, $v_i(P)\ne0$. Then $v=\prod v_i\in A$, $v(P)\neq 0$, $vb_i\in A$ for all $i$, so $vB\subseteq A$, i.e.\ $v\in\mathfrak c$ and $P\notin Z(\mathfrak c)$.

Hence the set of non-normal points is a proper closed subset of $X$. \hfill$\square$

\exercise{5.14}{\emph{Analytically Isomorphic Singularities.}
\begin{enumerate}[label=(\alph*),leftmargin=*]
\item If $P\in Y$ and $Q\in Z$ are analytically isomorphic plane curve singularities, show that the multiplicities $\mu_P(Y)$ and $\mu_Q(Z)$ are the same (Ex.~5.3).
\item Generalize the example in the text (5.6.3) to show that if $f=f_r+f_{r+1}+\dots\in k[[x,y]]$, and if the leading form $f_r$ of $f$ factors as $f_r=g_sh_t$, where $g_s,h_t$ are homogeneous of degrees $s$ and $t$ respectively, and have no common linear factor, then there are formal power series
\[
g=g_s+g_{s+1}+\cdots,\qquad h=h_t+h_{t+1}+\cdots
\]
in $k[[x,y]]$ such that $f=gh$.
\item Let $Y$ be defined by the equation $f(x,y)=0$ in $\mathbf A^2$, and let $P=(0,0)$ be a point of multiplicity $r$, so that when $f$ is expanded as a polynomial in $x$ and $y$, we have $f=f_r+$ higher terms. We say that $P$ is an \emph{ordinary $r$-fold point} if $f_r$ is a product of $r$ distinct linear factors. Show that any two ordinary double points are analytically isomorphic. Ditto for ordinary triple points. But show that there is a one-parameter family of mutually non-isomorphic ordinary $4$-fold points.
\item Assume $\operatorname{char}k\ne2$. Show that any double point of a plane curve is analytically isomorphic to the singularity at $(0,0)$ of the curve $y^2=x^r$, for a uniquely determined $r\ge2$. If $r=2$ it is a node (Ex.~5.6). If $r=3$ we call it a \emph{cusp}; if $r=4$ a \emph{tacnode}. See (V, 3.9.5) for further discussion. (Part (d) is starred in the book.)
\end{enumerate}}
\pf
Let $S=k[[x,y]]$, $\fm_S=(x,y)$. Points are moved to the origin by a translation. For a plane curve $Y=Z(f)$ ($f$ irreducible) with $P=O$ we use the standard facts: $\hat{\mathcal O}_{P,Y}\cong S/(f)$ (completion commutes with quotients of noetherian rings, $\hat{\mathcal O}_{O,\mathbf A^2}=S$; Thm~5.4A, Atiyah--Macdonald 10.13), and $\hat{\mathcal O}/\hat\fm^n\cong\mathcal O/\fm^n$, compatibly with multiplication so that the associated graded rings agree (Thm~5.4A(a), A--M 10.15, 10.22). Two singularities are analytically isomorphic if $\hat{\mathcal O}_{P,Y}\cong\hat{\mathcal O}_{Q,Z}$ as $k$-algebras; such an isomorphism of local rings preserves maximal ideals and their powers.

\textbf{Auxiliary fact (coordinates).} If $g,h\in\fm_S$ have linearly independent linear parts, then $x\mapsto g,\ y\mapsto h$ defines an automorphism of $S$ (formal inverse function theorem): the induced map on the associated graded rings $k[x,y]\to k[x,y]$ is the linear isomorphism, hence an isomorphism, and a continuous homomorphism of complete filtered rings inducing an isomorphism on associated graded rings is an isomorphism (A--M 10.23, 10.24). We then call $(g,h)$ a \emph{coordinate system} of $S$.

\textbf{(a)} Let $O=\mathcal O_{P,Y}$, $\fm$ its maximal ideal, $A=k[x,y]$, $\fm_A=(x,y)$, $f=f_r+\cdots$ with $r=\mu_P(Y)$, and $R_n=A/\fm_A^n$, $\dim R_n=n(n+1)/2$. Then $O/\fm^n=A/(f,\fm_A^n)=R_n/fR_n$. For $n\ge r$ the map $R_{n-r}\to R_n$, $p\mapsto fp$, is well defined ($f\fm_A^{n-r}\subseteq\fm_A^n$) and injective: if $p\notin\fm_A^{n-r}$ has lowest form $p_s$ with $s<n-r$, the lowest form of $fp$ is $f_rp_s\neq0$ (polynomial ring is a domain) of degree $r+s<n$, so $fp\notin\fm_A^n$. Hence for $n\ge r$
\[
\dim_kO/\fm^n=\frac{n(n+1)}2-\frac{(n-r)(n-r+1)}2=rn-\frac{r(r-1)}2,
\]
so $\dim_k\fm^n/\fm^{n+1}=r$ for all $n\ge r$; thus $r$ is the eventual value of $\dim_k\fm^n/\fm^{n+1}$. By the facts above this is the same for $\hat{\mathcal O}_{P,Y}$, and an isomorphism $\hat{\mathcal O}_{P,Y}\cong\hat{\mathcal O}_{Q,Z}$ preserves $\hat\fm^n/\hat\fm^{n+1}$ up to isomorphism. Hence $\mu_P(Y)=\mu_Q(Z)$.

\textbf{(b)} If $s=0$ or $t=0$ take $g=1$, $h=f$ (or vice versa), so let $s,t\ge1$; $r=s+t$. "No common linear factor" means $g_s,h_t$ are coprime in $k[x,y]$ (an irreducible form in two variables is linear since $k$ is algebraically closed).

\emph{Lemma.} For every $N\ge s+t-1$, every form $R$ of degree $N$ can be written $R=Ag_s+Bh_t$ with $A,B$ forms of degrees $N-s,N-t$. Indeed, let $S_j$ be the space of forms of degree $j$ ($\dim S_j=j+1$ for $j\ge0$, $0$ for $j<0$) and consider $\theta:S_{N-s}\oplus S_{N-t}\to S_N$, $(A,B)\mapsto Ag_s+Bh_t$. If $Ag_s+Bh_t=0$, then $h_t\mid Ag_s$, so $h_t\mid A$ (coprime in the UFD $k[x,y]$), $A=Ch_t$, $B=-Cg_s$ with $C\in S_{N-s-t}$; conversely such pairs are in the kernel. Hence $\dim\ker\theta=\dim S_{N-s-t}=N-s-t+1$ (as $N-s-t+1\ge0$), and
\[
\dim\operatorname{im}\theta=(N-s+1)+(N-t+1)-(N-s-t+1)=N+1=\dim S_N.
\]
So $\theta$ is onto.

Now construct $g_{s+m},h_{t+m}$ by induction on $m\ge1$. The degree $r+m$ part of $gh$ is
\[
g_sh_{t+m}+g_{s+m}h_t+\sum_{\substack{s<i<s+m\\ j=r+m-i}}g_ih_j .
\]
We need this to equal $f_{r+m}$, i.e.\ $h_tg_{s+m}+g_sh_{t+m}=R_m:=f_{r+m}-\sum_{s<i<s+m}g_ih_{r+m-i}$, where the sum involves only $g_i,h_j$ with $i<s+m$, $j<t+m$, already constructed. As $\deg R_m=r+m\ge s+t+1\ge s+t-1$, the Lemma gives $g_{s+m}=B$, $h_{t+m}=A$. The series $g,h$ so obtained satisfy $gh=f$ (degree by degree). (For $m=0$, $g_sh_t=f_r$ by hypothesis.)

\textbf{(c)} \emph{Double points.} $f_2=\ell_1\ell_2$ with $\ell_1,\ell_2$ non-proportional linear forms. By (b) ($s=t=1$), $f=gh$ with $g=\ell_1+\cdots$, $h=\ell_2+\cdots$. As $\ell_1,\ell_2$ are independent, $(g,h)$ is a coordinate system of $S$, i.e.\ there is an automorphism $\sigma$ of $S$ with $\sigma(x)=g,\sigma(y)=h$, and $\sigma(xy)=f$. Hence $S/(f)\cong S/(xy)$, so every ordinary double point is analytically isomorphic to $xy=0$, and any two are isomorphic.

\emph{Triple points.} $f_3=\ell_1\ell_2\ell_3$, distinct lines. Apply (b) to $f=g_1\cdot h$ ($g_1=\ell_1+\dots$, $h=\ell_2\ell_3+\cdots$; $\ell_1$ is coprime to $\ell_2\ell_3$), and again to $h$ to get $f=g_1g_2g_3$ with $g_i=\ell_i+$ higher terms. Take $X=g_1$, $Y=g_2$ as a coordinate system; then $g_3\in k[[X,Y]]$ with linear part $\ell_3=aX+bY+\ldots$ in these coordinates, where $a,b\neq 0$ because $\ell_3$ is proportional neither to $\ell_1$ nor to $\ell_2$. Write $g_3=\alpha X+\beta Y$ with $\alpha,\beta\in k[[X,Y]]$, $\alpha(0)=a$, $\beta(0)=b$ units. Then $X'=\alpha X$, $Y'=\beta Y$ have independent linear parts $aX,bY$, so $(X',Y')$ is a coordinate system, and
\[
f=XYg_3=\alpha^{-1}\beta^{-1}\,X'Y'(X'+Y').
\]
So $(f)=(X'Y'(X'+Y'))$ in the coordinates $X',Y'$, and $S/(f)\cong S/(xy(x+y))$: all ordinary triple points are analytically isomorphic.

\emph{Four-fold points.} For $\lambda\in k\setminus\{0,1\}$ let $f_\lambda=xy(x-y)(x-\lambda y)$, so $O$ is an ordinary $4$-fold point of $Y_\lambda=Z(f_\lambda)$ (union of four lines; the argument below only uses the leading form, so it applies equally to any plane curve with this leading form). We show: if $\hat{\mathcal O}_{O,Y_\lambda}\cong\hat{\mathcal O}_{O,Y_\mu}$ then $\mu$ lies in the orbit $\{\lambda,\,1/\lambda,\,1-\lambda,\,1/(1-\lambda),\,\lambda/(\lambda-1),\,(\lambda-1)/\lambda\}$ (at most six elements). Since $k$ is infinite there are then infinitely many isomorphism classes, a one-parameter family.

Let $f=f_4+\cdots$, $f'=f'_4+\cdots$ be two ordinary 4-fold points with $\hat{\mathcal O}\cong\hat{\mathcal O}'$. The graded ring $\operatorname{gr}(\mathcal O)=\bigoplus\fm^n/\fm^{n+1}$ is $k[x,y]/(f_4)$: indeed $(f)\cap\fm_A^n=f\fm_A^{n-4}$ (since $\operatorname{ord}(fp)=4+\operatorname{ord}p$), so $\fm^n/\fm^{n+1}=\fm_A^n/(\fm_A^{n+1}+f\fm_A^{n-4})$, i.e.\ forms of degree $n$ modulo $f_4\cdot$(forms of degree $n-4$). Hence an isomorphism of the completions induces a graded $k$-algebra isomorphism $\Phi:k[x,y]/(f_4)\to k[x',y']/(f'_4)$. In degree $1$ both sides are the full space of linear forms ($4\ge2$); let $\mathcal A:k[x,y]\to k[x',y']$ be the isomorphism of polynomial rings with $\mathcal A|_{\text{degree }1}=\Phi_1$. The two ring homomorphisms $k[x,y]\to k[x',y']/(f_4')$, $\Phi\circ\pi$ and $\pi'\circ\mathcal A$, agree on $x,y$, hence are equal, so $\mathcal A(f_4)\in(f'_4)$; being of degree $4$, $\mathcal A(f_4)=cf'_4$ with $c\in k^*$. Thus $\mathcal A$ induces a projectivity of $\mathbf P^1$ carrying the four roots of $f_4$ in $\mathbf P^1$ onto the four roots of $f'_4$. The roots of $f_\lambda$ are $(0:1),(1:0),(1:1),(\lambda:1)$. The cross ratio of four distinct points of $\mathbf P^1$ (the image of the 4th point under the projectivity sending the first three to $\infty,0,1$) is invariant under projectivities and takes the six values above when the points are reordered. The cross ratio of $(\infty,0,1,\lambda)$ is $\lambda$, so $\lambda$ lies in the orbit of $\mu$, as claimed.

\textbf{(d)} Let $\operatorname{char}k\ne2$ and $P=O$ a double point of $Y=Z(f)$, $f$ irreducible, $f=f_2+f_3+\cdots$, $f_2\ne0$, so $\hat{\mathcal O}_{P,Y}=S/(f)$.

\emph{Step 0: $f$ is not of the form $ug^2$ in $S$ with $u$ a unit and $g\in\fm_S$.} Since $k$ is algebraically closed, $f$ irreducible implies that not both $f_x,f_y$ are zero (else in characteristic $p$, $f\in k[x^p,y^p]$ would be a $p$-th power). Say $\partial f\neq0$ for a partial derivative $\partial$ in one variable; $\deg\partial f<\deg f$ so $f\nmid\partial f$, and $f,\partial f$ are coprime. By Ex.~5.4(a)'s argument, $\mathcal O_{P,\mathbf A^2}/(f,\partial f)$ has finite length, so $(f,\partial f)$ is $\fm_S$-primary in $S$. If $f=ug^2$ then $\partial f=g\cdot(\dots)\in(g)$, so $(f,\partial f)\subseteq(g)\subsetneq S$ with $S/(g)$ of dimension $1$, not $\fm_S$-primary. Contradiction.

\emph{Step 1: normal form.} Choose a direction $(a:b)$ with $f_2(a,b)\ne0$ ($f_2\ne0$, $k$ infinite) and, by a linear change of coordinates, assume $f_2(0,1)\neq0$. Then $f(0,y)=f_2(0,1)y^2+\cdots$ has order exactly $2$, so by the Weierstrass preparation theorem $f=u\cdot W$, $u$ a unit of $S$, $W=y^2+a(x)y+b(x)$ with $a,b\in xk[[x]]$. Since $\operatorname{char}k\ne2$ put $y_1=y+a(x)/2$; $(x,y_1)$ is a coordinate system and $W=y_1^2-c(x)$ with $c=a^2/4-b\in xk[[x]]$. The change of coordinates preserves the order function, so $\operatorname{ord}W=\operatorname{ord}f=2$, whence $\operatorname{ord}c\ge2$. Also $c\ne0$: otherwise $f=uy_1^2$, excluded by Step 0. So $c=x^r\varepsilon$ with $r=\operatorname{ord}c\ge2$ and $\varepsilon\in k[[x]]^\times$. As $\operatorname{char}k\ne2$ and $k$ is algebraically closed, $\varepsilon$ has a square root in $k[[x]]$ (binomial series / Hensel). Put $y_2=y_1\varepsilon^{-1/2}$; $(x,y_2)$ is a coordinate system and $y_1^2-x^r\varepsilon=\varepsilon(y_2^2-x^r)$. So $f=u\varepsilon(y_2^2-x^r)$ and
\[
\hat{\mathcal O}_{P,Y}\cong S/(y^2-x^r)\quad\text{for some }r\ge2.
\]
(Remark: if $f_2=y^2+a_1xy+b_2x^2$ up to the factor $u(0)$, then $c=(a_1^2/4-b_2)x^2+\dots$, so $r=2$ iff $a_1^2-4b_2\ne0$ iff $f_2$ has two distinct linear factors, i.e.\ iff $P$ is a node in the sense of Ex.~5.6.)

\emph{Step 2: uniqueness of $r$.} Let $R_r=S/(y^2-x^r)$ and suppose $R_r\cong R_s$ ($r,s\ge2$); we show $r=s$. Put $F=y^2-x^r$, $F'=y^2-x^s$ (both have order $2$). An isomorphism $\theta:R_r\to R_s$ preserves maximal ideals, and $\fm_R/\fm_R^2\cong\fm_S/\fm_S^2$ since $F\in\fm_S^2$. Let $\pi:S\to R_r$, $\pi':S\to R_s$ be the quotient maps and lift $\theta\pi(x),\theta\pi(y)\in\fm_{R_s}$ to elements $\sigma(x),\sigma(y)\in\fm_S$. They have independent linear parts (because $\theta$ induces a bijection $\fm_{R_r}/\fm_{R_r}^2\to\fm_{R_s}/\fm_{R_s}^2$, and these are $\fm_S/\fm_S^2$ and its analogue), so $x\mapsto\sigma(x),y\mapsto\sigma(y)$ defines an automorphism $\sigma$ of $S$ with $\pi'\sigma=\theta\pi$ (both are continuous homomorphisms agreeing on $x,y$). Taking kernels, $\sigma^{-1}(F'S)=FS$, i.e.\ $\sigma(F)S=F'S$. Consequently, for $g\in S$ with $\operatorname{ord}g=1$,
\[
\dim_kS/(F,g)=\dim_kS/(\sigma F,\sigma g)=\dim_kS/(F',\sigma g)
\]
and $\operatorname{ord}\sigma(g)=1$. Define $\nu(F)=\sup\{\dim_kS/(F,g):\ g\in S,\ \operatorname{ord}g=1\}\in\mathbf N\cup\{\infty\}$; then $\nu(F)=\nu(F')$.

\emph{Computation of $\nu$.} If $g$ has order $1$ then $S/(g)\cong k[[\tau]]$ for a suitable coordinate $\tau\in\{x,y\}$, and $\dim_kS/(F,g)=\operatorname{ord}_\tau(F|_{g=0})$ (or $\infty$ if this restriction is $0$).

\begin{itemize}
\item \emph{$r$ odd.} Write $g=\alpha x+\beta y+(\text{order}\ge2)$. If $\alpha\neq0$, then $g=0$ is $x=\phi(y)$, $\phi\in yk[[y]]$ and $F|=y^2-\phi^r$ has order $2$ ($r\ge3$). If $\alpha=0$, $\beta\ne0$, then $g=0$ is $y=\psi(x)$ with $\psi\in x^2k[[x]]$ (since $g$ has no linear $x$-term), and $F|=\psi^2-x^r$ has order $\min(2\operatorname{ord}\psi,r)\le r$, because $2\operatorname{ord}\psi\ne r$ ($r$ odd), and order $r$ if $\psi=0$. So all values are $\le r$, and $g=y$ gives $\dim S/(y^2-x^r,y)=\dim k[[x]]/(x^r)=r$. Hence $\nu(F)=r$.
\item \emph{$r=2m$ even.} Take $g_N=y-x^m+x^N$, $N>m$; its order is $1$. On $g_N=0$, $y=x^m-x^N$, and $F=(y-x^m)(y+x^m)\mapsto(-x^N)(2x^m-x^N)=-x^{N+m}(2-x^{N-m})$ of order $N+m$ ($\operatorname{char}k\neq2$). So $\nu(F)=\infty$.
\end{itemize}
Thus $\nu(F)=\nu(F')$ shows $r,s$ have the same parity, and if both are odd then $r=s$. If $r=2m$, $s=2m'$, then $R_r=S/(p_1p_2)$ with $p_{1,2}=y\mp x^m$; these are prime in the UFD $S$ (as $S/(p_i)\cong k[[x]]$) and distinct, so the minimal primes of $R_r$ are $\bar p_1,\bar p_2$, and $R_r/(\bar p_1+\bar p_2)=S/(y-x^m,y+x^m)=S/(y,x^m)$ has dimension $m$ ($\operatorname{char}k\ne2$). This number is intrinsic to the ring $R_r$ (the minimal primes are permuted by an isomorphism), hence $m=m'$ and $r=s$.

This proves that every double point of a plane curve is analytically isomorphic to $y^2=x^r$ for a unique $r\ge2$; $r=2$ is exactly the node. \hfill$\square$

\exercise{5.15}{\emph{Families of Plane Curves.} A homogeneous polynomial $f$ of degree $d$ in three variables $x,y,z$ has $\binom{d+2}2$ coefficients. Let these coefficients represent a point in $\mathbf P^N$, where $N=\binom{d+2}2-1=\frac12d(d+3)$.
\begin{enumerate}[label=(\alph*),leftmargin=*]
\item Show that this gives a correspondence between points of $\mathbf P^N$ and algebraic sets in $\mathbf P^2$ which can be defined by an equation of degree $d$. The correspondence is $1$--$1$ except in some cases where $f$ has a multiple factor.
\item Show under this correspondence that the (irreducible) nonsingular curves of degree $d$ correspond $1$--$1$ to the points of a nonempty Zariski-open subset of $\mathbf P^N$. [\emph{Hints:} (1) Use elimination theory (5.7A) applied to the homogeneous polynomials $\partial f/\partial x_0,\dots,\partial f/\partial x_n$; (2) use the previous exercises (Ex.~5.5, 5.8, 5.9) above.]
\end{enumerate}}
\pf
Let $V_d$ be the space of forms of degree $d$ in $x,y,z$, of dimension $N+1$; a non-zero form $f$ gives a point $[f]\in\mathbf P(V_d)=\mathbf P^N$.

\textbf{(a)} The map $\Phi:[f]\mapsto Z(f)\subseteq\mathbf P^2$ is well defined (rescaling $f$ does not change $Z(f)$), and its image is, by definition, the set of algebraic sets definable by an equation of degree $d$. We show it is injective on the set $\Sigma$ of squarefree $[f]$, and that for $[f]\in\Sigma$ no other point of $\mathbf P^N$ maps to $Z(f)$. Let $f$ be squarefree and $g$ any form of degree $d$ with $Z(g)=Z(f)$. By the (homogeneous) Nullstellensatz $I(Z(g))=\sqrt{(g)}=(g_{\mathrm{red}})$, where $g_{\mathrm{red}}$ is the product of the distinct irreducible factors of $g$, and $I(Z(f))=(f)$. So $(f)=(g_{\mathrm{red}})$ and $f=cg_{\mathrm{red}}$, $c\in k^*$. Now $\deg g_{\mathrm{red}}\le\deg g=d=\deg f$ with equality iff $g$ is squarefree, so $g=g_{\mathrm{red}}$ and $g=c^{-1}f$: $[g]=[f]$. Thus $\Phi$ is one-to-one on the squarefree locus (and does not identify a squarefree form with a non-squarefree one). Non-injectivity can only occur among forms with a multiple factor, and does occur there (e.g.\ $x^2y$ and $xy^2$ for $d=3$ both give $Z=\{xy=0\}$).

\textbf{(b)} Let $F=\sum a_{ijk}x^iy^jz^k$ be the generic form of degree $d$. Let $W\subseteq\mathbf P^N$ be the set of $[f]$ for which the four forms $f,\,f_x,\,f_y,\,f_z$ have a common zero in $\mathbf P^2$.

\emph{Claim 1: $[f]\notin W$ iff $f$ is irreducible and $Z(f)$ is nonsingular, i.e.\ $Z(f)$ is a nonsingular curve of degree $d$.} (A curve of degree $d$ is $Z(f)$ with $f$ irreducible, $I(Z(f))=(f)$; a power $f=g^e$, $e\ge2$, is excluded: $Z(g)$ is a nonsingular curve but of degree $\deg g<d$, and $f,f_x,f_y,f_z$ all vanish on $Z(g)$.) If $[f]\notin W$, then at every point of $Z(f)$ one of $f_x,f_y,f_z$ is non-zero, so by Ex.~5.9 $f$ is irreducible and $Z(f)$ is a nonsingular curve. Conversely, if $f$ is irreducible and $Z(f)$ nonsingular, then $I(Z(f))=(f)$ and by Ex.~5.8 ($n=2$, $r=1$) the matrix $\nabla f(P)$ has rank $1$ at every $P\in Z(f)$, so $f,f_x,f_y,f_z$ have no common zero. (Note that in characteristic $p\mid d$ one cannot omit $f$ here: the three partials may have a common zero off the curve, e.g.\ $xy+yz+z^2$ in characteristic $2$ at $(1:0:1)$.)

\emph{Claim 2: $W$ is Zariski-closed.} (For $d=1$ every line is a nonsingular curve and $f_x,f_y,f_z$ are constants not all zero, so $W=\emptyset$; assume $d\ge2$, so that all four forms below have positive degree.) Apply the elimination theory of 5.7A to four forms $\Phi_0,\dots,\Phi_3$ in $x,y,z$ with indeterminate coefficients, of degrees $d,d-1,d-1,d-1$: there are polynomials $g_1,\dots,g_m$ in these indeterminates, homogeneous in the coefficients of each $\Phi_i$ separately, such that for any values of the coefficients in $k$ the forms have a common zero in $\mathbf P^2$ iff all $g_\nu$ vanish. Substitute for the coefficients of $\Phi_0,\Phi_1,\Phi_2,\Phi_3$ the coefficients of $F,F_x,F_y,F_z$, which are integer multiples of the $a_{ijk}$ (linear forms in the $a$'s): we obtain polynomials $G_\nu(a_{ijk})$, homogeneous in the $a_{ijk}$, since each $g_\nu$ is homogeneous in the coefficients of each $\Phi_i$ separately and the substituted coefficients are linear forms in the $a_{ijk}$ (the degree of $G_\nu$ is the sum of the four degrees of $g_\nu$). By the elimination statement, $[f]\in W$ iff $G_\nu(f)=0$ for all $\nu$. So $W=Z(G_1,\dots,G_m)$ is closed in $\mathbf P^N$.

\emph{Claim 3: $U=\mathbf P^N\setminus W$ is non-empty.} By Ex.~5.5 there is a nonsingular curve of degree $d$ in $\mathbf P^2$ over $k$ for every $d>0$ and every characteristic (Fermat curve if $\operatorname{char}k\nmid d$, and $x^{d-1}y+y^{d-1}z+z^d$ if $\operatorname{char}k\mid d$); by Claim 1 its equation gives a point of $U$.

Hence $U$ is a non-empty open subset of $\mathbf P^N$; by Claim 1 its points are exactly the $[f]$ with $f$ irreducible and $Z(f)$ nonsingular, and since such $f$ are squarefree, part (a) shows that $[f]\mapsto Z(f)$ is injective on $U$ and every nonsingular curve of degree $d$ is obtained from exactly one point of $U$. So the irreducible nonsingular curves of degree $d$ correspond $1$--$1$ to the points of the non-empty open set $U$. \hfill$\square$

\section*{6. Nonsingular curves}

Throughout, $k$ is algebraically closed, ``curve'' means a variety of dimension~$1$, and numbers refer to Hartshorne, Chapter~I. For a function field $K/k$ of dimension~$1$, $C_K$ is the abstract nonsingular curve of Section~6 (the set of discrete valuation rings of $K/k$, cofinite topology). A discrete valuation ring (DVR) $R$ with $k\subseteq R\subsetneq K$ is a point of $C_K$. I use freely the following consequences of the text.

\begin{itemize}
\item (Overrings of a DVR.) If $R\subseteq R'\subseteq K=\operatorname{Frac}R$ with $R$ a DVR, then $R'=R$ or $R'=K$. Indeed if $x\in R'\setminus R$ then $x=u\pi^{-n}$ with $u\in R^\times$, $n>0$, so $\pi^{-1}=\pi^{n-1}u^{-1}x\in R'$ and then $R'=K$.
\item (Lemma~A.) \emph{If $X$ is a nonsingular projective curve with function field $K$, then every DVR $R$ of $K/k$ equals $\mathcal O_{Q,X}$ for some $Q\in X$; i.e.\ $X\to C_K$, $Q\mapsto\mathcal O_{Q,X}$, is an isomorphism.}
\end{itemize}
\emph{Proof of Lemma~A.} By (6.7) and its proof, $\theta:X\to C_K$, $Q\mapsto\mathcal O_Q$, is an isomorphism of $X$ onto an open subset $U\subseteq C_K$. The morphism $\theta^{-1}:U\to X$ goes to a projective variety, and $C_K\setminus U$ is finite, so applying (6.8) one point at a time (and gluing the extensions, which agree on $U$) we get a morphism $\psi:C_K\to X$ extending $\theta^{-1}$. Let $R\in C_K$, $Q=\psi(R)\in X$. Since $\psi$ is dominant, $\psi^*$ gives $\mathcal O_{Q,X}\subseteq R$ inside $K$. Now $\mathcal O_{Q,X}$ is a DVR and $R\ne K$, so $R=\mathcal O_{Q,X}$ by the overring remark, i.e.\ $R=\theta(Q)\in U$. Hence $U=C_K$. \hfill$\square$

\exercise{6.1}{Recall that a curve is \emph{rational} if it is birationally equivalent to $\mathbf P^1$ (Ex.~4.4). Let $Y$ be a nonsingular rational curve which is not isomorphic to $\mathbf P^1$.
\begin{enumerate}[label=(\alph*),leftmargin=*]
\item Show that $Y$ is isomorphic to an open subset of $\mathbf A^1$.
\item Show that $Y$ is affine.
\item Show that $A(Y)$ is a unique factorization domain.
\end{enumerate}}
\pf
(a) Let $K=K(Y)$. Since $Y$ is birational to $\mathbf P^1$, $K\cong k(t)=K(\mathbf P^1)$ (Cor.~4.5). By (6.7) $Y$ is isomorphic to an open subset $U$ of $C_K$, and $C_K\cong C_{k(t)}\cong\mathbf P^1$ by Lemma~A applied to $X=\mathbf P^1$. So $Y$ is isomorphic to a nonempty open subset $U'\subseteq\mathbf P^1$, and $U'\ne\mathbf P^1$ because $Y\not\cong\mathbf P^1$. Thus $\mathbf P^1\setminus U'$ is a nonempty finite set. The group $PGL(1)$ acts transitively on $\mathbf P^1$ (Ex.~6.6(a): $x\mapsto 1/(x-a)$ sends $a$ to $\infty$), so after an automorphism of $\mathbf P^1$ we may assume $\infty\notin U'$, i.e.\ $U'\subseteq\mathbf A^1$ is open: $U'=\mathbf A^1\setminus\{a_1,\dots,a_r\}$ with $r\ge0$ and the $a_i$ distinct.

(b) Put $g(x)=\prod_i(x-a_i)$ and $V=\{(x,z)\in\mathbf A^2:\ z\,g(x)=1\}$, a closed set. The projection $V\to U'$, $(x,z)\mapsto x$, is a morphism, and $U'\to V$, $x\mapsto(x,1/g(x))$ is a morphism (both coordinates are regular on $U'$); they are inverse to each other. So $Y\cong U'\cong V$ is affine.

(c) The polynomial $zg(x)-1$ is irreducible (it is of degree~$1$ in $z$ with coprime coefficients $g$ and $-1$) so $I(V)=(zg-1)$ and $A(Y)\cong k[x,z]/(zg-1)\cong k[x][1/g]$, the localization of the UFD $k[x]$ at the multiplicative set $\{g^n\}$. A localization of a UFD is a UFD (every nonzero element is a unit times a product of the primes of $k[x]$ not dividing $g$; these stay prime, and the others become units). Hence $A(Y)$ is a UFD. \hfill$\square$

\exercise{6.2}{(\emph{An elliptic curve.}) Let $Y$ be the curve $y^2=x^3-x$ in $\mathbf A^2$, and assume that the characteristic of $k$ is $\ne 2$. In this exercise we show that $Y$ is not a rational curve, hence $K(Y)$ is not a pure transcendental extension of $k$.
\begin{enumerate}[label=(\alph*),leftmargin=*]
\item Show that $Y$ is nonsingular, and deduce that $A=A(Y)=k[x,y]/(y^2-x^3+x)$ is an integrally closed domain.
\item Let $k[x]$ be the subring of $K=K(Y)$ generated by the image of $x$ in $A$. Show that $k[x]$ is a polynomial ring, and that $A$ is the integral closure of $k[x]$ in $K$.
\item Show that there is an automorphism $\sigma:A\to A$ which sends $y$ to $-y$ and leaves $x$ fixed. For $a\in A$ define the norm $N(a)=a\cdot\sigma(a)$. Show that $N(a)\in k[x]$, $N(1)=1$, and $N(ab)=N(a)N(b)$ for all $a,b\in A$.
\item Using the norm, show that the units of $A$ are precisely the nonzero elements of $k$. Show that $x$ and $y$ are irreducible elements of $A$. Show that $A$ is not a unique factorization domain.
\item Prove that $Y$ is not a rational curve (Ex.~6.1).
\end{enumerate}
(See II, 8.20.3 and III, Ex.~5.3 for other proofs of this important result.)}
\pf
Put $f=y^2-x^3+x$.

(a) $f$ is irreducible: it is monic of degree~$2$ in $y$, so a factorization would be $(y-a)(y+a)$ with $a\in k[x]$ (Gauss) and $a^2=x^3-x$, impossible since the degree of the right side is odd. Hence $(f)$ is prime, $I(Y)=(f)$, $A$ is a domain, and $Y$ is a curve. The partial derivatives are $f_x=1-3x^2$ and $f_y=2y$. At a singular point both vanish; $2y=0$ gives $y=0$ (as $\operatorname{char}k\ne2$), then $x^3-x=0$, so $x\in\{0,1,-1\}$, and $1-3x^2=0$ gives $3x^2=1$: for $x=0$ this says $0=1$, for $x=\pm1$ it says $3=1$, i.e.\ $2=0$; both are false. So $Y$ is nonsingular. By (5.1) every local ring $A_{\fm}$ ($\fm$ maximal) is a regular local ring of dimension~$1$, hence a DVR (6.2A) and in particular integrally closed in $K$. Now $A=\bigcap_{\fm}A_\fm$ inside $K$, and an element of $K$ integral over $A$ is integral over each $A_\fm$, hence lies in each $A_\fm$, hence in $A$. So $A$ is integrally closed.

(b) Division by the polynomial $f$, monic in $y$ of degree~$2$, shows that every element of $A$ can be written uniquely as $p(x)+q(x)y$ with $p,q\in k[x]$ (uniqueness: a nonzero multiple of $f$ has $y$-degree $\ge2$). Thus $A=k[x]\oplus k[x]y$, and in particular $k[x]\to A$ is injective (a nonzero polynomial in $x$ alone is not in $(f)$), so $k[x]$ is a polynomial ring in one variable. $A$ is integral over $k[x]$ (it is generated by $y$ with $y^2=x^3-x$), and $A$ is integrally closed in $K=\operatorname{Frac}A$ by (a). So the integral closure of $k[x]$ in $K$ contains $A$, and is contained in $A$ because anything integral over $k[x]$ is integral over $A$. Hence it equals $A$.

(c) The $k$-algebra map $k[x,y]\to A$, $x\mapsto x$, $y\mapsto-y$ kills $f$ (as $(-y)^2-x^3+x=f$), so induces $\sigma:A\to A$, and $\sigma^2=\mathrm{id}$, so $\sigma$ is an automorphism. For $a=p+qy$ we have $\sigma(a)=p-qy$ and
\[
N(a)=p^2-q^2y^2=p^2-(x^3-x)q^2\in k[x].
\]
Since $\sigma$ is a ring homomorphism of a commutative ring, $N(ab)=ab\sigma(a)\sigma(b)=N(a)N(b)$, and $N(1)=1$.

(d) \emph{Degrees of norms.} Let $a=p+qy\ne0$. If $q\ne0$, then $2\deg p$ is even while $\deg\big((x^3-x)q^2\big)=2\deg q+3$ is odd, so no cancellation occurs and $\deg N(a)=\max(2\deg p,\,2\deg q+3)\ge3$. If $q=0$, then $N(a)=p^2$ has even degree. Hence:
\begin{itemize}
\item $\deg N(a)\ne1$ for every $a\ne0$;
\item $N(a)\in k^\times$ iff $q=0$ and $p\in k^\times$, i.e.\ iff $a\in k^\times$.
\end{itemize}
\emph{Units.} Nonzero constants are units. If $a$ is a unit then $N(a)N(a^{-1})=1$, so $N(a)\in k[x]^\times=k^\times$ and hence $a\in k^\times$.

\emph{Irreducibility.} $x$ and $y$ are nonzero nonunits (not in $k$). Suppose $x=ab$. Then $N(a)N(b)=N(x)=x^2$, so $\deg N(a)+\deg N(b)=2$ and neither degree is~$1$; thus one of them is~$0$, say $N(a)\in k^\times$, and then $a\in k^\times$ is a unit. So $x$ is irreducible. Suppose $y=ab$. Then $N(a)N(b)=N(y)=-y^2=x-x^3$, of degree~$3$; the degrees are $(0,3),(3,0),(1,2),(2,1)$, and the last two are excluded, so again one factor is a unit. So $y$ is irreducible.

\emph{Not a UFD.} In $A$ we have $y^2=x(x^2-1)$, so $x\mid y^2$. But $x\nmid y$: if $y=x(p+qy)$ then comparing coefficients in $A=k[x]\oplus k[x]y$ gives $xq=1$ in $k[x]$, absurd. So $x$ is an irreducible element which is not prime, and $A$ is not a UFD.

(e) Suppose $Y$ were rational. $Y$ is nonsingular by (a), and $Y\not\cong\mathbf P^1$: $Y$ is affine, while an affine variety isomorphic to $\mathbf P^1$ would have $A=\mathcal O(\mathbf P^1)=k$ (Thm.~3.4(a)), i.e.\ be a point. By Ex.~6.1(c), $A(Y)$ would be a UFD, contradicting (d). So $Y$ is not rational. By Cor.~4.5, if $K(Y)\cong k(t)$ then $Y$ would be birational to $\mathbf P^1$; hence $K(Y)$ is not a pure transcendental extension of $k$. \hfill$\square$

\exercise{6.3}{Show by example that the result of (6.8) is false if either (a) $\dim X\ge2$, or (b) $Y$ is not projective.}
\pf
Recall (6.8): if $X$ is an abstract nonsingular curve, $P\in X$, $Y$ a projective variety and $\varphi:X-P\to Y$ a morphism, then $\varphi$ extends to a morphism $X\to Y$.

(a) Let $X=\mathbf A^2$ (nonsingular, $\dim 2$), $P=(0,0)$, $Y=\mathbf P^1$, and $\varphi:X-P\to\mathbf P^1$, $(x,y)\mapsto(x:y)$, a morphism since $x,y$ do not vanish simultaneously. Suppose $\bar\varphi:X\to\mathbf P^1$ is a morphism extending $\varphi$. For a line $\ell\subseteq\mathbf A^2$ through $P$, $\varphi$ is constant on $\ell\setminus P$, with value $c_\ell$, and $c_\ell\ne c_{\ell'}$ for $\ell\ne\ell'$. The set $\ell\cap\bar\varphi^{-1}(c_\ell)$ is closed in $\ell\cong\mathbf A^1$ and contains the infinite set $\ell\setminus P$, so it is all of $\ell$. Hence $\bar\varphi(P)=c_\ell$ for every $\ell$, contradicting $c_\ell\neq c_{\ell'}$.

(b) Let $X=\mathbf A^1$ (an abstract nonsingular curve by (6.7)), $P=0$, and $Y=\mathbf A^1$, which is quasi-projective but not projective. Let $\varphi:X-P\to Y$, $\varphi(x)=1/x$, a morphism. An extension $\bar\varphi:\mathbf A^1\to\mathbf A^1$ would be a polynomial $g$ with $xg(x)=1$ for all $x\ne0$; the polynomial $xg(x)-1$ would have infinitely many zeros, hence be zero, but its value at $0$ is $-1$. So no extension exists. \hfill$\square$

\exercise{6.4}{Let $Y$ be a nonsingular projective curve. Show that every nonconstant rational function $f$ on $Y$ defines a surjective morphism $\varphi:Y\to\mathbf P^1$, and that for every $P\in\mathbf P^1$, $\varphi^{-1}(P)$ is a finite set of points.}
\pf
Let $K=K(Y)$ and $f\in K\setminus k$.

\emph{The morphism.} Let $U\subseteq Y$ be the open set where $f$ is regular, and $\varphi_U:U\to\mathbf A^1\subset\mathbf P^1$ the morphism it defines. $Y\setminus U$ is finite ($Y$ is an irreducible curve). By (6.7) $Y$ is an abstract nonsingular curve, so applying (6.8) to $U\cup\{P_i\}\supseteq U$ for each $P_i\in Y\setminus U$ and gluing the extensions (they agree on $U$) we get a morphism $\varphi:Y\to\mathbf P^1$ with $\varphi|_U=f$. Then $\varphi^*:K(\mathbf P^1)=k(t)\to K$ sends $t\mapsto f$. Since $f\notin k$ and $k$ is algebraically closed, $f$ is transcendental over $k$, so $\varphi^*$ is injective, and $K$ is a finitely generated extension of $k$ of transcendence degree $1$, hence a finite extension of $k(f)$.

\emph{Surjectivity.} Let $Q\in\mathbf P^1$ and let $O=\mathcal O_{Q,\mathbf P^1}\subseteq k(t)$, a DVR (5.1, 6.2A). It is a Dedekind domain, so by (6.3A) the integral closure $B$ of $O$ in the finite extension $K\supseteq k(f)=k(t)$ is a Dedekind domain; it is not a field (otherwise $O$ would be a field). Let $\fn$ be a maximal ideal of $B$; then $R=B_\fn$ is a DVR of $K$ containing $k$, and $\fn\cap O$ is a maximal ideal of $O$ (integral extension), so $\fm_R\cap O=\fm_Q$. By Lemma~A, $R=\mathcal O_{P,Y}$ for some $P\in Y$. We claim $\varphi(P)=Q$. Put $R'=R\cap k(t)$. It contains $O$ and also $\mathcal O_{\varphi(P)}$ (since $\varphi^*:\mathcal O_{\varphi(P)}\to\mathcal O_{P,Y}$ is the inclusion inside $K$). $R'\ne k(t)$ because $\fm_R\cap k(t)\supseteq\fm_Q\ne0$ consists of nonunits of $R$. By the overring remark $R'=O$ and $R'=\mathcal O_{\varphi(P)}$. So $\mathcal O_Q=\mathcal O_{\varphi(P)}$ and $\varphi(P)=Q$ by Lemma~6.4. Thus $\varphi$ is surjective.

\emph{Finiteness of fibres.} $\varphi$ is not constant ($f\notin k$), so for each $P\in\mathbf P^1$ the closed set $\varphi^{-1}(P)$ is a proper closed subset of the irreducible curve $Y$, hence finite. \hfill$\square$

\exercise{6.5}{Let $X$ be a nonsingular projective curve. Suppose that $X$ is a (locally closed) subvariety of a variety $Y$ (Ex.~3.10). Show that $X$ is in fact a closed subset of $Y$. (See II, Ex.~4.4 for a generalization.)}
\pf
Let $Z$ be the closure of $X$ in $Y$. Since $X$ is locally closed, $X$ is open in $Z$; $Z$ is a variety, with $\dim Z=\dim X=1$ and $K(Z)=K(X)=K$. Suppose $P\in Z\setminus X$.

Choose an affine open neighbourhood $V$ of $P$ in $Z$, with coordinate ring $A=A(V)$, a finitely generated $k$-domain of dimension~$1$ with fraction field $K$, and let $\fm\subset A$ be the maximal ideal of $P$, so $\mathcal O_{P,Z}=A_\fm$. Let $A'$ be the integral closure of $A$ in $K$. By (3.9A) $A'$ is a finite $A$-module, hence a noetherian integrally closed domain of dimension~$1$, i.e.\ a Dedekind domain. By lying over there is a maximal ideal $\fn\subset A'$ with $\fn\cap A=\fm$. Then $R=A'_\fn$ is a DVR of $K/k$ containing $A_\fm=\mathcal O_{P,Z}$.

By Lemma~A applied to the nonsingular projective curve $X$, $R=\mathcal O_{Q,X}$ for some $Q\in X$. As $X$ is open in $Z$, $\mathcal O_{Q,X}=\mathcal O_{Q,Z}$. So $\mathcal O_{P,Z}\subseteq\mathcal O_{Q,Z}$ as subrings of $K$, and Lemma~6.4 (applied to the quasi-projective variety $Z$) gives $P=Q\in X$, a contradiction. Hence $Z=X$ and $X$ is closed in $Y$. \hfill$\square$

\exercise{6.6}{(\emph{Automorphisms of $\mathbf P^1$.}) Think of $\mathbf P^1$ as $\mathbf A^1\cup\{\infty\}$. Then we define a \emph{fractional linear transformation} of $\mathbf P^1$ by sending $x\mapsto(ax+b)/(cx+d)$, for $a,b,c,d\in k$, $ad-bc\ne0$.
\begin{enumerate}[label=(\alph*),leftmargin=*]
\item Show that a fractional linear transformation induces an automorphism of $\mathbf P^1$ (i.e.\ an isomorphism of $\mathbf P^1$ with itself). We denote the group of all these by $PGL(1)$.
\item Let $\operatorname{Aut}\mathbf P^1$ denote the group of automorphisms of $\mathbf P^1$. Show that $\operatorname{Aut}\mathbf P^1\cong\operatorname{Aut}_k k(x)$, the group of $k$-automorphisms of the field $k(x)$.
\item Now show that every automorphism of $k(x)$ is a fractional linear transformation, and deduce that $PGL(1)\to\operatorname{Aut}\mathbf P^1$ is an isomorphism.
\end{enumerate}
(Note: we will see later (II, 7.1.1) that a similar result holds for $\mathbf P^n$: every automorphism is given by a linear change of the homogeneous coordinates.)}
\pf
Use homogeneous coordinates $(u:v)$ on $\mathbf P^1$, with $x=u/v$ on $\mathbf A^1=\{v\ne0\}$ and $\infty=(1:0)$.

(a) For $M=\begin{pmatrix}a&b\\c&d\end{pmatrix}\in GL_2(k)$ define $\varphi_M(u:v)=(au+bv:cu+dv)$. The two linear forms have no common zero on $k^2\setminus0$ because $M$ is invertible, so $\varphi_M$ is a morphism $\mathbf P^1\to\mathbf P^1$ (given by homogeneous polynomials of the same degree without common zero). Clearly $\varphi_M\circ\varphi_{M'}=\varphi_{MM'}$ and $\varphi_I=\mathrm{id}$, so $\varphi_M$ is an automorphism with inverse $\varphi_{M^{-1}}$. On $\mathbf A^1$ it is $x\mapsto(ax+b)/(cx+d)$ (with value $\infty$ where $cx+d=0$, and $\varphi_M(\infty)=a/c$, resp.\ $\infty$ if $c=0$): this is exactly the fractional linear transformation. So each fractional linear transformation is an automorphism, and $PGL(1)$ is the image of the homomorphism $GL_2(k)\to\operatorname{Aut}\mathbf P^1$, $M\mapsto\varphi_M$. Its kernel consists of the scalar matrices: if $\varphi_M=\mathrm{id}$ then $M$ fixes the points $(1:0),(0:1),(1:1)$, which forces $c=0$, $b=0$, $a=d$. So $PGL(1)\cong GL_2(k)/k^\times$.

(b) By Cor.~6.12 the functor $Y\mapsto K(Y)$ is an equivalence (contravariant) between nonsingular projective curves with dominant morphisms and function fields of dimension~$1$ with $k$-homomorphisms. In particular $\varphi\mapsto\varphi^*$ gives a bijection $\operatorname{Hom}_{\rm dom}(\mathbf P^1,\mathbf P^1)\to\operatorname{Hom}_k(k(x),k(x))$ with $(\varphi\psi)^*=\psi^*\varphi^*$, and it restricts to a bijection on invertible elements, i.e.\ $\operatorname{Aut}\mathbf P^1\to\operatorname{Aut}_kk(x)$. This is a bijective anti-homomorphism; composing with inversion, $\varphi\mapsto(\varphi^*)^{-1}$ is a group isomorphism $\operatorname{Aut}\mathbf P^1\cong\operatorname{Aut}_kk(x)$.

(c) Let $\sigma\in\operatorname{Aut}_kk(x)$ and write $\sigma(x)=u=f/g$ with $f,g\in k[x]$ coprime; $u\notin k$ since $\sigma$ is injective.

\emph{Claim:} $[k(x):k(u)]=\max(\deg f,\deg g)$. Indeed $u$ is transcendental over $k$, so $k[u,T]$ is a polynomial ring, and $x$ is a root of $F(T)=f(T)-ug(T)\in k[u][T]$, whose $T$-degree is $\max(\deg f,\deg g)\ge1$. $F$ is irreducible in $k[u,T]$: it has degree~$1$ in $u$, so in a factorization one factor lies in $k[T]$ and divides both $f$ and $g$, hence is a constant. By Gauss's lemma $F$ is irreducible in $k(u)[T]$, so $[k(x):k(u)]=\deg_TF$.

Now $\sigma$ is surjective, so $k(x)=\sigma(k(x))=k(\sigma(x))=k(u)$, and therefore $\max(\deg f,\deg g)=1$. Thus $f=ax+b$, $g=cx+d$ with $a,b,c,d\in k$; and $ad-bc\ne0$, for otherwise $(a,b)$ and $(c,d)$ would be proportional and $f/g\in k$. So $\sigma(x)=(ax+b)/(cx+d)$.

For $M$ as above, $\varphi_M^*(x)=x\circ\varphi_M=(ax+b)/(cx+d)=\sigma(x)$, and a $k$-automorphism of $k(x)$ is determined by the image of $x$, so $\sigma=\varphi_M^*$. Hence every element of $\operatorname{Aut}_kk(x)$ is $\varphi_M^*$ for some $M$. Given an arbitrary $\varphi\in\operatorname{Aut}\mathbf P^1$, $\varphi^*=\varphi_M^*$ for some $M$ and, by injectivity in (b), $\varphi=\varphi_M\in PGL(1)$. So $PGL(1)=\operatorname{Aut}\mathbf P^1$, i.e.\ the injective map $PGL(1)\to\operatorname{Aut}\mathbf P^1$ (inclusion of the group of fractional linear transformations) is an isomorphism. \hfill$\square$

\exercise{6.7}{Let $P_1,\dots,P_r,Q_1,\dots,Q_s$ be distinct points of $\mathbf A^1$. If $\mathbf A^1-\{P_1,\dots,P_r\}$ is isomorphic to $\mathbf A^1-\{Q_1,\dots,Q_s\}$, show that $r=s$. Is the converse true? (Cf.\ Ex.~3.1.)}
\pf
Write $\Sigma=\{P_1,\dots,P_r,\infty\}$, $\Sigma'=\{Q_1,\dots,Q_s,\infty\}\subset\mathbf P^1$, $U=\mathbf P^1\setminus\Sigma$, $V=\mathbf P^1\setminus\Sigma'$.

\emph{Claim.} Every isomorphism $\varphi:U\to V$ is the restriction of an automorphism $\Phi$ of $\mathbf P^1$; in particular $\Phi(\Sigma)=\Sigma'$.

\emph{Proof.} $U,V$ are open in $\mathbf P^1$, hence abstract nonsingular curves, and $\mathbf P^1$ is projective. Apply (6.8) to $\varphi:U\to V\subseteq\mathbf P^1$, one point of the finite set $\Sigma$ at a time, and glue: we get a morphism $\Phi:\mathbf P^1\to\mathbf P^1$ extending $\varphi$. Likewise $\varphi^{-1}$ extends to $\Psi$. $\Psi\circ\Phi$ is a morphism $\mathbf P^1\to\mathbf P^1$ which is the identity on the dense open set $U$, and $\mathrm{id}$ is another one, so by the uniqueness in (6.8), applied to one point of $\Sigma$ at a time, $\Psi\circ\Phi=\mathrm{id}$ on all of $\mathbf P^1$. So $\Psi\Phi=\mathrm{id}$, and similarly $\Phi\Psi=\mathrm{id}$. Thus $\Phi$ is a bijection with $\Phi(U)=V$, so $\Phi(\Sigma)=\mathbf P^1\setminus V=\Sigma'$. \hfill$\square$

Hence an isomorphism $U\cong V$ forces $|\Sigma|=|\Sigma'|$, i.e.\ $r+1=s+1$, so $r=s$.

\emph{The converse.} By Ex.~6.6, $\Phi\in PGL(1)$, and conversely a Möbius transformation $g$ with $g(\Sigma)=\Sigma'$ restricts to an isomorphism $U\to V$. Recall that $PGL(1)$ is sharply $3$-transitive on $\mathbf P^1$: for distinct $a,b,c\in\mathbf P^1$ there is exactly one element sending $(a,b,c)\mapsto(0,1,\infty)$ (it is $z\mapsto\frac{(z-a)(b-c)}{(z-c)(b-a)}$, with the obvious modification if one of $a,b,c$ is $\infty$; uniqueness because a Möbius transformation fixing three points is the identity, as $cz^2+(d-a)z-b=0$ has at most two roots unless $c=b=0,\ a=d$).

\begin{itemize}
\item If $r=s\le2$, then $|\Sigma|=|\Sigma'|\le3$, so $\Sigma,\Sigma'$ can be enlarged to $3$-element sets, and a Möbius transformation sends the enlargement of $\Sigma$ onto that of $\Sigma'$ compatibly (choose the bijection so that $\Sigma\mapsto\Sigma'$). Hence $\mathbf A^1-\{P_i\}\cong\mathbf A^1-\{Q_j\}$: the converse holds for $r=s\le 2$.
\item If $r=s\ge3$ the converse is \emph{false}. Fix $\Sigma$ with $|\Sigma|=r+1\ge4$. If $g\in PGL(1)$ and $g(\Sigma)=\Sigma'\supseteq\{0,1,\infty\}$, then $g$ maps some ordered triple of distinct points of $\Sigma$ to $(0,1,\infty)$; there are at most $(r+1)r(r-1)$ such triples, each determines $g$ uniquely, so only finitely many sets $\Sigma'\supseteq\{0,1,\infty\}$ are of the form $g(\Sigma)$. But there are infinitely many $(r+1)$-element subsets of $\mathbf P^1$ containing $\{0,1,\infty\}$ ($k$ is infinite). Choose such a $\Sigma'$ not of the form $g(\Sigma)$ and let $Q_j$ be the finite points of $\Sigma'$, so $s=r$. By the Claim and Ex.~6.6, $\mathbf A^1-\{P_i\}\not\cong\mathbf A^1-\{Q_j\}$.
\end{itemize}
So the converse is true exactly for $r=s\le2$ and false for every $r=s\ge3$. \hfill$\square$

\section*{7. Intersections in projective space}

Throughout, $k$ is algebraically closed (so infinite), $S=k[x_0,\dots,x_n]$, $S(Y)=S/I(Y)$ for $Y\subseteq\mathbf P^n$, and numbers refer to Hartshorne, Chapter~I. For a finitely generated graded module $M$, $h_M(l)=\dim_kM_l$ is its Hilbert function and $P_M$ its Hilbert polynomial (Thm.~7.5); $h_Y,P_Y$ are those of $S(Y)$, and for $\dim Y=r$, $\deg P_Y=r$ with leading coefficient $\deg Y/r!$. Two facts used repeatedly: (i) a linear change of coordinates preserves all the notions involved; (ii) if $Y=\mathbf P(V)$ is a linear subspace of dimension $r$ then $S(Y)\cong k[z_0,\dots,z_r]$ and $h_Y(l)=\binom{l+r}{r}$, so $\deg Y=1$.

\exercise{7.1}{
\begin{enumerate}[label=(\alph*),leftmargin=*]
\item Find the degree of the $d$-uple embedding of $\mathbf P^n$ in $\mathbf P^N$ (Ex.~2.12). [Answer: $d^n$.]
\item Find the degree of the Segre embedding of $\mathbf P^r\times\mathbf P^s$ in $\mathbf P^N$ (Ex.~2.14). [Answer: $\binom{r+s}{r}$.]
\end{enumerate}}
\pf
(a) Let $\psi:k[y_\alpha]\to S$ send $y_\alpha\mapsto x^\alpha$, where $\alpha$ runs over the exponents of degree $d$ ($N+1=\binom{n+d}{n}$ of them). The image $Y=v_d(\mathbf P^n)$ of the $d$-uple embedding is $Z(\ker\psi)$; and $I(Y)=\ker\psi$, because a polynomial $F(y)$ vanishes on $Y$ (equivalently on all points $(x^\alpha)$, $x\in k^{n+1}$) iff $F(x^\alpha)=0$ as a polynomial in $x$ ($k$ is infinite), and $\ker\psi$ is a prime ideal. So $S(Y)\cong\operatorname{Im}\psi$, graded so that $y_\alpha$ has degree~$1$. The image is spanned by products of $l$ monomials of degree $d$, i.e.\ by all monomials of degree $ld$ (each such monomial is a product of $l$ monomials of degree $d$). Hence $S(Y)_l\cong S_{ld}$ and
\[
h_Y(l)=\binom{ld+n}{n}=\frac{d^n}{n!}\,l^n+(\text{lower terms}).
\]
Since $\dim Y=n$, $\deg Y=n!\cdot d^n/n!=d^n$.

(b) Let $S_1=k[x_0,\dots,x_r]$, $T=k[y_0,\dots,y_s]$ and $\psi:k[z_{ij}]\to S_1\otimes T=k[x,y]$, $z_{ij}\mapsto x_iy_j$. As above $I(Y)=\ker\psi$ for $Y=$ the Segre image (a polynomial vanishing at all points $(x_iy_j)$ is killed by $\psi$ since $k$ is infinite), so $S(Y)\cong\operatorname{Im}\psi$, which in degree $l$ is spanned by the products of $l$ of the $x_iy_j$, i.e.\ by all monomials $x^\alpha y^\beta$ with $|\alpha|=|\beta|=l$. These monomials are linearly independent, so
\[
h_Y(l)=\binom{l+r}{r}\binom{l+s}{s}=\frac{l^{r+s}}{r!\,s!}+(\text{lower terms}).
\]
As $\dim Y=r+s$, $\deg Y=(r+s)!/(r!\,s!)=\binom{r+s}{r}$. \hfill$\square$

\exercise{7.2}{Let $Y$ be a variety of dimension $r$ in $\mathbf P^n$, with Hilbert polynomial $P_Y$. We define the \emph{arithmetical genus} of $Y$ to be $p_a(Y)=(-1)^r(P_Y(0)-1)$. This is an important invariant which (as we will see later in III, Ex.~5.3) is independent of the projective embedding of $Y$.
\begin{enumerate}[label=(\alph*),leftmargin=*]
\item Show that $p_a(\mathbf P^n)=0$.
\item If $Y$ is a plane curve of degree $d$, show that $p_a(Y)=\frac12(d-1)(d-2)$.
\item More generally, if $H$ is a hypersurface of degree $d$ in $\mathbf P^n$, then $p_a(H)=\binom{d-1}{n}$.
\item If $Y$ is a complete intersection (Ex.~2.17) of surfaces of degrees $a,b$ in $\mathbf P^3$, then $p_a(Y)=\frac12ab(a+b-4)+1$.
\item Let $Y^r\subseteq\mathbf P^n$, $Z^s\subseteq\mathbf P^m$ be projective varieties, and embed $Y\times Z\subseteq\mathbf P^n\times\mathbf P^m\to\mathbf P^N$ by the Segre embedding. Show that
\[
p_a(Y\times Z)=p_a(Y)p_a(Z)+(-1)^sp_a(Y)+(-1)^rp_a(Z).
\]
\end{enumerate}}
\pf
For a polynomial of the form $\binom{z+m}{m}=\frac1{m!}(z+1)\cdots(z+m)$ we have value $1$ at $z=0$.

(a) $P_{\mathbf P^n}(z)=\binom{z+n}{n}$ (proof of 7.6(c)), so $P(0)=1$ and $p_a=(-1)^n(1-1)=0$.

(c) As in the proof of 7.6(d), $I(H)=(f)$ with $\deg f=d$ and $P_H(z)=\binom{z+n}{n}-\binom{z-d+n}{n}$ (polynomial identity; $\dim H=n-1$). At $z=0$ the second binomial is
\[
\binom{n-d}{n}=\frac{(n-d)(n-d-1)\cdots(1-d)}{n!}=(-1)^n\frac{(d-1)(d-2)\cdots(d-n)}{n!}=(-1)^n\binom{d-1}{n},
\]
where $\binom{d-1}{n}=0$ if $d\le n$ (the numerator then contains the factor $0$). So $P_H(0)=1-(-1)^n\binom{d-1}n$ and
\[
p_a(H)=(-1)^{n-1}\big(-(-1)^n\binom{d-1}{n}\big)=\binom{d-1}{n}.
\]

(b) is (c) with $n=2$: $\binom{d-1}{2}=\frac12(d-1)(d-2)$.

(d) Let $I(Y)=(f,g)$ with $\deg f=a$, $\deg g=b$, $S=k[x_0,\dots,x_3]$, $\dim Y=1$. Then $f,g$ have no common factor (a common nonconstant factor $h$ would give $Y\supseteq Z(h)$, of dimension $\ge2$). The sequence
\[
0\to S(-a-b)\xrightarrow{\ w\mapsto(gw,-fw)\ }S(-a)\oplus S(-b)\xrightarrow{\ (u,v)\mapsto uf+vg\ }S\to S/(f,g)\to0
\]
is exact: the first map is injective; the composite is zero; if $uf+vg=0$ then $f\mid vg$, so $f\mid v$, say $v=-fw$, and then $u=gw$. Hence
\[
P_Y(z)=\binom{z+3}{3}-\binom{z-a+3}{3}-\binom{z-b+3}{3}+\binom{z-a-b+3}{3}.
\]
Using $\binom{3-c}{3}=\frac{(3-c)(2-c)(1-c)}{6}=-\frac{(c-1)(c-2)(c-3)}{6}$ and writing $\rho(c)=(c-1)(c-2)(c-3)=c^3-6c^2+11c-6$,
\[
P_Y(0)=1+\tfrac16\big(\rho(a)+\rho(b)-\rho(a+b)\big).
\]
Now $a^3+b^3-(a+b)^3=-3ab(a+b)$, $a^2+b^2-(a+b)^2=-2ab$, and the linear terms cancel, so $\rho(a)+\rho(b)-\rho(a+b)=-3ab(a+b)+12ab-6$. Hence $P_Y(0)=1-\frac12ab(a+b)+2ab-1=-\frac12ab(a+b-4)$ and $p_a(Y)=(-1)^1(P_Y(0)-1)=\frac12ab(a+b-4)+1$.

(e) Let $S(Y)=S/I(Y)$, $S(Z)=S'/I(Z)$ be the homogeneous coordinate rings. By Ex.~3.15, $S(Y)\otimes_kS(Z)$ is the coordinate ring of $C(Y)\times C(Z)$ (product of the affine cones). The Segre embedding sends $((x),(y))\mapsto(x_iy_j)$, and the affine cone of $Y\times Z\subseteq\mathbf P^N$ is the image of $C(Y)\times C(Z)$ under $\mu:(x,y)\mapsto(x_iy_j)$. Let $\theta:k[z_{ij}]\to S(Y)\otimes S(Z)$, $z_{ij}\mapsto\bar x_i\otimes\bar y_j$. Then $F\in\ker\theta$ iff $F\circ\mu$ vanishes on $C(Y)\times C(Z)$ iff $F$ vanishes on the cone of $Y\times Z$, i.e.\ $\ker\theta=I(Y\times Z)$. So $S(Y\times Z)$ is the subring of $S(Y)\otimes S(Z)$ generated by the $\bar x_i\otimes\bar y_j$, which in degree $l$ is $S(Y)_l\otimes S(Z)_l$ (spanned by products of $l$ such elements). Therefore $h_{Y\times Z}(l)=h_Y(l)h_Z(l)$, and $P_{Y\times Z}=P_YP_Z$, $\dim(Y\times Z)=r+s$. Write $P_Y(0)=1+(-1)^rp_a(Y)$, $P_Z(0)=1+(-1)^sp_a(Z)$. Then
\begin{align*}
p_a(Y\times Z)&=(-1)^{r+s}\big(P_Y(0)P_Z(0)-1\big)\\
&=(-1)^{r+s}\big((-1)^rp_a(Y)+(-1)^sp_a(Z)+(-1)^{r+s}p_a(Y)p_a(Z)\big)\\
&=(-1)^sp_a(Y)+(-1)^rp_a(Z)+p_a(Y)p_a(Z).\qquad\square
\end{align*}

\exercise{7.3}{(\emph{The dual curve.}) Let $Y\subseteq\mathbf P^2$ be a curve. We regard the set of lines in $\mathbf P^2$ as another projective space $(\mathbf P^2)^*$, by taking $(a_0,a_1,a_2)$ as homogeneous coordinates of the line $L:a_0x_0+a_1x_1+a_2x_2=0$. For each nonsingular point $P\in Y$, show that there is a unique line $T_P(Y)$ whose intersection multiplicity with $Y$ at $P$ is $>1$. This is the \emph{tangent line} to $Y$ at $P$. Show that the mapping $P\mapsto T_P(Y)$ defines a morphism of $\operatorname{Reg}Y$ (the set of nonsingular points of $Y$) into $(\mathbf P^2)^*$. The closure of the image of this morphism is called the \emph{dual curve} $Y^*\subseteq(\mathbf P^2)^*$.}
\pf
Let $Y=Z(f)$ with $f\in S=k[x_0,x_1,x_2]$ irreducible of degree $d$, and $f_i=\partial f/\partial x_i$. For a line $L=Z(l)\not\supseteq Y$ the intersection multiplicity at $P$ is $i(Y,L;P)=\mu_{\fp}(S/(f,l))$, $\fp=I(P)$ (Thm.~7.7). If $d=1$ then $Y$ is itself a line, and the only line containing $Y$ is $Y$; the multiplicity is only defined for $L\ne Y$, it equals $1$ at the single point $L\cap Y$, and the natural tangent line is $T_P(Y)=Y$ (the constant map). So assume $d\ge2$; then no line lies in $Y$ ($f$ irreducible of degree $\ge2$ is not divisible by a linear form), and $i(Y,L;P)$ is defined for all lines $L$.

\emph{Step 1: computing $i(Y,L;P)$.} Let $P=[p]\in Y\cap L$ and pick $q\in k^3$ with $L=\{[sp+tq]\}$ ($p,q$ independent). The ring map $\rho:S\to k[s,t]$, $x_i\mapsto sp_i+tq_i$, is surjective (the forms $sp_i+tq_i$ span $ks+kt$) and kills $l$, so $\ker\rho=(l)$ (the kernel is a homogeneous ideal containing the linear form $l$, and $S/(l)\cong k[s,t]$ has the same Hilbert function as $k[s,t]$). Under this identification $S/(f,l)\cong k[s,t]/(g)$ with $g(s,t)=f(sp+tq)$ a binary form of degree $d$, and $g\ne0$ since $L\not\subseteq Y$. The prime $\fp=I(P)$ corresponds to $(t)$ (the point $P$ is $(s,t)=(1,0)$). The localization $k[s,t]_{(t)}$ is a DVR with uniformizer $t$, so if $g=t^mh$ with $t\nmid h$ then $h$ is a unit there, and
\[
i(Y,L;P)=\operatorname{length}\big(k[s,t]_{(t)}/(t^m)\big)=m=\operatorname{ord}_t g.
\]
(The length of $M_\fp$ over $S_\fp$ equals that over the quotient ring, since the module structure factors through it.)

\emph{Step 2: the tangent line.} Expanding in $t$ (a polynomial identity valid in every characteristic),
\[
g(s,t)=f(sp+tq)=f(sp)+t\sum_i f_i(sp)q_i+t^2(\cdots)=s^df(p)+t\,s^{d-1}\sum_if_i(p)q_i+t^2(\cdots),
\]
because $f_i$ is homogeneous of degree $d-1$. Since $P\in Y$, $f(p)=0$, and so
\[
i(Y,L;P)>1\iff\sum_if_i(p)q_i=0 \iff \nabla f(p)\cdot q=0.
\]
Moreover $\nabla f(p)\cdot p=d\,f(p)=0$ (Euler's formula, valid in any characteristic). Thus $\nabla f(p)\cdot v=0$ for all $v\in\operatorname{span}(p,q)=\hat L$, the $2$-dimensional subspace $\{a\cdot v=0\}$ if $L=\{a_0x_0+a_1x_1+a_2x_2=0\}$. So $i(Y,L;P)>1$ iff $\nabla f(p)\in\hat L^{\perp}=k\cdot a$, i.e.\ iff $(a_0,a_1,a_2)$ is proportional to $\nabla f(p)$ when $\nabla f(p)\ne0$; if $\nabla f(p)=0$ then every line through $P$ has multiplicity $>1$.

\emph{Step 3: nonsingular $\iff\nabla f(p)\ne0$.} Say $p_i\ne0$ and use the affine chart $x_i=1$ with $F(y)=f|_{x_i=1}$ (defining $Y\cap U_i$, with $I=(F)$). Then $\partial F/\partial y_j=f_j$ evaluated at the point, for $j\ne i$. By Euler, $\sum_jp_jf_j(p)=0$, so if $f_j(p)=0$ for all $j\ne i$ then $f_i(p)=0$. Hence $P$ is nonsingular on the curve $Y\cap U_i$ (Jacobian criterion, rank $=n-r=1$) iff $\nabla f(p)\ne0$.

Consequently, for a nonsingular point $P=[p]$ there is exactly one line with $i(Y,L;P)>1$, namely
\[
T_P(Y)=(f_0(P):f_1(P):f_2(P))\in(\mathbf P^2)^*.
\]
\emph{Morphism.} $\operatorname{Reg}Y$ is open in $Y$ (Thm.~5.3), and $P\mapsto(f_0(P):f_1(P):f_2(P))$ is given by three homogeneous polynomials of the same degree $d-1$ without common zero on $\operatorname{Reg}Y$; so it is a morphism $\operatorname{Reg}Y\to(\mathbf P^2)^*$ (on the open set $\{f_j\neq0\}$ the coordinates $f_i/f_j$ are regular). \hfill$\square$

\exercise{7.4}{Given a curve $Y$ of degree $d$ in $\mathbf P^2$, show that there is a nonempty open subset $U$ of $(\mathbf P^2)^*$ in its Zariski topology such that for each $L\in U$, $L$ meets $Y$ in exactly $d$ points. [\emph{Hint}: Show that the set of lines in $(\mathbf P^2)^*$ which are either tangent to $Y$ or pass through a singular point of $Y$ is contained in a proper closed subset.] This result shows that we could have defined the degree of $Y$ to be the number $d$ such that almost all lines in $\mathbf P^2$ meet $Y$ in $d$ points, where ``almost all'' refers to a nonempty open subset of the set of lines, when this set is identified with the dual projective space $(\mathbf P^2)^*$.}
\pf
If $d=1$, $Y$ is a line and $U=(\mathbf P^2)^*\setminus\{Y\}$ (open) works: every other line meets $Y$ in exactly $1=d$ point. Let $d\ge2$.

The set $\operatorname{Sing}Y$ is a proper closed subset of the curve $Y$ (Thm.~5.3), hence finite, say $\{Q_1,\dots,Q_m\}$. For a point $Q=[q]\in\mathbf P^2$, the set $\Lambda_Q=\{a\in(\mathbf P^2)^*:a\cdot q=0\}$ of lines through $Q$ is a line in $(\mathbf P^2)^*$, hence closed of dimension~$1$. By Ex.~7.3, $\varphi:\operatorname{Reg}Y\to(\mathbf P^2)^*$, $P\mapsto T_P(Y)$, is a morphism from an irreducible variety of dimension~$1$, so its image closure $Y^*$ is an irreducible closed subset with $k(Y^*)\hookrightarrow K(Y)$, whence $\dim Y^*\le1$. Put
\[
\Gamma=Y^*\cup\Lambda_{Q_1}\cup\dots\cup\Lambda_{Q_m},
\]
a closed subset of $(\mathbf P^2)^*$ of dimension $\le1<2$, hence proper, and $U=(\mathbf P^2)^*\setminus\Gamma\ne\emptyset$ is open.

Let $L\in U$. Since $d\ge2$, $L\not\supseteq Y$, so $L\cap Y$ is finite and nonempty (Thm.~7.2). By Thm.~7.7 applied to $Y$ and the hypersurface $L$ ($\deg L=1$, and the components of $Y\cap L$ are points of degree~$1$),
\[
\sum_{P\in L\cap Y}i(Y,L;P)=d.
\]
Every $P\in L\cap Y$ is nonsingular on $Y$ ($L\notin\Lambda_{Q_j}$), and $L\ne T_P(Y)$ (else $L\in Y^*$); by Ex.~7.3 therefore $i(Y,L;P)=1$ (it is $\ge1$ since $P$ is a minimal prime of $S/(I_Y+I_L)$, and not $>1$). Hence $L\cap Y$ consists of exactly $d$ points. \hfill$\square$

\exercise{7.5}{
\begin{enumerate}[label=(\alph*),leftmargin=*]
\item Show that an irreducible curve $Y$ of degree $d>1$ in $\mathbf P^2$ cannot have a point of multiplicity $\ge d$ (Ex.~5.3).
\item If $Y$ is an irreducible curve of degree $d>1$ having a point of multiplicity $d-1$, then $Y$ is a rational curve (Ex.~6.1).
\end{enumerate}}
\pf
Let $Y=Z(f)$, $f\in k[x_0,x_1,x_2]$ irreducible homogeneous of degree $d>1$ and let $P\in Y$. Recall (Ex.~5.3) the multiplicity $\mu_P(Y)$: in an affine chart containing $P$, translated so that $P$ is the origin, with equation $F=F_\mu+F_{\mu+1}+\dots$ ($F_j$ homogeneous of degree $j$, $F_\mu\ne0$), $\mu_P(Y)=\mu$; equivalently, the largest $\mu$ with $F\in\fm_P^\mu$, which is invariant under isomorphisms of local rings, in particular under a linear change of coordinates in $\mathbf P^2$. So we may assume $P=(1:0:0)$. In the chart $x_0=1$ with coordinates $x=x_1/x_0$, $y=x_2/x_0$, $F(x,y)=f(1,x,y)$, and writing $f=\sum_{j}x_0^{d-j}f_j(x_1,x_2)$ with $f_j$ binary forms of degree $j$ we get $F=\sum_jf_j(x,y)$, so $F_j=f_j(x,y)$ and $f_j=0$ for $j<\mu=\mu_P(Y)$. Also $f_0=0$ as $P\in Y$.

(a) If $\mu\ge d$ then, since $j\le d$, $f=f_d(x_1,x_2)$ is a binary form of degree $d$, which factors into $d$ linear forms ($k$ algebraically closed). Irreducibility of $f$ forces $d=1$, contrary to $d>1$.

(b) Now $\mu=d-1$: $f=x_0f_{d-1}(x_1,x_2)+f_d(x_1,x_2)$ and $F=f_{d-1}(x,y)+f_d(x,y)$. Both $f_{d-1}$ and $f_d$ are nonzero ($f_d=0$ would make $x_0\mid f$, contradicting irreducibility for $d>1$). In the function field $K(Y)=k(x,y)$ ($Y\cap\{x_0\ne0\}$ is a nonempty open subset) put $t=y/x$; here $x\ne0$ in $K(Y)$ since $Y\not\subseteq\{x_1=0\}$ (a line), and $t$ is not in $k$ for the same reason ($t=c$ would put $Y$ in the line $x_2=cx_1$), so $t$ is transcendental over $k$. Substituting $y=tx$,
\[
0=F(x,tx)=x^{d-1}f_{d-1}(1,t)+x^df_d(1,t)=x^{d-1}\big(f_{d-1}(1,t)+x\,f_d(1,t)\big).
\]
Since $x\ne0$ we get $f_{d-1}(1,t)+xf_d(1,t)=0$. The polynomial $f_d(1,T)\in k[T]$ is nonzero (its coefficients are those of $f_d$), so $f_d(1,t)\ne0$ ($t$ transcendental) and
\[
x=-\frac{f_{d-1}(1,t)}{f_d(1,t)}\in k(t),\qquad y=tx\in k(t).
\]
Thus $K(Y)=k(x,y)\subseteq k(t)\subseteq K(Y)$, i.e.\ $K(Y)=k(t)$ is a pure transcendental extension of $k$ of degree~$1$, so $K(Y)\cong K(\mathbf P^1)$ and $Y$ is birational to $\mathbf P^1$ (Cor.~4.5): $Y$ is rational. \hfill$\square$

\exercise{7.6}{(\emph{Linear varieties.}) Show that an algebraic set $Y$ of pure dimension $r$ (i.e.\ every irreducible component of $Y$ has dimension $r$) has degree~$1$ if and only if $Y$ is a linear variety (Ex.~2.11). [\emph{Hint}: First, use (7.7) and treat the case $\dim Y=1$. Then do the general case by cutting with a hyperplane and using induction.]}
\pf
Recall: a linear variety of dimension $r$ in $\mathbf P^n$ is a set $Y$ with $I(Y)$ generated by $n-r$ independent linear forms, i.e.\ $Y=\mathbf P(V)$ for a subspace $V\subseteq k^{n+1}$ of dimension $r+1$.

($\Leftarrow$) $S(Y)\cong k[z_0,\dots,z_r]$, so $h_Y(l)=\binom{l+r}{r}$ and $\deg Y=1$.

($\Rightarrow$) \emph{Step 0: $Y$ is irreducible.} Let $Y=Y_1\cup\dots\cup Y_m$ be the decomposition into irreducible components, each of dimension $r$. By (7.6(b)) applied inductively ($Y_1\cup\dots\cup Y_{j-1}$ and $Y_j$ both have dimension $r$ and their intersection lies in $\bigcup_{i<j}(Y_i\cap Y_j)$, of dimension $<r$), $\deg Y=\sum_i\deg Y_i$, with each $\deg Y_i$ a positive integer by (7.6(a)). So $\deg Y=1$ forces $m=1$: $Y$ is a variety. If $r=0$, $Y$ is a point (a linear variety), so let $r\ge1$.

Observation: if $H$ is a hyperplane with $Y\not\subseteq H$, then every irreducible component $Z_j$ of $Y\cap H$ has dimension exactly $r-1$ ($\ge r-1$ by Thm.~7.2, $\le r-1$ as a proper closed subset of $Y$), and Thm.~7.7 gives
\[
\sum_ji(Y,H;Z_j)\deg Z_j=\deg Y\cdot\deg H=1.
\]
Each $i(Y,H;Z_j)=\operatorname{length}M_{\fp_j}\ge1$ ($\fp_j$ is a minimal prime of $M$), and $\deg Z_j\ge1$. So $Y\cap H$ is irreducible (a single $Z_1$), $\deg Z_1=1$, $i=1$.

\emph{Case $r=1$.} Then $Z_1$ is a point, so $Y\cap H$ is a single point for every hyperplane $H\not\supseteq Y$. Take two distinct points $P,Q\in Y$ ($Y$ is infinite). Every hyperplane $H$ containing $P,Q$ must contain $Y$ (otherwise $Y\cap H\supseteq\{P,Q\}$ has $\ge2$ points). The intersection of all hyperplanes containing the line $\overline{PQ}$ is $\overline{PQ}$, so $Y\subseteq\overline{PQ}\cong\mathbf P^1$; being closed, irreducible of dimension~$1$, $Y=\overline{PQ}$, a line.

\emph{Case $r\ge2$, induction.} Let $P\ne Q$ in $Y$. Since $\dim Y\ge2$, $Y\not\subseteq\overline{PQ}$, so some hyperplane $H\supseteq\overline{PQ}$ does not contain $Y$ (otherwise $Y\subseteq\bigcap H=\overline{PQ}$). By the observation $Y\cap H$ is a variety of dimension $r-1$ and degree~$1$, hence a linear variety by induction, and it contains $P,Q$, hence the line $\overline{PQ}$. So $Y$ contains the line through any two of its points. If $V\subseteq k^{n+1}$ is the set of $0$ together with all representatives of points of $Y$, then $V$ is closed under scalars and under addition (for independent $v,w\in V$, $[v+w]\in\overline{[v][w]}\subseteq Y$; for dependent $v,w$ it is clear), so $V$ is a subspace and $Y=\mathbf P(V)$ is a linear variety. \hfill$\square$

\exercise{7.7}{Let $Y$ be a variety of dimension $r$ and degree $d>1$ in $\mathbf P^n$. Let $P\in Y$ be a nonsingular point.
\begin{enumerate}[label=(\alph*),leftmargin=*]
\item Show that $X$ is a variety of dimension $r+1$, where $X$ is the closure of the union of all lines $PQ$, where $Q\in Y$, $Q\ne P$.
\item Show that $\deg X<d$. [\emph{Hint}: Use induction on $\dim Y$.]
\end{enumerate}}
\pf
\emph{Remark on method.} I do not follow the suggested induction; instead I give a direct commutative-algebra argument via Hilbert functions (the standard facts used are listed where they occur: finiteness of integral closure (3.9A), Krull's description of a noetherian normal domain as the intersection of its localizations at height-one primes, and the dimension theory of Section~1).

Since $d>1$ and a point has degree $1$ we have $r\ge1$. After a linear change of coordinates $P=(1:0:\dots:0)$. Put $S'=k[x_1,\dots,x_n]\subset S$, $I=I(Y)$ (prime), $S_Y=S/I$, $I'=I\cap S'$ (a homogeneous prime of $S'$), and let $B=S'/I'\subseteq S_Y$ be the subring generated by $\bar x_1,\dots,\bar x_n$. Then $S_Y=B[\bar x_0]$, and $B$, $S_Y$ are graded domains generated in degree~$1$ with $\dim S_Y=r+1$ (Krull dimension $=$ dimension of the affine cone), so $\dim B\in\{r,r+1\}$ ($B\subseteq S_Y=B[x_0]$ gives $\operatorname{trdeg}$ of fraction fields differing by at most one). Let $W=Z(I')\subseteq\mathbf P^{n-1}$, so $S(W)=B$. Write $x_0$ for $\bar x_0$.

Since $P\in Y$:
\[
(\star)\qquad F(1,0,\dots,0)=0\ \text{ for all }F\in I.
\]

\emph{Step 1: $X=Z(I'S)$, a cone over $W$.} Let $p=(1,0,\dots,0)$ and $C(Y)\subseteq k^{n+1}$ the affine cone, with $I(C(Y))=I$. The union of the lines $PQ$ ($Q\in Y\setminus P$) is $U_1=\{[y+tp]:y\in C(Y),\ t\in k,\ y+tp\ne0\}$ (if $y\in kp$ then $[y+tp]=P$, which lies on every such line, and $Y\setminus P\ne\emptyset$). Clearly $I'S\subseteq I(U_1)$, since $g\in I'$ depends only on $x_1,\dots,x_n$ and vanishes on the projection of $C(Y)$. Conversely let $F\in I(U_1)$ be homogeneous, $F=\sum_{i\le m}x_0^iF_i(x')$ with $F_m\neq0$. Then $F(x_0+t,x')\in I$ for every $t\in k$; expanding, $F(x_0+t,x')=\sum_{j=0}^mt^jG_j$ with $G_j=\sum_{i\ge j}\binom ijx_0^{i-j}F_i$, and since $k$ is infinite (Vandermonde) each $G_j\in I$; in particular $G_m=F_m\in I\cap S'=I'$. Then $F-x_0^mF_m\in I(U_1)$ is homogeneous with smaller $x_0$-degree, so by induction $F\in I'S$. Thus $I(U_1)=I'S$ and $X=\overline{U_1}=Z(I'S)$. Since $S/I'S=B[x_0]$ is a polynomial ring over the domain $B$, $I'S$ is prime, so $X$ is a variety, $S(X)=B[x_0]$, and $\dim X=\dim_{\rm Krull}B[x_0]-1=\dim B$. As $Y\subseteq X$ we get $r\le\dim X=\dim B\le r+1$.

\emph{Step 2: $\dim B=r+1$.} Suppose $\dim B=r$. Then $\dim X=r=\dim Y$ and $Y\subseteq X$ irreducible closed, so $Y=X$, i.e.\ $I=I'S$. Then $S_Y=B[x_0]$ and the coordinate ring of $Y\cap\{x_0\ne0\}$ is $k[y_1,\dots,y_n]/I'(y)\cong B$ (ungraded), with $P$ corresponding to the maximal ideal $B_+=\bigoplus_{l>0}B_l$. $P$ is nonsingular of dimension $r$, so $\mathcal O_{P,Y}=B_{B_+}$ is a regular local ring of dimension $r$ (Thm.~5.1), i.e.\ $\dim_k B_+/B_+^2=r$. But $B_+/B_+^2=B_1$ ($B$ is generated by $B_1$), so $\dim_kB_1=r$. The surjection $k[z_1,\dots,z_r]\to B$ sending the $z_i$ to a basis of $B_1$ is a surjection of domains of the same dimension $r$, so it is an isomorphism ($k[z]/\ker$ has smaller dimension unless $\ker=0$). Hence $I'$ is generated by $n-r$ independent linear forms and $I=I'S$ too; $Y$ is a linear variety, $\deg Y=1$, contradicting $d>1$. So $\dim B=r+1$.

This proves (a): $X$ is a variety, $\dim X=\dim B=r+1$.

\emph{Step 3: $\deg X=e:=\deg W$.} $\dim W=r$, so $h_B(t)=P_W(t)=\frac e{r!}t^r+\dots$ for $t\gg0$. As $S(X)=B[x_0]=\bigoplus_jx_0^jB$, $h_X(l)=\sum_{t=0}^lh_B(t)=\frac e{(r+1)!}l^{r+1}+O(l^r)$, so $\deg X=e$.

\emph{Step 4: $d\ge e+1$.} Let $F_j=\sum_{i\le j}x_0^iB\subseteq S_Y$ ($F_{-1}=0$), graded $B$-submodules, $S_Y=\bigcup F_j$. Let $\fa_j=\{b\in B:\ x_0^jb\in F_{j-1}\}$, a homogeneous ideal; $\fa_0=0$, $F_j/F_{j-1}\cong(B/\fa_j)(-j)$, and $\fa_j\subseteq\fa_{j+1}$ (as $x_0^{j+1}b=x_0\cdot x_0^jb\in x_0F_{j-1}\subseteq F_j$). Let $\fa=\bigcup_j\fa_j$, a homogeneous ideal of $B$. Since $(S_Y)_l=\sum_{j\le l}x_0^jB_{l-j}$,
\[
h_Y(l)=\sum_{j=0}^l\dim(B/\fa_j)_{l-j}\ \ge\ h_B(l)+\sum_{t=0}^{l-1}h_{B/\fa}(t).\tag{1}
\]

\emph{Claim: $\dim(B/\fa)=r$ (Krull dimension).} First, $x_0$ is algebraic over $K_B=\operatorname{Frac}B$ since $\operatorname{trdeg}\operatorname{Frac}S_Y=r+1=\operatorname{trdeg}K_B$; let $L=K_B(x_0)=\operatorname{Frac}S_Y$. A nonzero relation of $x_0$ over $B$ may be taken homogeneous (take a homogeneous component whose leading coefficient $c_0\ne0$); it says $c_0\in\fa_j$ for some $j$, so $\fa\ne0$.

Let $C$ be the integral closure of $B$ in $L$; it is a finite $B$-module (3.9A), a noetherian normal domain, a finitely generated $k$-algebra with $\dim C=\dim B=r+1$, and $\operatorname{Frac}C=L$ (some $bx_0$, $b\in B\setminus0$, is integral over $B$). \emph{$x_0\notin C$:} a monic relation $x_0^N+c_1x_0^{N-1}+\dots+c_N=0$ with $c_i\in B$ could be replaced by its degree-$N$ component, giving $F=x_0^N+\sum_{i\ge1}\tilde c_ix_0^{N-i}\in I$ with $\tilde c_i\in S'_i$, and then $F(1,0,\dots,0)=1\ne0$ contradicts $(\star)$. By Krull ($C=\bigcap_{\operatorname{ht}\fq=1}C_\fq$) there is a prime $\fq$ of $C$ of height~$1$ with $x_0\notin C_\fq$. Put $\fp=\fq\cap B$. If some $b\in\fa\setminus\fp$ existed, say $b\in\fa_j$, then $x_0^j=\sum_{i<j}(c_i/b)x_0^i$ with $c_i\in B$, so $x_0$ would be integral over $B_\fp\subseteq C_\fq$; as $C_\fq$ is integrally closed in $L$ (a DVR), $x_0\in C_\fq$, a contradiction. Hence $\fa\subseteq\fp$. Finally $\dim C/\fq=\dim C-\operatorname{ht}\fq=r$ (Thm.~1.8A) and $C/\fq$ is integral over $B/\fp$, so $\dim B/\fp=r$, whence $\dim B/\fa\ge\dim B/\fp=r$; and $\dim B/\fa\le r$ since $\fa\ne0$ and $B$ is a domain of dimension $r+1$. This proves the claim.

Now $B/\fa=S'/\fb$ for a homogeneous ideal $\fb\supseteq I'$ and, by Thm.~7.5, $\deg P_{B/\fa}=\dim Z(\fb)\subseteq\mathbf P^{n-1}$; since $\dim_{\rm Krull}S'/\fb=r\ge1$, $Z(\fb)\neq\emptyset$ and $\dim Z(\fb)=r-1$. So $P_{B/\fa}(t)=\frac{c}{(r-1)!}t^{r-1}+\dots$ with $c$ a positive integer (by 7.3(a) the leading coefficient is an integer divided by $(r-1)!$, and it is positive since $h_{B/\fa}\ge0$ and the degree is exactly $r-1$). Therefore $\sum_{t=0}^{l-1}h_{B/\fa}(t)=\frac{c}{r!}l^r+O(l^{r-1})$. Comparing leading coefficients of $l^r$ in (1), where $h_Y(l)=\frac d{r!}l^r+\dots$ and $h_B(l)=\frac e{r!}l^r+\dots$, gives $d\ge e+c\ge e+1$.

So $\deg X=e\le d-1<d$, which is (b). \hfill$\square$

\exercise{7.8}{Let $Y\subseteq\mathbf P^n$ be a variety of degree $2$. Show that $Y$ is contained in a linear subspace $L$ of dimension $r+1$ in $\mathbf P^n$ ($r=\dim Y$). Thus $Y$ is isomorphic to a quadric hypersurface in $\mathbf P^{r+1}$ (Ex.~5.12).}
\pf
Since a point has degree $1$, $r\ge1$. The singular locus of $Y$ is a proper closed subset (Thm.~5.3), so $Y$ has a nonsingular point $P$. By Ex.~7.7, the closure $X$ of the union of the lines $PQ$ ($Q\in Y\setminus P$) is a variety of dimension $r+1$ with $\deg X<2$, i.e.\ $\deg X=1$ (degrees are positive integers, 7.6(a)). By Ex.~7.6, $X$ is a linear variety of dimension $r+1$, say $X=L\cong\mathbf P^{r+1}$ (Ex.~2.11), and $Y\subseteq X=L$ since $P$ and every $Q\in Y\setminus P$ lie on the line $PQ$.

Under the identification $L\cong\mathbf P^{r+1}$, $Y$ is an irreducible closed subset of dimension $r$, i.e.\ a hypersurface in $\mathbf P^{r+1}$: $I(Y)$ is a prime of height~$1$ in the UFD $k[z_0,\dots,z_{r+1}]$, hence $I(Y)=(g)$ with $g$ irreducible (Prop.~1.13). After a linear change of coordinates $L=\{x_{r+2}=\dots=x_n=0\}$, so $S(Y)\cong k[z_0,\dots,z_{r+1}]/(g)$ and by (7.6(d)) $\deg Y=\deg g$, so $\deg g=2$ and $Y$ is a quadric hypersurface in $L\cong\mathbf P^{r+1}$. \hfill$\square$

\end{document}
